% Sensibilisation sur le rôle joué par les végétaux verts à travers la photosynthèse au sein de l'environnement — SVT 1ère D — Cours signé OnBuch+
% Programme officiel MINESEC. Compiler avec : tectonic ND06-role-vegetaux-verts-photosynthese.tex
\newcommand{\DOCMATIERE}{SVT}
\newcommand{\DOCNIVEAU}{1ère D}
\newcommand{\DOCMODULE}{Module 1 : Le Monde Vivant}
\newcommand{\DOCLECON}{ND06}
\newcommand{\DOCTITRE}{Sensibilisation sur le rôle joué par les végétaux verts à travers la photosynthèse au sein de l'environnement}
% OnBuch+ — préambule « cours signés » ENRICHI (compilé avec tectonic / XeLaTeX)
% Système visuel « Pop » d'OnBuch+ : fond crème (jamais blanc), bordures encre
% épaisses, ombres « dures » décalées, titres Archivo Black en capitales.
% Version MATHÉMATIQUES : graphes (pgfplots/tikz), boîtes pédagogiques
% (définition, propriété, méthode, attention, le savais-tu, exercice résolu,
% exercice, TP, bilan), page de garde et signature OnBuch+ ancrée en bas.
%
% Macros attendues dans main.tex AVANT \input{preamble} :
%   \DOCMATIERE  \DOCNIVEAU  \DOCMODULE  \DOCLECON (numéro)  \DOCTITRE
\documentclass[11pt]{article}

\usepackage{fontspec}
\usepackage[french]{babel}
\usepackage[a4paper,margin=2cm,top=3.4cm,bottom=2.8cm,headheight=1.4cm,footskip=1.3cm]{geometry}
\usepackage{amsmath,amssymb,amsfonts}
\usepackage{mathtools} % \xmapsto, \coloneqq... souvent utilisés par les modèles
\usepackage{cancel} % simplifications barrées (bilans de forces)
% \abs{z} (module, valeur absolue) et \norm{u} (norme), utilisés par les modèles
\providecommand{\abs}{}\let\abs\relax\DeclarePairedDelimiter{\abs}{\lvert}{\rvert}
\providecommand{\norm}{}\let\norm\relax\DeclarePairedDelimiter{\norm}{\lVert}{\rVert}
\usepackage[table]{xcolor} % [table] : fournit aussi \rowcolors (alternance de lignes)
\usepackage{pagecolor}
\usepackage{tikz}
\usetikzlibrary{arrows.meta,positioning,calc,shapes.geometric,shapes.misc,shapes.arrows,decorations.pathmorphing,decorations.pathreplacing,decorations.markings,patterns,shadows,mindmap,trees,fit,backgrounds,matrix,intersections,angles,quotes}
\usepackage{pgfplots}
\usepackage[version=4]{mhchem} % équations nucléaires \ce{^{235}_{92}U}
\usepackage[european,siunitx]{circuitikz} % schémas électriques
\pgfplotsset{compat=1.18}
\usepgfplotslibrary{fillbetween,groupplots} % aires entre courbes (calcul intégral)
\usepackage{enumitem}
\usepackage{fancyhdr}
% [most] charge toutes les bibliothèques tcolorbox (listings, minted, raster,
% documentation...) et gonfle inutilement la mémoire TeX sur un document long
% (~300 boîtes) : on ne charge que ce qui est réellement utilisé.
\usepackage[skins,breakable]{tcolorbox}
\usepackage{graphicx}
\usepackage{colortbl}
\usepackage{booktabs}
\usepackage{tabularx}
\usepackage{diagbox} % case d'en-tête diagonale des tableaux à double entrée
% Colonnes extensibles centrée / gauche / droite (C, L, R), souvent utilisées par les modèles
\newcolumntype{C}{>{\centering\arraybackslash}X}
\newcolumntype{L}{>{\raggedright\arraybackslash}X}
\newcolumntype{R}{>{\raggedleft\arraybackslash}X}
\usepackage{array}
\usepackage{multicol}
\usepackage{siunitx}
% parse-numbers=false : accepte aussi \SI{2,5 \times 10^{-3}}{...}
\sisetup{locale=FR,output-decimal-marker={,},group-minimum-digits=4,parse-numbers=false}
\DeclareSIUnit\ohm{\ensuremath{\symup{\Omega}}} % U+2126 absent de la police
\DeclareSIUnit\division{div} % réglages d'oscilloscope (ms/div, V/div)
\DeclareSIUnit\tr{tr} % tours (vitesse de rotation, ex. \SI{1500}{\tr\per\minute})
\DeclareSIUnit\ch{ch} % cheval-vapeur (puissance moteur, unité usuelle hors SI)
\DeclareSIUnit\franccfa{FCFA} % franc CFA (coûts, contexte camerounais)
\DeclareSIUnit\dioptre{\ensuremath{\delta}} % vergence d'une lentille (dioptrie)
\usepackage{qrcode}
% \degree : commande fréquemment utilisée par les modèles pour un angle ou une
% température en dehors de \SI{}{} (non fournie par les paquets chargés).
\providecommand{\degree}{^{\circ}}
% \qed (fin de démonstration), utilisé par les modèles sans amsthm
\providecommand{\qed}{\hfill$\square$}
% \barre : double barre des valeurs interdites dans un tableau de variations (tkz-tab)
\providecommand{\barre}{\ensuremath{\|}}
% environnement proof (amsthm non chargé) utilisé par les modèles
\ifdefined\proof\else\newenvironment{proof}[1][Démonstration]{\par\noindent\textit{#1.}\ }{\qed\par}\fi
\usepackage{lastpage}
\usepackage{hyperref}
\hypersetup{hidelinks}

% ============================================================
% Palette « Pop » OnBuch+ — mêmes hex que la classe P (plus_ui.dart)
% ============================================================
\definecolor{poporange}{HTML}{F59321}
\definecolor{popdark}{HTML}{E07A0C}
\definecolor{popcream}{HTML}{FFF6E8}
\definecolor{popink}{HTML}{1C1714}
\definecolor{poppurple}{HTML}{7A5AE0}
\definecolor{popgreen}{HTML}{2E9E6A}
\definecolor{poppink}{HTML}{E0457B}
\definecolor{popblue}{HTML}{2F7DE1}
\definecolor{popgold}{HTML}{C98A12}
\definecolor{popmuted}{HTML}{8A7F76}
\definecolor{popline}{HTML}{EADFD2}
\definecolor{popgray}{HTML}{8A7F76}
\colorlet{popgrey}{popgray}
% Alias tolérants (le modèle nomme parfois les couleurs différemment)
\colorlet{popwhite}{white}
\colorlet{popblack}{popink}
\colorlet{popred}{poppink}
\colorlet{popyellow}{popgold}
% Teintes claires pour fonds de tableaux / zones
\colorlet{popblueL}{popblue!10!white}
\colorlet{poporangeL}{poporange!14!white}
\colorlet{popgreenL}{popgreen!12!white}
\colorlet{poppurpleL}{poppurple!12!white}
\colorlet{poppinkL}{poppink!12!white}
\colorlet{popgoldL}{popgold!14!white}

\pagecolor{popcream}

% ============================================================
% Typographie
% ============================================================
\defaultfontfeatures{Ligatures=TeX}
\setmainfont{PlusJakartaSans-Medium.ttf}[Path=../../../fonts/,BoldFont=PlusJakartaSans-Bold.ttf]
% Maths en sans-serif (Fira Math) avec chiffres et lettres droites de Plus
% Jakarta, pour que les formules se fondent dans le texte.
\usepackage[math-style=ISO,bold-style=ISO]{unicode-math}
\setmathfont{FiraMath-Regular.otf}[Path=../../../fonts/]
\setmathfont{PlusJakartaSans-Medium.ttf}[Path=../../../fonts/,range={up/num,up/{Latin,latin},\mathup}]
\usepackage{icomma} % virgule décimale sans espace en mode math
% Caractères mathématiques Unicode absents des polices de texte (Plus Jakarta,
% Archivo Black) : rendus par leur équivalent mathématique, en texte comme en
% formule — sinon ils s'affichent en carré vide (ex. « Arithmétique dans ℤ »).
\usepackage{newunicodechar}
\newunicodechar{ℝ}{\ensuremath{\mathbb{R}}}
\newunicodechar{ℤ}{\ensuremath{\mathbb{Z}}}
\newunicodechar{ℂ}{\ensuremath{\mathbb{C}}}
\newunicodechar{ℕ}{\ensuremath{\mathbb{N}}}
\newunicodechar{ℚ}{\ensuremath{\mathbb{Q}}}
\newunicodechar{𝔻}{\ensuremath{\mathbb{D}}}
\newunicodechar{α}{\ensuremath{\alpha}}
\newunicodechar{β}{\ensuremath{\beta}}
\newunicodechar{γ}{\ensuremath{\gamma}}
\newunicodechar{δ}{\ensuremath{\delta}}
\newunicodechar{ε}{\ensuremath{\varepsilon}}
\newunicodechar{θ}{\ensuremath{\theta}}
\newunicodechar{λ}{\ensuremath{\lambda}}
\newunicodechar{μ}{\ensuremath{\mu}}
\newunicodechar{σ}{\ensuremath{\sigma}}
\newunicodechar{τ}{\ensuremath{\tau}}
\newunicodechar{φ}{\ensuremath{\varphi}}
\newunicodechar{ψ}{\ensuremath{\psi}}
\newunicodechar{ω}{\ensuremath{\omega}}
\newunicodechar{Σ}{\ensuremath{\Sigma}}
% majuscules produites par \MakeUppercase dans les titres (α-aminés -> Α)
\newunicodechar{Α}{\ensuremath{\alpha}}
\newunicodechar{Β}{\ensuremath{\beta}}
\newunicodechar{Γ}{\ensuremath{\Gamma}}
\newunicodechar{Δ}{\ensuremath{\Delta}}
\newunicodechar{Ε}{\ensuremath{\varepsilon}}
\newunicodechar{Θ}{\ensuremath{\Theta}}
\newunicodechar{Λ}{\ensuremath{\Lambda}}
\newunicodechar{Μ}{\ensuremath{\mu}}
\newunicodechar{Τ}{\ensuremath{\tau}}
\newunicodechar{Φ}{\ensuremath{\Phi}}
\newunicodechar{Ψ}{\ensuremath{\Psi}}
\newunicodechar{Ω}{\ensuremath{\Omega}}
\newunicodechar{↦}{\ensuremath{\mapsto}}
\newunicodechar{∈}{\ensuremath{\in}}
\newunicodechar{∉}{\ensuremath{\notin}}
\newunicodechar{∧}{\ensuremath{\wedge}}
\newunicodechar{⇔}{\ensuremath{\Leftrightarrow}}
\newunicodechar{⟺}{\ensuremath{\Longleftrightarrow}}
\newunicodechar{⇒}{\ensuremath{\Rightarrow}}
\newunicodechar{⟹}{\ensuremath{\Longrightarrow}}
\newunicodechar{∘}{\ensuremath{\circ}}
\newunicodechar{∪}{\ensuremath{\cup}}
\newunicodechar{∩}{\ensuremath{\cap}}
\newunicodechar{≡}{\ensuremath{\equiv}}
\newunicodechar{⊂}{\ensuremath{\subset}}
\newunicodechar{⊥}{\ensuremath{\perp}}
\newunicodechar{∃}{\ensuremath{\exists}}
\newunicodechar{∀}{\ensuremath{\forall}}
\newunicodechar{∛}{\ensuremath{\sqrt[3]{\,}}}
\newunicodechar{∜}{\ensuremath{\sqrt[4]{\,}}}
\newunicodechar{⃗}{\ensuremath{{}^{\to}}}
\newunicodechar{ⁿ}{\ifmmode^{n}\else\textsuperscript{\ensuremath{n}}\fi}
\newunicodechar{ˣ}{\ifmmode^{x}\else\textsuperscript{\ensuremath{x}}\fi}
\newunicodechar{ᵇ}{\ifmmode^{b}\else\textsuperscript{\ensuremath{b}}\fi}
\newunicodechar{ʳ}{\ifmmode^{r}\else\textsuperscript{\ensuremath{r}}\fi}
\newunicodechar{ᵘ}{\ifmmode^{u}\else\textsuperscript{\ensuremath{u}}\fi}
\newunicodechar{ʸ}{\ifmmode^{y}\else\textsuperscript{\ensuremath{y}}\fi}
\newunicodechar{ᵃ}{\ifmmode^{a}\else\textsuperscript{\ensuremath{a}}\fi}
\newunicodechar{ᵉ}{\ifmmode^{e}\else\textsuperscript{\ensuremath{e}}\fi}
\newunicodechar{ᵏ}{\ifmmode^{k}\else\textsuperscript{\ensuremath{k}}\fi}
\newunicodechar{ᵖ}{\ifmmode^{p}\else\textsuperscript{\ensuremath{p}}\fi}
\newunicodechar{ᴺ}{\ifmmode^{N}\else\textsuperscript{\ensuremath{N}}\fi}
\newunicodechar{ᶜ}{\ifmmode^{c}\else\textsuperscript{\ensuremath{c}}\fi}
\newunicodechar{ʰ}{\ifmmode^{h}\else\textsuperscript{\ensuremath{h}}\fi}
\newunicodechar{ᵠ}{\ifmmode^{q}\else\textsuperscript{\ensuremath{q}}\fi}
\newunicodechar{ₙ}{\ifmmode_{n}\else\textsubscript{\ensuremath{n}}\fi}
\newunicodechar{ₐ}{\ifmmode_{a}\else\textsubscript{\ensuremath{a}}\fi}
\newunicodechar{ᵢ}{\ifmmode_{i}\else\textsubscript{\ensuremath{i}}\fi}
\newunicodechar{ᵣ}{\ifmmode_{r}\else\textsubscript{\ensuremath{r}}\fi}
\newunicodechar{ₖ}{\ifmmode_{k}\else\textsubscript{\ensuremath{k}}\fi}
\newunicodechar{ᵦ}{\ifmmode_{\beta}\else\textsubscript{\ensuremath{\beta}}\fi}
\newunicodechar{ₓ}{\ifmmode_{x}\else\textsubscript{\ensuremath{x}}\fi}
\newfontfamily\popdisplay{ArchivoBlack-Regular.ttf}[Path=../../../fonts/]
\newfontfamily\poplogo{SpaceGrotesk-Bold.ttf}[Path=../../../fonts/]
\newfontfamily\popsemi{PlusJakartaSans-SemiBold.ttf}[Path=../../../fonts/]
\newfontfamily\popxbold{PlusJakartaSans-ExtraBold.ttf}[Path=../../../fonts/]
\newfontfamily\popxboldls{PlusJakartaSans-ExtraBold.ttf}[Path=../../../fonts/,LetterSpace=3.0]

\setlist[enumerate]{leftmargin=1.7em,itemsep=5pt,topsep=4pt}
\setlist[itemize]{leftmargin=1.7em,itemsep=4pt,topsep=4pt,label={\color{poporange}\textbullet}}
\setlength{\parindent}{0pt}
\setlength{\parskip}{5pt}
\renewcommand{\arraystretch}{1.25}

% pgfplots : style Pop par défaut
\pgfplotsset{
  every axis/.append style={axis line style={popink,line width=1pt},tick style={popink},
    grid style={popline,line width=0.5pt},label style={font=\small},tick label style={font=\footnotesize}},
  popcurve/.style={poporange,line width=1.8pt,no markers,smooth},
}
% Même style disponible dans un \draw TikZ simple (hors pgfplots)
\tikzset{popcurve/.style={poporange,line width=1.8pt}}
\tikzset{popgrid/.style={popline,very thin}}
\colorlet{popcurve}{poporange} % le nom du style est parfois utilisé comme couleur

% ============================================================
% Pill / squiggle
% ============================================================
\newcommand{\poppill}[3]{%
  \tikz[baseline=-0.55ex]\node[draw=popink,line width=1.2pt,rounded corners=2.6pt,%
    fill=#2,inner xsep=6.5pt,inner ysep=3pt]{%
    \popxboldls\fontsize{7}{7}\selectfont\color{#3}\MakeUppercase{#1}};%
}
\newcommand{\popsquiggle}[1]{%
  \tikz[baseline=-0.4ex]{
    \draw[#1,line width=2.6pt,line cap=round]
      plot [smooth] coordinates {(0,0.09) (0.34,0.01) (0.68,0.21) (1.02,0.01) (1.36,0.10)};
  }%
}
% Mot-clé surligné (marqueur orange sous le texte)
\newcommand{\cle}[1]{{\popxbold\color{popdark}#1}}

% ============================================================
% En-tête / pied
% ============================================================
\pagestyle{fancy}
\fancyhf{}
\renewcommand{\headrulewidth}{0pt}
\renewcommand{\footrulewidth}{0pt}
\lhead{\raisebox{-0.15ex}{\poplogo\fontsize{13}{13}\selectfont\color{popink}OnBuch\color{poporange}+}}
\rhead{\poppill{\DOCMATIERE\ \textbullet\ \DOCNIVEAU\ \textbullet\ Leçon \DOCLECON}{white}{popink}}
\lfoot{\popxbold\fontsize{7.6}{7.6}\selectfont\color{popmuted}\MakeUppercase{Cours signé OnBuch+}\ \textbullet\ {\popsemi\DOCTITRE}}
\rfoot{\tikz[baseline=-0.6ex]\node[rounded corners=8.5pt,fill=popink,minimum height=17pt,minimum width=17pt,inner xsep=5pt,inner sep=0]{\popxbold\fontsize{8}{8}\selectfont\color{popcream}\thepage\,/\,\pageref*{LastPage}};}

% ============================================================
% Titres
% ============================================================
\newcommand{\chaptitre}[2]{%
  \begin{tcolorbox}[enhanced,colback=poporange,colframe=popink,boxrule=2.6pt,arc=11pt,
      drop shadow=popink,left=15pt,right=15pt,top=12pt,bottom=11pt,before skip=3pt,after skip=18pt]
    \poppill{#2}{popink}{popcream}\par\vspace{9pt}
    {\raggedright\hyphenpenalty=10000\exhyphenpenalty=10000\popdisplay\fontsize{22}{24}\selectfont\color{popcream}\MakeUppercase{#1}\par}
  \end{tcolorbox}
}
% Section numérotée : pastille orange avec le numéro + titre Archivo
\newcounter{popsec}
\newcommand{\coursec}[1]{%
  \stepcounter{popsec}\setcounter{popsub}{0}%
  \par\vspace{16pt}\noindent
  \begin{minipage}{\linewidth}
  \tikz[baseline=-0.7ex]\node[draw=popink,line width=1.6pt,fill=poporange,rounded corners=4pt,minimum size=22pt,inner sep=0]{\popdisplay\fontsize{12}{12}\selectfont\color{popcream}\thepopsec};%
  \hspace{8pt}{\popdisplay\fontsize{15}{17}\selectfont\color{popink}\MakeUppercase{#1}}\par
  \vspace{4pt}\noindent{\color{poporange}\rule{36pt}{3.6pt}}
  \end{minipage}\par\nopagebreak\vspace{8pt}
}
\newcounter{popsub}
\newcommand{\courssub}[1]{%
  \stepcounter{popsub}%
  \par\vspace{9pt}\noindent{\popxbold\fontsize{11.5}{13}\selectfont\color{popink}{\color{poporange}\thepopsec.\thepopsub}\hspace{6pt}#1}\par\nopagebreak\vspace{4pt}
}

% ============================================================
% Boîtes pédagogiques — même recette : fond blanc, bordure épaisse colorée,
% ombre dure encre, badge plein. Titre optionnel [#1].
% ============================================================
\tcbset{popbase/.style={breakable,enhanced,colback=white,boxrule=2.3pt,arc=9pt,
  shadow={3pt}{-3pt}{0pt}{popink},left=14pt,right=14pt,top=11pt,bottom=11pt,before skip=11pt,after skip=15pt,
  fontupper=\popsemi\fontsize{10.6}{14.5}\selectfont\color{popink},
  fontlower=\popsemi\fontsize{10.6}{14.5}\selectfont\color{popink}}}
% Coche dessinée (le glyphe ✓ n'existe pas dans les polices du document)
\newcommand{\popcheck}[1]{\tikz[baseline=-0.5ex]\draw[#1,line width=1.6pt,line cap=round,line join=round](-0.13,0.0)--(-0.04,-0.09)--(0.13,0.1);}
\providecommand{\coche}{\popcheck{popgreen}}
\providecommand{\resultat}[1]{\par\smallskip\noindent\fcolorbox{popgreen}{popgreenL}{\parbox{\dimexpr\linewidth-2\fboxsep-2\fboxrule\relax}{\textbf{Résultat.} #1}}\par\smallskip}
\newcommand{\popboxhead}[3]{\poppill{#1}{#2}{white}\ifx\relax#3\relax\else\hspace{6pt}{\popxbold\fontsize{10.6}{12}\selectfont\color{popink}#3}\fi\par\vspace{7pt}}

\newtcolorbox{aretenir}[1][]{popbase,colframe=popgreen,before upper={\popboxhead{À retenir}{popgreen}{#1}}}
\newtcolorbox{exemplebox}[1][]{popbase,colframe=poporange,before upper={\popboxhead{Exemple}{poporange}{#1}}}
\newtcolorbox{definition}[1][]{popbase,colframe=popblue,before upper={\popboxhead{Définition}{popblue}{#1}}}
\newtcolorbox{propriete}[1][]{popbase,colframe=poppurple,before upper={\popboxhead{Propriété}{poppurple}{#1}}}
\newtcolorbox{methode}[1][]{popbase,colframe=popgold,before upper={\popboxhead{Méthode}{popgold}{#1}}}
\newtcolorbox{attention}[1][]{popbase,colframe=poppink,before upper={\popboxhead{Attention}{poppink}{#1}}}
\newtcolorbox{experience}[1][]{popbase,colframe=popblue,colback=popblueL,before upper={\popboxhead{Expérience}{popblue}{#1}}}
% « Le savais-tu ? » : carte violette pleine (contexte, culture, Cameroun)
\newtcolorbox{savaistu}[1][]{popbase,colback=poppurple,colframe=popink,
  fontupper=\popsemi\fontsize{10.4}{14.2}\selectfont\color{white},
  before upper={\poppill{Le savais-tu\,?}{white}{poppurple}\ifx\relax#1\relax\else\hspace{6pt}{\popxbold\color{white}#1}\fi\par\vspace{7pt}}}
% Exercice résolu : énoncé en haut, \tcblower puis la solution sur fond crème
\newcounter{popexo}\newcounter{popexr}
\newtcolorbox{exoresolu}[1][]{popbase,colframe=poporange,colbacklower=popcream,
  segmentation style={popink,line width=1.2pt,dashed},
  before upper={\stepcounter{popexr}\popboxhead{Exercice résolu \thepopexr}{poporange}{#1}},
  before lower={\poppill{Solution}{popink}{popcream}\par\vspace{6pt}}}
% Exercice (non résolu en place) — la correction est dans la partie Corrigés
\newtcolorbox{exercice}[1][]{popbase,colframe=popink,
  before upper={\stepcounter{popexo}\popboxhead{Exercice \thepopexo}{popink}{#1}}}
% Corrigé
\newtcolorbox{corrige}[1][]{popbase,colframe=popgreen,colback=popgreenL,
  before upper={\popboxhead{Corrigé}{popgreen}{#1}}}
% Objectifs / prérequis / bilan
\newtcolorbox{objectifs}{popbase,colframe=popink,colback=poporangeL,
  before upper={\popboxhead{Objectifs de la leçon}{popink}{}}}
\newtcolorbox{prerequis}{popbase,colframe=popink,colback=popblueL,
  before upper={\popboxhead{Prérequis}{popblue}{}}}
\newtcolorbox{bilan}{popbase,colframe=popink,colback=white,boxrule=2.8pt,
  before upper={\popboxhead{Fiche bilan}{poporange}{L'essentiel à connaître pour l'examen}}}
% Figure encadrée avec légende
\newtcolorbox{popfigure}[1][]{enhanced,breakable=false,colback=white,colframe=popline,boxrule=1.4pt,arc=8pt,
  left=10pt,right=10pt,top=10pt,bottom=8pt,before skip=10pt,after skip=12pt,halign=center,
  lower separated=false,halign lower=center,
  fontlower=\popsemi\fontsize{9}{11.5}\selectfont\color{popmuted},
  #1}
\newcommand{\legende}[1]{\par\vspace{6pt}{\popsemi\fontsize{9}{11.5}\selectfont\color{popmuted}#1}}

% ============================================================
% Signature OnBuch+
% ============================================================
\newcommand{\poplogoinline}[1]{{\poplogo\fontsize{#1}{#1}\selectfont\color{popink}OnBuch\color{poporange}+}}
\newcommand{\signatureonbuch}{%
  \begin{tcolorbox}[enhanced,colback=popink,colframe=popink,boxrule=2.2pt,arc=10pt,
      drop shadow=poporange,left=14pt,right=14pt,top=10pt,bottom=10pt,before skip=0pt,after skip=0pt]
    \tikz[baseline=-0.7ex]\node[circle,fill=popgreen,draw=popcream,line width=1.3pt,minimum size=22pt,inner sep=0]{\popcheck{white}};%
    \hspace{9pt}\begin{minipage}[c]{0.62\linewidth}
      {\popxbold\fontsize{12}{13}\selectfont\color{popcream}Signé\ \textbullet\ {\poplogo OnBuch\color{poporange}+}}\par\vspace{2pt}
      {\raggedright\popsemi\fontsize{8}{10.5}\selectfont\color{popline}\MakeUppercase{Équipe pédagogique OnBuch+}\\\MakeUppercase{Conforme au programme officiel MINESEC}\par}
    \end{minipage}\hfill
    {\poplogo\fontsize{22}{22}\selectfont\color{popcream}OnBuch\color{poporange}+}
  \end{tcolorbox}%
}

% ============================================================
% Page de garde
% #1 titre · #2 matière · #3 niveau · #4 module · #5 n° leçon · #6 durée ·
% #7 description / objectifs (texte ou itemize)
% ============================================================
\newcommand{\pagegarde}[7]{%
  \thispagestyle{empty}
  \newgeometry{a4paper,margin=1.7cm,top=1.6cm,bottom=1.4cm}%
  \begin{tikzpicture}[remember picture,overlay]
    \fill[poporange!18] ([xshift=-3.2cm,yshift=-3.6cm]current page.north east) circle (4.6cm);
    \fill[poppurple!14] ([xshift=2.4cm,yshift=7.6cm]current page.south west) circle (3.8cm);
  \end{tikzpicture}%
  {\poplogo\fontsize{30}{30}\selectfont\color{popink}OnBuch\color{poporange}+}\hfill
  \raisebox{0.6ex}{\poppill{Cours signé \textbullet\ Premium}{popink}{popcream}}\par
  \vspace{1pt}\popsquiggle{poppurple}\par
  \vspace{8pt}
  \begin{tcolorbox}[enhanced,colback=poporange,colframe=popink,boxrule=2.8pt,arc=13pt,
      drop shadow=popink,left=18pt,right=18pt,top=12pt,bottom=13pt,after skip=10pt]
    \poppill{#2 \textbullet\ #3}{popink}{popcream}\hspace{5pt}\poppill{#4}{white}{popink}\par\vspace{8pt}
    {\popdisplay\fontsize{14}{14}\selectfont\color{popink}LEÇON #5}\par\vspace{4pt}
    {\raggedright\hyphenpenalty=10000\exhyphenpenalty=10000\popdisplay\fontsize{25}{27}\selectfont\color{popcream}\MakeUppercase{#1}\par}
    \vspace{8pt}
    {\popxbold\fontsize{9.5}{9.5}\selectfont\color{popink}\MakeUppercase{Durée conseillée : #6}}
  \end{tcolorbox}
  \begin{tcolorbox}[enhanced,colback=white,colframe=popink,boxrule=2.2pt,arc=10pt,
      drop shadow=popink,left=16pt,right=16pt,top=10pt,bottom=10pt,after skip=8pt]
    \poppill{Dans cette leçon}{poppurple}{white}\par\vspace{5pt}
    \popsemi\fontsize{9.3}{12.6}\selectfont\color{popink}#7
  \end{tcolorbox}
  \noindent\begin{minipage}[t]{0.487\linewidth}
    \begin{tcolorbox}[enhanced,colback=popgreen,colframe=popink,boxrule=2.2pt,arc=10pt,
        drop shadow=popink,left=11pt,right=11pt,top=10pt,bottom=10pt,height=3.1cm,valign=center]
      \tcbox[colback=white,colframe=popink,boxrule=1.3pt,arc=5pt,boxsep=0pt,nobeforeafter,
        left=3pt,right=3pt,top=3pt,bottom=3pt,valign=center]{\qrcode[height=1.9cm]{https://play.google.com/store/apps/details?id=com.luvvix.onbuch&hl=fr}}%
      \hspace{8pt}\begin{minipage}[c]{0.5\linewidth}
        {\popdisplay\fontsize{11}{12.5}\selectfont\color{white}\MakeUppercase{Télécharge OnBuch}}\par\vspace{4pt}
        {\raggedright\popsemi\fontsize{8.4}{11}\selectfont\color{white}Gratuit \textbullet\ Google Play\\Tuteur IA, annales, concours, résultats\par}
      \end{minipage}
    \end{tcolorbox}
  \end{minipage}\hfill
  \begin{minipage}[t]{0.487\linewidth}
    \begin{tcolorbox}[enhanced,colback=poppurple,colframe=popink,boxrule=2.2pt,arc=10pt,
        drop shadow=popink,left=11pt,right=11pt,top=10pt,bottom=10pt,height=3.1cm,valign=center]
      \tikz[baseline=-0.7ex]\node[circle,fill=white,draw=popink,line width=1.3pt,minimum size=1.6cm,inner sep=0]{\popdisplay\fontsize{24}{24}\selectfont\color{poppurple}+};%
      \hspace{8pt}\begin{minipage}[c]{0.6\linewidth}
        {\popdisplay\fontsize{11}{12.5}\selectfont\color{white}\MakeUppercase{Passe à OnBuch+}}\par\vspace{4pt}
        {\raggedright\popsemi\fontsize{8.4}{11}\selectfont\color{white}Tous les cours signés, Tuteur IA illimité, fiches PDF — onglet OnBuch+ de l'app\par}
      \end{minipage}
    \end{tcolorbox}
  \end{minipage}
  \vspace{8pt}
  \popcontenu
  \vfill
  \signatureonbuch
  \restoregeometry
  \newpage
}

% Rangée « Ce cours contient » (page de garde)
\newcommand{\poptile}[3]{% #1 couleur #2 gros texte #3 légende
  \begin{tcolorbox}[enhanced,colback=#1,colframe=popink,boxrule=2pt,arc=9pt,shadow={2.5pt}{-2.5pt}{0pt}{popink},
      left=6pt,right=6pt,top=9pt,bottom=9pt,halign=center,valign=center,height=1.75cm,width=\linewidth,nobeforeafter]
    {\popdisplay\fontsize{12.5}{13}\selectfont\color{white}\MakeUppercase{#2}}\par\vspace{3pt}
    {\centering\popsemi\fontsize{7.8}{9.5}\selectfont\color{white}#3\par}
  \end{tcolorbox}}
\newcommand{\popcontenu}{%
  \par\noindent{\popxboldls\fontsize{7.5}{7.5}\selectfont\color{popmuted}CE COURS SIGNÉ CONTIENT}\par\vspace{7pt}
  \noindent\begin{minipage}[t]{0.235\linewidth}\poptile{popblue}{Cours}{complet, illustré, pas à pas}\end{minipage}\hfill
  \begin{minipage}[t]{0.235\linewidth}\poptile{poporange}{Méthodes}{et exercices résolus}\end{minipage}\hfill
  \begin{minipage}[t]{0.235\linewidth}\poptile{poppink}{Exercices}{type examen + corrigés}\end{minipage}\hfill
  \begin{minipage}[t]{0.235\linewidth}\poptile{popgreen}{Fiche bilan}{l'essentiel à retenir}\end{minipage}\par}

% Sommaire
\newtcolorbox{sommaire}{enhanced,colback=white,colframe=popink,boxrule=2.2pt,arc=10pt,
  drop shadow=popink,left=18pt,right=18pt,top=13pt,bottom=13pt,before skip=6pt,after skip=20pt,
  before upper={\poppill{Sommaire}{poppurple}{white}\par\vspace{9pt}\popsemi\fontsize{10.6}{14.5}\selectfont\color{popink}}}

% Signature de fin de cours — poussée tout en bas de la dernière page
\newcommand{\findecours}{%
  \par\vfill\nopagebreak
  \begin{center}\popsquiggle{poporange}\end{center}\vspace{4pt}
  \signatureonbuch
}
\AtBeginDocument{\def\llbracket{\mathopen{[\mkern-3.5mu[}}\def\rrbracket{\mathclose{]\mkern-3.5mu]}}} % intervalles d'entiers (glyphe ⟦ absent de la police)
% Glyphes absents de Fira Math : redéfinis à partir de glyphes disponibles
\AtBeginDocument{%
  \def\setminus{\mathbin{\backslash}}\let\smallsetminus\setminus
  \def\vdots{\mathord{\vbox{\baselineskip3.2pt\lineskiplimit0pt\kern4pt\hbox{.}\hbox{.}\hbox{.}}}}%
  \def\ddots{\mathinner{\mkern1mu\raise6.5pt\hbox{.}\mkern2mu\raise3.6pt\hbox{.}\mkern2mu\raise0.7pt\hbox{.}\mkern1mu}}%
  \def\triangle{\mathord{\scalebox{0.9}{$\symup{\Delta}$}}}%
  \def\nabla{\mathord{\raisebox{\depth}{\scalebox{1}[-1]{$\symup{\Delta}$}}}}%
  \def\checkmark{\popcheck{popdark}}%
  \def\star{\mathbin{\ast}}\def\bigstar{\mathord{\ast}}%
  \def\diamond{\mathbin{\scalebox{0.6}{$\Diamond$}}}%
  \def\bigcup{\mathop{\mathchoice{\raisebox{-0.3em}{\scalebox{1.7}{$\cup$}}}{\scalebox{1.25}{$\cup$}}{\cup}{\cup}}\displaylimits}%
  \def\bigcap{\mathop{\mathchoice{\raisebox{-0.3em}{\scalebox{1.7}{$\cap$}}}{\scalebox{1.25}{$\cap$}}{\cap}{\cap}}\displaylimits}%
  \def\lesseqqgtr{\mathrel{\lessgtr}}\def\bigoplus{\mathop{\oplus}\displaylimits}%
}
\providecommand{\ssi}{si et seulement si} % abréviation fréquente
\AtBeginDocument{\providecommand{\wideparen}{\overparen}} % arc de cercle

%% FIGFIX-BEGIN v4 — correctifs globaux des figures (docs/onbuch-generation/tools/figfix)
\usepackage{adjustbox}\usepackage{etoolbox}
% toute figure TikZ/pgfplots plus large que la ligne est réduite à la largeur disponible
\BeforeBeginEnvironment{tikzpicture}{\begin{adjustbox}{max width=\linewidth}}
\AfterEndEnvironment{tikzpicture}{\end{adjustbox}}
% légendes pgfplots lisibles par-dessus les courbes
\pgfplotsset{every axis/.append style={legend style={fill=white,fill opacity=0.92,text opacity=1,draw=black!15,inner sep=3pt}}}
% étiquettes d'arêtes (syntaxe edge["..."]) avec fond blanc
\tikzset{every edge quotes/.append style={fill=white,inner sep=1.2pt}}
%% FIGFIX-END
\begin{document}

\pagegarde{Sensibilisation sur le rôle joué par les végétaux verts à travers la photosynthèse au sein de l'environnement}{SVT}{1ère D}{Module 1 : Le Monde Vivant}{ND06}{10 h}{Cette leçon étudie les conditions, les mécanismes et le bilan global de la photosynthèse chez les végétaux chlorophylliens. Elle montre ensuite le rôle fondamental des végétaux verts dans les cycles du carbone et de l'oxygène et les conséquences de la déforestation sur les équilibres écologiques.
\vspace{4pt}
\begin{multicols}{2}\small
\begin{itemize}
\item À la fin de cette leçon, je sais définir la photosynthèse et citer ses conditions nécessaires
\item À la fin de cette leçon, je sais décrire la structure d'un chloroplaste et localiser les pigments photosynthétiques
\item À la fin de cette leçon, je sais exploiter des résultats expérimentaux mettant en évidence les échanges gazeux photosynthétiques
\item À la fin de cette leçon, je sais écrire et interpréter l'équation-bilan globale de la photosynthèse
\item À la fin de cette leçon, je sais expliquer le rôle des végétaux verts dans les cycles du carbone et de l'oxygène
\item À la fin de cette leçon, je sais argumenter sur les conséquences de la déforestation pour l'équilibre de l'environnement
\end{itemize}\end{multicols}}

\begin{sommaire}
\begin{multicols}{2}
\begin{enumerate}
\item Mise en évidence expérimentale de la photosynthèse
\item Le siège de la photosynthèse : le chloroplaste et les pigments
\item Bilan global de la photosynthèse
\item Les végétaux verts et les cycles du carbone et de l'oxygène
\item Déforestation et déséquilibres environnementaux
\item Activité d'intégration
\item Méthodes et exercices résolus
\item Exercices
\item Corrigés des exercices
\item Fiche bilan
\end{enumerate}
\end{multicols}
\end{sommaire}

\begin{objectifs}
\begin{itemize}
\item À la fin de cette leçon, je sais définir la photosynthèse et citer ses conditions nécessaires
\item À la fin de cette leçon, je sais décrire la structure d'un chloroplaste et localiser les pigments photosynthétiques
\item À la fin de cette leçon, je sais exploiter des résultats expérimentaux mettant en évidence les échanges gazeux photosynthétiques
\item À la fin de cette leçon, je sais écrire et interpréter l'équation-bilan globale de la photosynthèse
\item À la fin de cette leçon, je sais expliquer le rôle des végétaux verts dans les cycles du carbone et de l'oxygène
\item À la fin de cette leçon, je sais argumenter sur les conséquences de la déforestation pour l'équilibre de l'environnement
\item À la fin de cette leçon, je sais interpréter un graphique de dégagement d'O₂ ou d'absorption de CO₂ en fonction de l'éclairement
\end{itemize}
\end{objectifs}

\begin{prerequis}
\begin{itemize}
\item Notion de métabolisme et de nutrition chez les êtres vivants (classe de Seconde)
\item Structure générale de la cellule végétale (paroi, vacuole, chloroplastes)
\item Notions de gaz de l'air : dioxygène O₂, dioxyde de carbone CO₂
\item Écriture et équilibrage d'une équation chimique simple
\item Lecture et exploitation de tableaux de résultats et de courbes
\end{itemize}
\end{prerequis}

\newpage

\coursec{Mise en évidence expérimentale de la photosynthèse}

À Yaoundé comme à Douala, tout le monde l'a remarqué : les quartiers très boisés, comme les abords du mont Fébé, restent plus frais et l'air y semble plus « respirable » que dans les zones densément urbanisées et bitumées. Les journaux parlent de la forêt du bassin du Congo comme du « poumon » de l'Afrique centrale. Un élève de Première D se pose alors des questions légitimes : pourquoi une forêt rafraîchit-elle et purifie-t-elle l'air ? Que se passerait-il si les coupes de bois et l'extension des champs de cacao continuaient de réduire nos forêts ? Pour répondre, il faut d'abord comprendre \emph{comment} les végétaux verts fabriquent leur propre matière organique et modifient la composition de l'atmosphère. Ce phénomène porte un nom : la \cle{photosynthèse}. Dans cette section, nous allons la mettre en évidence expérimentalement, étape par étape, comme le ferait un chercheur au laboratoire.

\courssub{Qu'est-ce que la photosynthèse ?}

\textbf{Observation de départ.} Une plante en pot placée dans une chambre éclairée grandit, produit de nouvelles feuilles, grossit. Pourtant, si l'on pèse le sol du pot au début et après plusieurs mois, sa masse n'a presque pas diminué. D'où vient donc la matière de la plante ? Pas (ou très peu) du sol. L'expérience historique de Van Helmont (saule planté dans un pot pesé pendant 5 ans) avait déjà montré que l'essentiel de la matière végétale ne provient pas de la terre. Elle provient de l'\cle{air} et de l'\cle{eau}, transformés grâce à l'énergie de la \cle{lumière}.

\begin{definition}[La photosynthèse]
La \cle{photosynthèse} est l'ensemble des réactions par lesquelles les végétaux \cle{chlorophylliens} (plantes vertes, algues) fabriquent de la \cle{matière organique} (des glucides, notamment du glucose puis de l'amidon) à partir de \cle{dioxyde de carbone} (\ce{CO2}) prélevé dans l'air et d'\cle{eau} (\ce{H2O}) absorbée par les racines, en utilisant l'énergie lumineuse captée par la \cle{chlorophylle}. Cette fabrication s'accompagne d'un rejet de \cle{dioxygène} (\ce{O2}) dans l'atmosphère.
\end{definition}

\begin{attention}[La plante respire aussi, jour et nuit !]
Erreur très fréquente au Probatoire : croire que la plante « fait de la photosynthèse le jour et respire la nuit ». C'est faux. La \cle{respiration cellulaire} (consommation d'\ce{O2}, rejet de \ce{CO2}) a lieu \textbf{jour et nuit}, dans toutes les cellules vivantes de la plante. La photosynthèse, elle, ne se déroule \textbf{qu'à la lumière}. Le jour, l'intensité de la photosynthèse dépasse largement celle de la respiration : le bilan net est donc une absorption de \ce{CO2} et un rejet d'\ce{O2}. La nuit, seule la respiration fonctionne : la plante rejette alors du \ce{CO2}.
\end{attention}

\courssub{La lumière est-elle indispensable ? L'expérience du cache opaque}

\textbf{Hypothèse.} Si la lumière est nécessaire à la fabrication de la matière organique, alors une partie de feuille privée de lumière ne devrait pas contenir d'amidon, contrairement à la partie éclairée.

\begin{experience}[Mise en évidence du rôle de la lumière]
\textbf{Matériel :} une plante en pot (géranium par exemple), papier d'aluminium ou cache opaque, trombones, alcool à \SI{90}{\degree}, bain-marie, \cle{eau iodée}, bécher, pince.

\textbf{Protocole :}
\begin{enumerate}
\item Placer la plante à l'obscurité pendant 48 h (désamidonnage : la plante consomme ses réserves d'amidon).
\item Fixer sur une feuille encore attachée à la plante un cache opaque (papier d'aluminium) qui masque une partie de la feuille.
\item Exposer la plante à la lumière pendant quelques heures.
\item Prélever la feuille et réaliser le test à l'eau iodée (voir méthode ci-dessous).
\end{enumerate}

\textbf{Observations :} après révélation, la partie éclairée de la feuille se colore en \textbf{bleu-violacé} (présence d'amidon), tandis que la partie restée sous le cache garde la teinte jaune-brun de l'eau iodée (absence d'amidon).

\textbf{Interprétation :} seule la zone ayant reçu la lumière a fabriqué de l'amidon. \textbf{Conclusion :} la lumière est une condition \emph{nécessaire} à la photosynthèse.
\end{experience}

\begin{methode}[Réaliser le test à l'eau iodée (technique exigible)]
\begin{enumerate}
\item \textbf{Décoloration :} plonger la feuille dans l'alcool à \SI{90}{\degree} chauffé \emph{au bain-marie} (jamais à la flamme : l'alcool est inflammable !). La chlorophylle se dissout dans l'alcool qui devient vert ; la feuille blanchit.
\item \textbf{Rinçage :} rincer délicatement la feuille à l'eau tiède pour éliminer l'alcool (qui durcit la feuille) et la réhydrater.
\item \textbf{Révélation :} étaler la feuille dans une boîte de Pétri et ajouter quelques gouttes d'\cle{eau iodée}.
\item \textbf{Lecture :} les zones contenant de l'amidon se colorent en \textbf{bleu-violacé} (parfois décrit comme bleu-noir) ; les zones sans amidon restent jaune-brun.
\end{enumerate}
\end{methode}

\begin{savaistu}[L'eau iodée, le révélateur star du Probatoire]
Le test à l'\cle{eau iodée} est \emph{le} révélateur classique de la photosynthèse dans les sujets du Probatoire D. L'iode de l'eau iodée s'insère dans l'hélice formée par les chaînes de glucose de l'\cle{amidon}, ce qui produit la couleur bleu-violacé caractéristique. Retiens bien : eau iodée + amidon $\rightarrow$ bleu-violacé. Ce même test sert d'ailleurs en cuisine et en laboratoire à détecter l'amidon dans les aliments (patate, manioc, igname : test positif immédiat !).
\end{savaistu}

\courssub{La chlorophylle est-elle indispensable ? L'expérience de la feuille panachée}

\textbf{Observation de départ.} Certaines plantes d'ornement (géranium panaché, chlorophytum, dracaena) possèdent des feuilles \cle{panachées} : vertes au centre, blanches ou jaunâtres sur les bords, car ces zones blanches sont dépourvues de chlorophylle.

\textbf{Hypothèse.} Si la chlorophylle est nécessaire, seules les zones vertes fabriqueront de l'amidon.

\textbf{Protocole et résultats.} On expose une feuille panachée à la lumière pendant quelques heures, on la dessine \emph{avant} le test (pour repérer les zones vertes et blanches), puis on la décolore à l'alcool bouillant et on la teste à l'eau iodée.

\textbf{Observations :} la carte des zones bleu-violacé (amidon présent) se superpose \emph{exactement} à la carte des zones vertes dessinée avant le test. Les zones blanches restent jaune-brun.

\textbf{Conclusion :} l'amidon ne se forme que dans les parties de la feuille contenant de la \cle{chlorophylle}. Le pigment vert est donc indispensable : il capte l'énergie lumineuse. C'est pourquoi on dit que la photosynthèse est le fait des végétaux \emph{chlorophylliens}.

\courssub{La plante éclairée dégage-t-elle un gaz ? L'expérience de l'élodée}

\textbf{Observation de départ.} Dans un aquarium bien éclairé, on voit parfois de petites bulles se former sur les feuilles des plantes aquatiques et monter à la surface. Quel est ce gaz ?

\begin{experience}[Dégagement gazeux d'une plante aquatique éclairée]
\textbf{Matériel :} un rameau d'\cle{élodée} (ou d'hydrille, plantes aquatiques), un entonnoir, un tube à essai, un cristallisoir rempli d'eau (additionnée d'un peu de \ce{CO2} ou d'eau gazeuse), une source lumineuse, une buchette en braise.

\textbf{Protocole :} placer l'élodée sous un entonnoir renversé dans le cristallisoir ; coiffer l'entonnoir d'un tube à essai rempli d'eau et renversé ; éclairer fortement le dispositif. Réaliser un témoin identique placé à l'obscurité.

\textbf{Observations :} à la lumière, des bulles gazeuses se dégagent des tiges et feuilles, montent et s'accumulent au sommet du tube, chassant l'eau. À l'obscurité, aucun dégagement notable. Lorsqu'on approche une buchette en braise du gaz recueilli, la braise \textbf{se ravive} (elle peut même se rallumer).

\textbf{Interprétation :} le gaz qui ravive une braise est le \cle{dioxygène} (\ce{O2}). \textbf{Conclusion :} la plante verte éclairée rejette du dioxygène, produit de la photosynthèse.
\end{experience}

\begin{popfigure}
\begin{center}
\begin{tikzpicture}[scale=0.95, every node/.style={font=\small}]
% cristallisoir
\draw[thick, popblue] (-3,0) -- (-3,3.2) (3,0) -- (3,3.2) (-3,0) -- (3,0);
\draw[popblue!50, fill=popblue!15] (-2.9,0.05) rectangle (2.9,2.4);
\node[popblue] at (2.1,0.35) {eau};
% entonnoir
\draw[thick, popdark] (-1.6,2.2) -- (-0.35,3.4) -- (-0.35,4.6) (1.6,2.2) -- (0.35,3.4) -- (0.35,4.6);
% elodée
\draw[thick, popgreen] (-0.9,0.3) .. controls (-0.7,1.2) and (-1.1,1.8) .. (-0.6,2.5);
\draw[thick, popgreen] (0.2,0.3) .. controls (0.4,1.3) and (0,1.9) .. (0.3,2.6);
\draw[popgreen, fill=popgreenL] (-0.75,1.2) ellipse (0.28 and 0.09);
\draw[popgreen, fill=popgreenL] (0.25,1.6) ellipse (0.28 and 0.09);
\draw[popgreen, fill=popgreenL] (-0.55,2.0) ellipse (0.24 and 0.08);
% tube à essai renversé
\draw[thick, popdark] (-0.55,3.4) -- (-0.55,6.2) arc (180:0:0.55) -- (0.55,3.4);
% gaz accumulé
\draw[poporange!60, fill=poporangeL] (-0.5,5.3) rectangle (0.5,6.05);
\draw[poporange!60] (-0.5,5.3) -- (0.5,5.3);
\node[popdark] at (0,5.65) {\footnotesize \ce{O2}};
% bulles
\foreach \x/\y/\r in {-0.3/2.8/0.07, 0.1/3.1/0.09, -0.15/3.6/0.07, 0.2/4.0/0.1, -0.1/4.5/0.08, 0.15/4.9/0.09}
  \draw[popblue] (\x,\y) circle (\r);
% lumière
\draw[thick, popgold, ->] (-4.2,5.5) -- (-3.0,4.8);
\draw[thick, popgold, ->] (-4.2,4.3) -- (-3.0,3.9);
\node[popgold] at (-4.3,5.9) {lumière};
% légendes pointées
\draw[thin, popmuted] (0.55,5.6) -- (2.6,5.9) node[right, popdark] {tube à essai renversé};
\draw[thin, popmuted] (-0.6,2.4) -- (-3.4,2.9) node[left, popdark] {élodée};
\draw[thin, popmuted] (1.2,2.5) -- (3.6,2.9) node[right, popdark] {entonnoir};
\draw[thin, popmuted] (0.2,4.2) -- (2.8,4.5) node[right, popdark] {bulles d'\ce{O2}};
\end{tikzpicture}
\end{center}
\legende{Dispositif de mise en évidence du dégagement de dioxygène par une plante aquatique (élodée) éclairée. Les bulles s'accumulent au sommet du tube renversé ; le gaz recueilli ravive une braise : c'est du \ce{O2}.}
\end{popfigure}

\courssub{La plante éclairée absorbe-t-elle du dioxyde de carbone ?}

\textbf{Hypothèse.} Si le \ce{CO2} est une matière première de la photosynthèse, sa teneur doit diminuer dans le milieu entourant une plante éclairée.

\begin{experience}[Absorption du \ce{CO2} par la plante éclairée]
On enferme une plante verte dans une enceinte contenant de l'eau additionnée de \cle{BTH} (bleu de bromothymol), indicateur qui est \textbf{jaune} en présence de \ce{CO2} (milieu acide) et \textbf{bleu} en son absence. On prépare deux enceintes identiques : l'une exposée à la lumière, l'autre placée à l'obscurité.

\textbf{Observations :} à la lumière, le BTH jaune redevient progressivement \textbf{bleu} : le \ce{CO2} a disparu du milieu. À l'obscurité, le BTH reste jaune (voire jaunit davantage, car la respiration de la plante rejette du \ce{CO2}).

\textbf{Conclusion :} la plante éclairée \emph{absorbe} le dioxyde de carbone du milieu. Le \ce{CO2} est bien une matière première de la photosynthèse. (On peut obtenir le même résultat avec l'\cle{eau de chaux}, qui se trouble en présence de \ce{CO2} : autour de la plante éclairée, l'eau de chaux reste limpide.)
\end{experience}

\courssub{Quels facteurs influencent l'intensité de la photosynthèse ?}

Mesurons maintenant l'\cle{intensité photosynthétique}, c'est-à-dire la quantité de matière organique fabriquée par unité de temps, que l'on peut estimer par le volume d'\ce{O2} dégagé (ou le nombre de bulles) par une élodée par minute.

\textbf{Résultats expérimentaux.} Lorsqu'on fait varier l'\cle{éclairement} (intensité lumineuse) en rapprochant ou éloignant la lampe, on obtient la courbe suivante :

\begin{popfigure}
\begin{center}
\begin{tikzpicture}
\begin{axis}[
  width=0.85\linewidth, height=6.2cm,
  xlabel={Éclairement (lux)},
  ylabel={Intensité photosynthétique (mL d'\ce{O2}/h)},
  xmin=0, xmax=10500, ymin=0, ymax=13,
  xtick={0,2000,4000,6000,8000,10000},
  ytick={0,2,4,6,8,10,12},
  grid=both, grid style={popline!40},
  axis lines=left, thick,
  legend style={at={(0.97,0.25)}, anchor=east, font=\small}
]
\addplot[popcurve, domain=0:10000, samples=200] {11*(1-exp(-x/2200))};
\addlegendentry{plante éclairée (\ce{CO2} suffisant)}
% point de compensation
\addplot[only marks, mark=*, poporange, mark size=2pt] coordinates {(250,1.18)};
\draw[dashed, poporange] (axis cs:250,0) -- (axis cs:250,1.18) -- (axis cs:0,1.18);
\node[poporange, font=\footnotesize, anchor=west] at (axis cs:400,0.6) {point de compensation};
% plateau
\draw[dashed, popmuted] (axis cs:0,11) -- (axis cs:10000,11);
\node[popmuted, font=\footnotesize, anchor=south east] at (axis cs:9900,11) {plateau (saturation)};
\end{axis}
\end{tikzpicture}
\end{center}
\legende{Intensité photosynthétique en fonction de l'éclairement. La courbe croît d'abord presque linéairement, puis atteint un plateau : au-delà d'un certain éclairement, un autre facteur (le \ce{CO2}, la température) devient limitant.}
\end{popfigure}

\textbf{Interprétation de la courbe.}
\begin{itemize}
\item Pour les faibles éclairements, l'intensité photosynthétique augmente proportionnellement à la lumière : la lumière est alors le \cle{facteur limitant}.
\item Au-delà d'un certain éclairement, la courbe plafonne (\emph{plateau de saturation}) : la lumière n'est plus limitante ; c'est un autre facteur qui bloque, par exemple la \cle{concentration en \ce{CO2}} de l'air (seulement 0,04 \% environ) ou la température. Si l'on enrichit alors l'atmosphère en \ce{CO2}, le plateau s'élève : c'est le principe utilisé dans les \emph{serres} où l'on diffuse du \ce{CO2} pour augmenter les rendements.
\item Au \cle{point de compensation}, l'\ce{O2} produit par la photosynthèse compense exactement celui consommé par la respiration : le bilan gazeux net est nul.
\end{itemize}

\begin{definition}[Facteur limitant]
Un \cle{facteur limitant} est un facteur écologique (lumière, \ce{CO2}, température, eau) dont la valeur, trop faible, empêche l'intensité photosynthétique d'augmenter, même si tous les autres facteurs sont favorables. L'intensité de la photosynthèse est fixée par le facteur le plus défavorable.
\end{definition}

\begin{exemplebox}[Volume d'\ce{O2} dégagé par une élodée]
Une élodée est placée dans le dispositif précédent pendant 30 min. Sous un éclairement faible (lampe à 60 cm), on recueille \SI{2,5}{mL} de gaz ; sous un éclairement fort (lampe à 20 cm), on recueille \SI{7,5}{mL}. Le volume d'\ce{O2} dégagé par heure passe donc de \SI{5}{mL/h} à \SI{15}{mL/h} : multiplier l'éclairement a multiplié par trois l'intensité photosynthétique, tant que la lumière reste le facteur limitant.
\end{exemplebox}

\begin{exoresolu}[Exploiter des résultats expérimentaux]
\textbf{Énoncé.} Une élève de 1ère D place un rameau d'élodée dans le dispositif à entonnoir et compte les bulles dégagées par minute sous trois conditions : (A) obscurité ; (B) lumière faible ; (C) lumière forte. Elle obtient : A = 0 bulle/min ; B = 12 bulles/min ; C = 35 bulles/min. Le gaz recueilli en C ravive une braise. 1) Identifier le gaz dégagé. 2) Comparer A et C et conclure. 3) Comparer B et C et conclure.
\tcblower
\textbf{Solution détaillée.}
\begin{enumerate}
\item Le seul gaz courant qui ravive une braise est le \cle{dioxygène} \ce{O2}. Le gaz dégagé par l'élodée éclairée est donc du dioxygène.
\item En A (obscurité), aucune bulle ; en C (lumière forte), 35 bulles/min. Toutes les autres conditions étant identiques, seule la présence de lumière varie : on conclut que le dégagement d'\ce{O2}, donc la photosynthèse, \emph{nécessite la lumière}.
\item En passant de B (lumière faible) à C (lumière forte), le dégagement passe de 12 à 35 bulles/min, soit une multiplication par environ 2,9. On conclut que, dans cet intervalle, \emph{l'intensité photosynthétique augmente avec l'éclairement} : la lumière y est le facteur limitant.
\end{enumerate}
\end{exoresolu}

\courssub{Conclusion : les conditions, les produits et le bilan global de la photosynthèse}

L'ensemble de ces expériences converge vers un bilan clair, qu'il faut savoir restituer sans hésiter :

\begin{aretenir}[Conditions, produits et bilan global de la photosynthèse]
\textbf{Conditions nécessaires (matières premières et facteurs) :}
\begin{itemize}
\item la \cle{lumière} (expérience du cache opaque : pas d'amidon à l'ombre) ;
\item la \cle{chlorophylle} (expérience de la feuille panachée : pas d'amidon dans les zones blanches) ;
\item le \cle{dioxyde de carbone} \ce{CO2} (expérience au BTH : il est absorbé à la lumière) ;
\item l'\cle{eau} \ce{H2O}, absorbée par les racines et transportée jusqu'aux feuilles.
\end{itemize}
\textbf{Produits :}
\begin{itemize}
\item de la \cle{matière organique} (glucose \ce{C6H12O6}, stocké sous forme d'\cle{amidon}, révélé en bleu-violacé par l'eau iodée) ;
\item du \cle{dioxygène} \ce{O2}, rejeté dans l'atmosphère (gaz qui ravive une braise).
\end{itemize}
\textbf{Bilan global (équation chimique) :}
\[
\ce{6 CO2 + 6 H2O ->[$\text{lumière}$][\text{chlorophylle}] C6H12O6 + 6 O2}
\]
\end{aretenir}

\begin{attention}[Pièges de rédaction au Probatoire]
\begin{itemize}
\item Ne jamais écrire que la plante « produit de l'oxygène la nuit » : à l'obscurité, elle \emph{consomme} de l'\ce{O2} (respiration).
\item Ne pas confondre les révélateurs : l'\cle{eau iodée} révèle l'\emph{amidon} (bleu-violacé) ; le \cle{BTH} ou l'\cle{eau de chaux} révèlent le \ce{CO2} ; la \emph{braise ravivée} identifie l'\ce{O2}.
\item Dans l'expérience du cache opaque, ne pas oublier le \emph{désamidonnage} préalable (48 h à l'obscurité) : sans lui, l'amidon déjà présent fausse le test.
\item La décoloration se fait à l'alcool \emph{chauffé au bain-marie}, jamais à la flamme directe.
\item L'équation bilane s'écrit avec des coefficients stœchiométriques entiers (6, 6, 1, 6) et les conditions au-dessus de la flèche.
\end{itemize}
\end{attention}

Nous savons désormais \emph{que} la photosynthèse a lieu, \emph{dans quelles conditions} et \emph{quel est son bilan global}. Reste à découvrir \emph{où} elle se déroule dans la cellule : c'est l'objet de la section suivante, consacrée au chloroplaste et aux pigments photosynthétiques.

\coursec{Le siège de la photosynthèse : le chloroplaste et les pigments}

Dans la section précédente, les expériences historiques (Van Helmont, Priestley, Ingenhousz, Sachs) nous ont appris un fait capital : seules les parties \textbf{vertes} des végétaux, et seulement \textbf{à la lumière}, dégagent du dioxygène et fabriquent de la matière organique. Une question surgit alors naturellement : \emph{pourquoi précisément les parties vertes ?} Qu'y a-t-il dans une feuille verte qui n'existe pas dans une racine blanche ? Pour répondre, il faut descendre à l'échelle de la cellule et y découvrir l'usine photosynthétique : le \cle{chloroplaste}, et ses molécules « pièges à lumière » : les \cle{pigments photosynthétiques}.

\courssub{Localisation des cellules photosynthétiques dans la feuille}

\textbf{Observation.} Coupe une feuille de manioc ou de cocotier : la face supérieure est vert foncé, la face inférieure vert pâle. Sous le microscope optique, une coupe transversale de feuille révèle que cette différence de couleur correspond à une différence de structure interne.

\begin{experience}[Observation d'une coupe transversale de feuille]
\textbf{Matériel :} lame mince de coupe transversale de feuille (dicotylédone), microscope optique.

\textbf{Observations :} on distingue, de haut en bas :
\begin{itemize}
\item l'\textbf{épiderme supérieur}, recouvert d'une cuticule transparente qui laisse passer la lumière ;
\item le \textbf{parenchyme palissadique} : une ou deux assises de cellules allongées, serrées, \emph{très riches en chloroplastes} (c'est lui qui donne le vert foncé de la face supérieure) ;
\item le \textbf{parenchyme lacuneux} (ou spongieux) : cellules arrondies, moins riches en chloroplastes, séparées par de grands espaces d'air (méats) qui facilitent la diffusion des gaz ($\ce{CO2}$, $\ce{O2}$, vapeur d'eau) ;
\item l'\textbf{épiderme inférieur}, percé de \textbf{stomates}, petits orifices qui assurent les échanges gazeux avec l'atmosphère.
\end{itemize}

\textbf{Interprétation :} les chloroplastes ne sont pas répartis au hasard : ils sont concentrés dans les cellules les mieux exposées à la lumière. La feuille est un organe \emph{adapté} à la photosynthèse : grande surface, faible épaisseur, transparence de l'épiderme, réseau de méats pour l'approvisionnement en $\ce{CO2}$.

\textbf{Conclusion :} la photosynthèse se déroule dans les \cle{cellules chlorophylliennes} du parenchyme palissadique et du parenchyme lacuneux, et plus précisément à l'intérieur d'un organite : le chloroplaste.
\end{experience}

\begin{definition}[Chloroplaste]
Le \cle{chloroplaste} est un \textbf{organite} (structure intracellulaire spécialisée) présent dans le cytoplasme des cellules chlorophylliennes des végétaux verts. C'est le \textbf{siège de la photosynthèse} : il capte l'énergie lumineuse et l'utilise pour fabriquer de la matière organique (glucides) à partir de $\ce{CO2}$ et d'eau. Une cellule du parenchyme palissadique peut contenir plusieurs dizaines de chloroplastes.
\end{definition}

\courssub{Structure du chloroplaste}

Au microscope électronique, le chloroplaste apparaît comme une petite lentille ovale de 5 à \SI{10}{\micro m} de long, dont l'organisation interne est remarquablement ordonnée.

\begin{popfigure}
\begin{center}
\begin{tikzpicture}[scale=1.0]
% Enveloppe externe
\draw[thick,popgreen,fill=popgreenL!40] (0,0) ellipse (4.6 and 2.6);
\draw[thick,popgreen!70!black] (0,0) ellipse (4.3 and 2.35);
% Stroma dots
\foreach \x/\y in {-3.4/1.2,-2.6/-1.4,-1.8/1.6,-0.6/-1.7,0.6/1.8,1.8/-1.5,2.8/1.3,3.5/-0.9,-3.6/-0.6,3.7/0.7,0.1/-0.9,-1.2/0.9,1.4/0.8}
  \fill[popgreen!60!black] (\x,\y) circle (0.045);
% Granum 1
\begin{scope}[shift={(-2.2,0.3)}]
\foreach \i in {0,...,5}{\draw[thick,popdark,fill=popgoldL] (0,\i*0.18) ellipse (0.85 and 0.09);}
\end{scope}
% Granum 2
\begin{scope}[shift={(0.4,-0.4)}]
\foreach \i in {0,...,5}{\draw[thick,popdark,fill=popgoldL] (0,\i*0.18) ellipse (0.85 and 0.09);}
\end{scope}
% Granum 3
\begin{scope}[shift={(2.6,0.5)}]
\foreach \i in {0,...,4}{\draw[thick,popdark,fill=popgoldL] (0,\i*0.18) ellipse (0.75 and 0.09);}
\end{scope}
% Lamelles intergranaires
\draw[thick,popdark] (-1.35,0.75) -- (-0.45,0.15);
\draw[thick,popdark] (1.25,0.05) -- (1.85,0.85);
% Légendes
\draw[thin,popmuted] (4.3,1.4) -- (5.6,2.2) node[right,align=left]{\textbf{Double membrane}\\ (enveloppe)};
\draw[thin,popmuted] (3.9,-1.9) -- (5.6,-2.4) node[right,align=left]{\textbf{Stroma}\\ (gel interne)};
\draw[thin,popmuted] (-2.2,1.45) -- (-2.2,2.6) node[above,align=center]{\textbf{Granum}\\ (empilement de thylakoïdes)};
\draw[thin,popmuted] (0.4,0.75) -- (0.4,2.6) node[above,align=center]{\textbf{Thylakoïde}\\ (sac aplati)};
\draw[thin,popmuted] (-0.9,0.45) -- (-4.9,-2.3) node[below,align=center]{\textbf{Lamelle}\\ intergranaire};
\end{tikzpicture}
\end{center}
\legende{Structure du chloroplaste vue au microscope électronique : une double membrane délimite un gel interne, le stroma, dans lequel baignent des sachets aplatis, les thylakoïdes, empilés en granums et reliés entre eux par des lamelles intergranaires.}
\end{popfigure}

\begin{aretenir}[Les trois niveaux d'organisation du chloroplaste]
\begin{enumerate}
\item Une \textbf{double membrane} (enveloppe) isole le chloroplaste du cytoplasme et contrôle les échanges.
\item Le \textbf{stroma} : gel interne riche en enzymes, en ribosomes et en ADN chloroplastique ; c'est le lieu des réactions chimiques de synthèse des glucides (cycle de Calvin).
\item Les \textbf{thylakoïdes} : sachets membranaires aplatis, empilés en colonnes appelées \textbf{granums} (singulier : granum). C'est dans les \emph{membranes des thylakoïdes} que se trouvent les pigments photosynthétiques et que s'effectue la capture de la lumière.
\end{enumerate}
\end{aretenir}

\begin{propriete}[Relation structure-fonction (admise)]
L'empilement des thylakoïdes en granums multiplie considérablement la \textbf{surface membranaire} disponible dans un faible volume. Or les pigments sont insérés dans ces membranes : plus la surface membranaire est grande, plus la quantité de pigments — et donc la quantité de lumière captée — est élevée. \textbf{La structure du chloroplaste optimise la capture de l'énergie lumineuse.}
\end{propriete}

\courssub{Les pigments photosynthétiques}

Pourquoi la feuille est-elle verte ? Parce qu'elle contient des molécules colorées capables d'absorber certaines radiations lumineuses : les pigments.

\begin{definition}[Pigments photosynthétiques]
Les \cle{pigments photosynthétiques} sont des molécules colorées, localisées dans les \textbf{membranes des thylakoïdes}, capables d'\textbf{absorber l'énergie lumineuse} et de la convertir en énergie chimique. On distingue :
\begin{itemize}
\item les \textbf{chlorophylles} : \cle{chlorophylle a} (vert bleuté) et \cle{chlorophylle b} (vert jaunâtre), pigments \emph{principaux} ;
\item les \textbf{caroténoïdes} (carotène orange, xanthophylle jaune), pigments \emph{accessoires} qui absorbent d'autres radiations et transfèrent l'énergie à la chlorophylle a.
\end{itemize}
\end{definition}

\textbf{Comment savoir quelles couleurs un pigment absorbe ?} On éclaire une solution de chlorophylle extraite de feuilles avec de la lumière blanche, puis on mesure, pour chaque longueur d'onde, la proportion de lumière absorbée : on obtient le \cle{spectre d'absorption}.

\begin{popfigure}
\begin{center}
\begin{tikzpicture}
\begin{axis}[
width=\linewidth, height=7cm,
xlabel={Longueur d'onde $\lambda$ (nm)},
ylabel={Absorption (\%)},
xmin=380, xmax=750, ymin=0, ymax=100,
xtick={400,450,500,550,600,650,700},
ytick={0,20,40,60,80,100},
grid=both, grid style={popline!40},
legend style={at={(0.5,0.98)},anchor=north,draw=popline,fill=popcream},
axis background/.style={fill=white},
]
\addplot[very thick,popblue,smooth] coordinates {(400,62)(420,78)(430,90)(440,86)(450,70)(470,40)(500,18)(530,12)(550,10)(570,14)(590,22)(610,35)(630,55)(650,72)(662,82)(675,70)(690,35)(710,8)(730,2)};
\addlegendentry{Chlorophylle a}
\addplot[very thick,popgreen!70!black,smooth] coordinates {(400,45)(420,60)(440,72)(455,80)(465,74)(480,50)(500,28)(530,16)(550,12)(580,18)(600,30)(620,45)(640,58)(645,62)(660,50)(680,22)(700,6)(720,2)};
\addlegendentry{Chlorophylle b}
\end{axis}
\end{tikzpicture}
\end{center}
\legende{Spectres d'absorption des chlorophylles a et b (allure simplifiée). Les deux pigments absorbent fortement dans le bleu (vers 430--460 nm) et dans le rouge (vers 640--665 nm), mais très peu dans le vert (vers 500--600 nm) : la lumière verte est réfléchie, d'où la couleur des feuilles.}
\end{popfigure}

\begin{propriete}[Absorption de la lumière par la chlorophylle (admise)]
La chlorophylle absorbe principalement la lumière \textbf{rouge} et la lumière \textbf{bleue}, et absorbe très peu la lumière \textbf{verte}, qu'elle \textbf{réfléchit}. C'est cette lumière verte réfléchie qui atteint notre œil : \textbf{les feuilles nous paraissent vertes parce qu'elles n'utilisent pas le vert}. L'énergie des radiations rouges et bleues absorbées est utilisée pour réaliser la photosynthèse.
\end{propriete}

\begin{attention}[Chloroplaste ou chlorophylle : ne pas confondre !]
Erreur très fréquente au Probatoire :
\begin{itemize}
\item le \textbf{chloroplaste} est un \emph{organite} (une structure de la cellule), le \emph{lieu} de la photosynthèse ;
\item la \textbf{chlorophylle} est un \emph{pigment} (une molécule) contenue \emph{dans} le chloroplaste, qui capte la lumière.
\end{itemize}
On ne dit donc jamais « la photosynthèse a lieu dans la chlorophylle » : elle a lieu \emph{dans le chloroplaste}, \emph{grâce à} la chlorophylle. De même, on écrit « cellules chlorophylliennes » (qui contiennent de la chlorophylle) et non « cellules chloroplastiques ».
\end{attention}

\courssub{Mise en évidence des pigments : la chromatographie sur papier}

Une feuille verte contient-elle un seul pigment ? Pour le savoir, on réalise une \cle{chromatographie sur papier}, technique de séparation des molécules d'un mélange.

\begin{experience}[Chromatographie des pigments d'une feuille]
\textbf{Matériel :} feuilles fraîches (épinard, manioc), mortier, alcool (éthanol), papier filtre, solvant (mélange éther de pétrole / acétone), cuve.

\textbf{Protocole :}
\begin{enumerate}
\item Broyer des feuilles dans l'alcool : on obtient un \textbf{extrait chlorophyllien} vert foncé.
\item Déposer une petite tache concentrée d'extrait sur une ligne tracée au crayon, à \SI{2}{cm} du bas d'une bande de papier filtre (la \textbf{ligne de dépôt}).
\item Suspendre la bande dans la cuve de sorte que le bas du papier trempe dans le solvant, \emph{sans que la tache touche le solvant}.
\item Laisser le solvant monter par capillarité, puis retirer et marquer immédiatement le \textbf{front du solvant}.
\end{enumerate}

\textbf{Observations :} la tache verte initiale s'est séparée en plusieurs taches colorées, de bas en haut : chlorophylle b (vert jaunâtre), chlorophylle a (vert bleuté), xanthophylle (jaune), carotène (jaune-orangé, tout en haut).

\textbf{Interprétation :} chaque pigment migre à une vitesse qui dépend de sa solubilité dans le solvant et de son affinité pour le papier : plus un pigment est soluble et peu retenu par le papier, plus il monte haut. L'extrait de feuille contient donc \textbf{plusieurs pigments}.

\textbf{Conclusion :} la couleur verte des feuilles résulte du mélange de plusieurs pigments, dont les chlorophylles a et b (dominantes) et des caroténoïdes.
\end{experience}

\begin{methode}[Lire et exploiter un chromatogramme de pigments]
\begin{enumerate}
\item Repérer la \textbf{ligne de dépôt} (point de départ) et le \textbf{front du solvant} (niveau maximal atteint par le solvant).
\item Identifier chaque \textbf{tache} par sa couleur et sa position : du bas vers le haut : chlorophylle b, chlorophylle a, xanthophylle, carotène.
\item Pour chaque tache, mesurer la distance $d$ parcourue par le pigment et la distance $D$ parcourue par le solvant, depuis la ligne de dépôt.
\item Calculer le \textbf{rapport frontal} $R_f$ :
\[ R_f = \frac{d}{D} \qquad (0 < R_f < 1) \]
\item Comparer les $R_f$ à des valeurs de référence pour identifier les pigments : le carotène a le $R_f$ le plus élevé (il migre le plus loin), la chlorophylle b le plus faible.
\end{enumerate}
\end{methode}

\begin{exoresolu}[Calcul d'un rapport frontal]
\textbf{Énoncé.} Lors d'une chromatographie d'extrait de feuilles de manioc, le front du solvant a parcouru \SI{12}{cm} depuis la ligne de dépôt. Une tache vert bleuté a migré de \SI{7,2}{cm}. (1) Calculer le rapport frontal de cette tache. (2) Identifier le pigment correspondant.

\tcblower
\textbf{Solution.}
(1) Par définition :
\[ R_f = \frac{d}{D} = \frac{7{,}2}{12} = 0{,}6. \]
(2) La tache est vert bleuté et présente un rapport frontal intermédiaire : il s'agit de la \textbf{chlorophylle a}. (Rappel : la chlorophylle b, vert jaunâtre, migre moins loin ; le carotène migre le plus haut.)
\end{exoresolu}

\begin{savaistu}[Pourquoi les feuilles jaunissent-elles ?]
En saison sèche, en automne sous les climats tempérés, ou chez une plante carencée, la chlorophylle, molécule fragile, se \textbf{dégrade} plus vite qu'elle n'est renouvelée. Les \textbf{caroténoïdes}, plus stables, étaient présents depuis le début mais masqués par le vert dominant de la chlorophylle : ils deviennent alors visibles et la feuille \textbf{jaunit} ou vire à l'orange. C'est aussi ce qu'on observe sur les feuilles de maïs carencées en magnésium dans certaines parcelles épuisées : le magnésium est un atome \emph{constitutif} de la molécule de chlorophylle ; sans lui, la plante ne peut plus en fabriquer et pâlit — signe d'alerte bien connu des agronomes camerounais.
\end{savaistu}

\begin{exercice}[Pour vérifier tes acquis]
(1) Recopie et complète : « La photosynthèse se déroule dans le \ldots, organite délimité par une \ldots\ membrane, contenant un gel appelé \ldots\ et des sachets aplatis, les \ldots, empilés en \ldots. »
(2) Sur un chromatogramme, une tache a migré de \SI{4,8}{cm} et le solvant de \SI{12}{cm}. Calcule son $R_f$ et propose une identification.
(3) Explique, en deux phrases, pourquoi une feuille éclairée en lumière verte uniquement réalise une photosynthèse très faible.
\end{exercice}

\begin{corrige}[Exercice n°1]
(1) « La photosynthèse se déroule dans le \textbf{chloroplaste}, organite délimité par une \textbf{double} membrane, contenant un gel appelé \textbf{stroma} et des sachets aplatis, les \textbf{thylakoïdes}, empilés en \textbf{granums}. »

(2) $R_f = \dfrac{4{,}8}{12} = 0{,}4$. Rapport frontal faible et position basse : il s'agit vraisemblablement de la \textbf{chlorophylle b} (à confirmer par la couleur vert jaunâtre de la tache).

(3) La chlorophylle absorbe très peu la lumière verte (elle la réfléchit) : éclairée uniquement en vert, la feuille capte très peu d'énergie lumineuse. La photosynthèse, qui dépend de l'énergie lumineuse absorbée, est donc très ralentie.
\end{corrige}

\begin{aretenir}[L'essentiel de la section]
\begin{itemize}
\item La photosynthèse a lieu dans les cellules chlorophylliennes des \textbf{parenchymes palissadique et lacuneux} de la feuille, au sein d'un organite : le \textbf{chloroplaste}.
\item Le chloroplaste comprend une \textbf{double membrane}, un \textbf{stroma} et des \textbf{thylakoïdes} empilés en \textbf{granums} ; la grande surface des membranes des thylakoïdes optimise la capture de la lumière.
\item Les \textbf{pigments} (chlorophylles a et b, caroténoïdes), logés dans les membranes des thylakoïdes, absorbent l'énergie lumineuse, surtout dans le \textbf{rouge} et le \textbf{bleu} ; ils réfléchissent le \textbf{vert}, d'où la couleur des feuilles.
\item La \textbf{chromatographie sur papier} sépare les pigments ; on les identifie par leur couleur et leur rapport frontal $R_f = d/D$.
\end{itemize}
\end{aretenir}

Maintenant que nous connaissons l'« usine » (le chloroplaste) et ses « capteurs » (les pigments), il reste à comprendre ce que cette usine produit et consomme : c'est l'objet de la section suivante, consacrée au \emph{bilan global de la photosynthèse}.

\coursec{Bilan global de la photosynthèse}

Dans la section précédente, nous avons découvert \emph{où} se déroule la photosynthèse : dans le \cle{chloroplaste}, organite dont les pigments (chlorophylles et caroténoïdes) captent l'énergie lumineuse. Nous savons aussi, grâce aux expériences de mise en évidence, \emph{quels} gaz sont échangés : le végétal vert éclairé absorbe du dioxyde de carbone et rejette du dioxygène. Il reste à répondre à une question fondamentale : que deviennent ces gaz ? Quelle matière le végétal fabrique-t-il, et comment résumer cette transformation en une seule écriture chimique ? C'est l'objet de cette section : établir, équilibrer et interpréter le \cle{bilan global de la photosynthèse}.

\courssub{De l'observation à l'équation-bilan}

\textbf{Situation concrète.} Un maïs planté en saison des pluies à Bafoussam ne reçoit, pendant toute sa croissance, que de l'eau, de l'air et de la lumière. Pourtant, en trois mois, il produit des tiges, des feuilles et des épis représentant plusieurs kilogrammes de matière sèche. D'où vient cette matière ? Ni du sol (la masse du sol du pot ne diminue presque pas), ni de l'eau seule. C'est la question que s'est posée le médecin et chimiste flamand \cle{Van Helmont} au XVII\textsuperscript{e} siècle : ayant fait pousser un saule dans un pot pendant cinq ans en ne l'arrosant que d'eau, il constata que l'arbre avait gagné plus de 74 kg alors que le sol n'avait perdu qu'environ 60 g. Il en conclut — à tort — que toute la matière venait de l'eau. Les travaux ultérieurs (Priestley, Ingenhousz, puis les expériences modernes) ont permis de trancher : la matière du végétal est construite essentiellement à partir du \ce{CO2} de l'air, l'eau fournissant l'hydrogène et le dioxygène dégagé.

\textbf{Démarche d'investigation.}
\begin{itemize}
\item \emph{Observation :} une plante verte éclairée, placée en présence de \ce{CO2}, produit de la matière organique (l'amidon mis en évidence par l'eau iodée dans les feuilles, qui vire au bleu-noir) et dégage du \ce{O2} (gaz qui ravive une braise).
\item \emph{Hypothèse :} le végétal transforme le \ce{CO2} et l'eau, grâce à l'énergie lumineuse, en matière organique et en dioxygène.
\item \emph{Vérification expérimentale :} en privant la plante de \ce{CO2} (la potasse, KOH, absorbe le \ce{CO2} de l'enceinte), la production d'amidon cesse ; en mesurant les gaz échangés, on constate que le volume de \ce{O2} dégagé est égal au volume de \ce{CO2} absorbé.
\item \emph{Conclusion :} la photosynthèse est une transformation chimique dont on peut écrire l'équation-bilan.
\end{itemize}

\begin{propriete}[Équation-bilan de la photosynthèse]
La photosynthèse est l'ensemble des réactions par lesquelles un végétal chlorophyllien éclairé synthétise de la matière organique (glucose) à partir de dioxyde de carbone et d'eau, avec dégagement de dioxygène. Son bilan global s'écrit :
\[ \ce{6CO2 + 6H2O ->[$\text{lumière}$][\text{chlorophylle}] C6H12O6 + 6O2} \]
Les conditions de la réaction (\emph{lumière} et \emph{chlorophylle}) s'écrivent au-dessus et au-dessous de la flèche, et non parmi les réactifs : ce ne sont ni des réactifs ni des produits. Cette écriture est exigible au Probatoire, avec les coefficients exacts.
\end{propriete}

\begin{experience}[Démonstration guidée : vérification de l'équilibrage atome par atome]
Une équation-bilan correcte doit respecter la \cle{conservation des atomes} : chaque atome présent dans les réactifs se retrouve dans les produits. Vérifions pour \ce{6CO2 + 6H2O -> C6H12O6 + 6O2}.
\begin{enumerate}
\item \textbf{Carbone (C) :} à gauche, $6 \times 1 = 6$ atomes (dans 6 \ce{CO2}) ; à droite, 6 atomes (dans \ce{C6H12O6}). Équilibré.
\item \textbf{Hydrogène (H) :} à gauche, $6 \times 2 = 12$ atomes (dans 6 \ce{H2O}) ; à droite, 12 atomes (dans \ce{C6H12O6}). Équilibré.
\item \textbf{Oxygène (O) :} à gauche, $6 \times 2 + 6 \times 1 = 12 + 6 = 18$ atomes ; à droite, 6 atomes dans \ce{C6H12O6} et $6 \times 2 = 12$ atomes dans 6 \ce{O2}, soit $6 + 12 = 18$. Équilibré.
\end{enumerate}
Les trois éléments sont conservés : l'équation est correctement équilibrée.
\end{experience}

\begin{methode}[Écrire et équilibrer l'équation-bilan au Probatoire]
\begin{enumerate}
\item Écrire d'abord les espèces en présence sans coefficients : \ce{CO2 + H2O -> C6H12O6 + O2}, en précisant « lumière » et « chlorophylle » sur la flèche.
\item Équilibrer le \textbf{carbone} : 6 atomes dans le glucose $\Rightarrow$ coefficient 6 devant \ce{CO2}.
\item Équilibrer l'\textbf{hydrogène} : 12 atomes dans le glucose $\Rightarrow$ coefficient 6 devant \ce{H2O}.
\item Équilibrer l'\textbf{oxygène} en dernier : 18 atomes O à gauche ; le glucose en contient 6, il en reste 12 $\Rightarrow$ coefficient 6 devant \ce{O2}.
\item Vérifier atome par atome (C, H, O) avant de conclure.
\end{enumerate}
\end{methode}

\begin{attention}[Erreur classique : la provenance du dioxygène dégagé]
Beaucoup d'élèves écrivent que le \ce{O2} dégagé par la plante « provient du \ce{CO2} ». C'est \textbf{faux}. Les expériences utilisant l'isotope marqué \cle{$^{18}$O} (isotope lourd de l'oxygène, utilisé comme traceur) l'ont démontré : si l'on fournit à la plante de l'eau marquée \ce{H2^{18}O} et du \ce{CO2} non marqué, le dioxygène dégagé est du \ce{^{18}O2}. Le \ce{O2} dégagé provient donc de la \textbf{photolyse de l'eau} (du grec \emph{photo} : lumière, et \emph{lysis} : rupture — c'est la rupture de la molécule d'eau sous l'effet de la lumière). Le \ce{CO2}, lui, fournit le carbone et une partie de l'oxygène du glucose.
\end{attention}

\courssub{Interprétation de l'équation-bilan : fixation du carbone et conversion d'énergie}

\begin{definition}[Fixation du carbone]
Le carbone du glucose synthétisé provient exclusivement du \ce{CO2} atmosphérique (ou dissous dans l'eau pour les plantes aquatiques). On dit que la photosynthèse réalise la \cle{fixation du carbone minéral} (le carbone inorganique du \ce{CO2}) en \cle{carbone organique} (le carbone des molécules vivantes). C'est la porte d'entrée du carbone dans le monde vivant : toute la matière organique des écosystèmes, y compris notre propre nourriture, découle de cette fixation.
\end{definition}

L'équation-bilan a aussi une signification énergétique. La réaction inverse (la dégradation du glucose lors de la respiration) libère de l'énergie ; la photosynthèse en \emph{consomme} donc. Cette énergie est fournie par la \textbf{lumière}. Ainsi :

\begin{aretenir}[Double signification du bilan photosynthétique]
\begin{itemize}
\item \textbf{Bilan de matière :} de la matière minérale (\ce{CO2}, \ce{H2O}) est transformée en matière organique (\ce{C6H12O6}), avec rejet de \ce{O2}.
\item \textbf{Bilan énergétique :} l'\cle{énergie lumineuse} captée par les pigments est convertie en \cle{énergie chimique}, stockée dans les liaisons des molécules organiques. Le glucose est un « réservoir » d'énergie utilisable ensuite par la respiration cellulaire.
\end{itemize}
\end{aretenir}

\begin{popfigure}
\begin{center}
\begin{tikzpicture}[scale=1.0, every node/.style={font=\small}]
% Feuille stylisée
\draw[thick, popgreen, fill=popgreenL] (0,0) .. controls (1.2,1.4) and (3.2,1.6) .. (4.6,0)
   .. controls (3.2,-1.6) and (1.2,-1.4) .. (0,0) -- cycle;
\draw[popgreen!70!black, thick] (0,0) -- (4.6,0);
\draw[popgreen!70!black] (1.2,0) -- (1.7,0.55);
\draw[popgreen!70!black] (2.3,0) -- (2.9,0.7);
\draw[popgreen!70!black] (3.4,0) -- (3.9,0.5);
\draw[popgreen!70!black] (1.2,0) -- (1.7,-0.55);
\draw[popgreen!70!black] (2.3,0) -- (2.9,-0.7);
\draw[popgreen!70!black] (3.4,0) -- (3.9,-0.5);
% Pétiole
\draw[thick, popgreen!70!black] (0,0) -- (-1.0,0);
\node[below, popgreen!50!black] at (2.3,-1.3) {\textbf{Feuille (chloroplastes)}};
% Entrée lumière
\draw[-{Stealth[length=3mm]}, very thick, popgold] (2.3,3.0) -- (2.3,1.2);
\node[popgold!60!black, above] at (2.3,3.0) {\textbf{Lumière}};
% Entrée CO2
\draw[-{Stealth[length=3mm]}, very thick, popblue] (-3.4,1.2) -- (-0.6,0.5);
\node[popblue, left] at (-3.4,1.2) {\textbf{\ce{CO2} absorbé}};
% Entrée H2O
\draw[-{Stealth[length=3mm]}, very thick, poppurple] (-1.0,-2.2) -- (-0.7,-0.2);
\node[poppurple, below left] at (-1.0,-2.2) {\textbf{\ce{H2O} (par les racines)}};
% Sortie O2
\draw[-{Stealth[length=3mm]}, very thick, poporange] (4.6,0.5) -- (7.0,1.4);
\node[poporange!80!black, right] at (7.0,1.4) {\textbf{\ce{O2} dégagé}};
% Sortie glucose
\draw[-{Stealth[length=3mm]}, very thick, poppink] (4.6,-0.5) -- (7.0,-1.4);
\node[poppink!80!black, right] at (7.0,-1.4) {\textbf{Glucose \ce{C6H12O6}}};
\end{tikzpicture}
\end{center}
\legende{Schéma récapitulatif des échanges de la feuille photosynthétique : entrées (lumière, \ce{CO2} atmosphérique pénétrant par les stomates, eau absorbée par les racines et transportée par la sève brute) et sorties (dioxygène rejeté, glucose utilisé par la plante ou stocké sous forme d'amidon).}
\end{popfigure}

\courssub{Les deux phases de la photosynthèse (complément utile au Probatoire)}

L'équation-bilan résume le résultat global, mais la photosynthèse se déroule en réalité en \textbf{deux phases} localisées dans des régions différentes du chloroplaste (voir section 2).

\begin{definition}[Phase photochimique et phase biochimique]
\begin{itemize}
\item La \cle{phase photochimique} (ou phase claire) se déroule dans les \textbf{thylakoïdes} (empilements de sacs aplatis formant les grana). Elle nécessite directement la lumière : les pigments captent l'énergie lumineuse, qui sert à réaliser la \textbf{photolyse de l'eau} ($\rightarrow$ dégagement de \ce{O2}) et à produire l'énergie chimique (molécules d'\cle{ATP}) et le pouvoir réducteur (\cle{NADPH}) nécessaires à la suite.
\item La \cle{phase biochimique} (ou phase sombre) se déroule dans le \textbf{stroma} (le « gel » entourant les grana). Elle n'utilise pas directement la lumière : grâce à l'ATP et au NADPH fournis par la phase claire, des enzymes y fixent le \ce{CO2} et le transforment en glucose (ensemble de réactions appelé \cle{cycle de Calvin}).
\end{itemize}
\end{definition}

\begin{savaistu}[La phase « sombre » ne se déroule pas dans l'obscurité !]
Le mot « sombre » est trompeur : il signifie seulement que cette phase \emph{ne consomme pas directement la lumière}. En réalité, la phase biochimique se déroule \textbf{en même temps} que la phase claire, donc en plein jour, car elle utilise aussitôt l'ATP et le NADPH produits par les thylakoïdes. Les deux phases sont indissociables : sans lumière, la phase claire s'arrête, et la phase sombre s'arrête presque aussitôt faute d'énergie.
\end{savaistu}

\courssub{Application quantitative : calcul de la masse de \ce{CO2} absorbée}

L'équation-bilan permet de faire des calculs stœchiométriques simples, très appréciés au Probatoire.

\begin{exoresolu}[Masse de \ce{CO2} absorbée pour produire 180 g de glucose]
\textbf{Énoncé.} Un végétal vert synthétise par photosynthèse 180 g de glucose. On donne $M(\ce{CO2}) = 44$ g/mol et $M(\ce{C6H12O6}) = 180$ g/mol. Calculer la masse de \ce{CO2} absorbée et le volume de \ce{O2} dégagé dans les conditions normales de température et de pression ($V_m = 22{,}4$ L/mol).
\tcblower
\textbf{Solution.} D'après l'équation-bilan : $\ce{6CO2 + 6H2O -> C6H12O6 + 6O2}$, il faut 6 moles de \ce{CO2} pour produire 1 mole de glucose, et 6 moles de \ce{O2} sont alors dégagées.

\textbf{1. Quantité de glucose :}
\[ n(\text{glucose}) = \frac{m}{M} = \frac{180}{180} = 1 \text{ mol} \]

\textbf{2. Quantité de \ce{CO2} absorbée :} d'après les coefficients, $n(\ce{CO2}) = 6 \times n(\text{glucose}) = 6$ mol.

\textbf{3. Masse de \ce{CO2} :}
\[ m(\ce{CO2}) = n \times M = 6 \times 44 = 264 \text{ g} \]

\textbf{4. Volume de \ce{O2} dégagé :} $n(\ce{O2}) = 6$ mol, donc :
\[ V(\ce{O2}) = n \times V_m = 6 \times 22{,}4 = 134{,}4 \text{ L} \]

\textbf{Conclusion :} pour fabriquer 180 g de glucose, la plante a absorbé \textbf{264 g de \ce{CO2}} et rejeté \textbf{134,4 L de dioxygène}. Ce calcul illustre concrètement le rôle « purificateur d'air » des végétaux verts.
\end{exoresolu}

\courssub{Comparaison photosynthèse et respiration cellulaire}

La respiration cellulaire, que tous les êtres vivants (y compris les plantes, jour et nuit !) réalisent, est la transformation quasi inverse de la photosynthèse :
\[ \ce{C6H12O6 + 6O2 -> 6CO2 + 6H2O} \quad \text{(respiration)} \]

\begin{center}
\begin{tabularx}{\linewidth}{|l|X|X|}
\hline
\rowcolor{popblueL} \textbf{Critère} & \textbf{Photosynthèse} & \textbf{Respiration cellulaire} \\
\hline
Équation-bilan & $\ce{6CO2 + 6H2O -> C6H12O6 + 6O2}$ & $\ce{C6H12O6 + 6O2 -> 6CO2 + 6H2O}$ \\
\hline
Gaz absorbé & \ce{CO2} & \ce{O2} \\
\hline
Gaz dégagé & \ce{O2} & \ce{CO2} \\
\hline
Matière organique & Synthétisée (anabolisme) & Dégradée (catabolisme) \\
\hline
Énergie & Lumineuse $\rightarrow$ chimique (stockée) & Chimique $\rightarrow$ utilisable (ATP) \\
\hline
Moment (chez la plante) & Uniquement à la lumière & Jour \textbf{et} nuit \\
\hline
Organite & Chloroplaste & Mitochondrie (et cytoplasme) \\
\hline
Qui la réalise ? & Végétaux chlorophylliens, algues, cyanobactéries & Tous les êtres vivants \\
\hline
\end{tabularx}
\end{center}

\begin{attention}[Piège fréquent]
Ne jamais écrire que « la plante respire le jour et photosynthétise la nuit » : c'est l'inverse pour la photosynthèse, et la respiration, elle, a lieu \textbf{en permanence}. Le jour, la photosynthèse est plus intense que la respiration : le bilan net de la plante éclairée est donc une absorption de \ce{CO2} et un dégagement de \ce{O2}. La nuit, seule la respiration se manifeste : la plante absorbe alors du \ce{O2} et rejette du \ce{CO2}.
\end{attention}

\begin{exercice}[Pour s'entraîner]
Une feuille absorbe en une journée 88 g de \ce{CO2}. On donne $M(\ce{CO2}) = 44$ g/mol et $M(\ce{C6H12O6}) = 180$ g/mol.
\begin{enumerate}
\item Calculer la quantité de matière de \ce{CO2} absorbée.
\item En déduire la masse de glucose synthétisée et le volume de \ce{O2} dégagé ($V_m = 22{,}4$ L/mol).
\end{enumerate}
\end{exercice}

\begin{corrige}[Exercice : masse de glucose et volume de \ce{O2}]
\textbf{1.} $n(\ce{CO2}) = \dfrac{88}{44} = 2$ mol.

\textbf{2.} D'après l'équation-bilan, 6 mol de \ce{CO2} donnent 1 mol de glucose et 6 mol de \ce{O2}. Donc :
\[ n(\text{glucose}) = \frac{2}{6} = \frac{1}{3} \text{ mol} \quad \Rightarrow \quad m(\text{glucose}) = \frac{1}{3} \times 180 = 60 \text{ g} \]
\[ n(\ce{O2}) = n(\ce{CO2}) = 2 \text{ mol} \quad \Rightarrow \quad V(\ce{O2}) = 2 \times 22{,}4 = 44{,}8 \text{ L} \]
La feuille a synthétisé \textbf{60 g de glucose} et dégagé \textbf{44,8 L de dioxygène}.
\end{corrige}

\begin{aretenir}[L'essentiel de la section]
\begin{itemize}
\item Bilan global : $\ce{6CO2 + 6H2O ->[$\text{lumière}$][\text{chlorophylle}] C6H12O6 + 6O2}$, équilibré en C, H et O.
\item Le \ce{O2} dégagé provient de \textbf{l'eau} (photolyse), le carbone du glucose provient du \textbf{\ce{CO2}} : c'est la fixation du carbone.
\item La photosynthèse convertit l'énergie lumineuse en énergie chimique stockée dans la matière organique.
\item Deux phases : photochimique (thylakoïdes, lumière indispensable, production d'ATP et de NADPH) et biochimique (stroma, fixation du \ce{CO2} par le cycle de Calvin).
\item Photosynthèse et respiration sont complémentaires : leurs échanges gazeux sont inverses, ce qui assure l'équilibre gazeux de la biosphère — équilibre que nous étudierons dans la section suivante.
\end{itemize}
\end{aretenir}

\coursec{Les végétaux verts et les cycles du carbone et de l'oxygène}

Dans la section précédente, nous avons établi le \textbf{bilan global de la photosynthèse} :
\ce{6CO2 + 6H2O ->[$\text{lumière}$][\text{chlorophylles}] C6H12O6 + 6O2}.
Cette équation résume une prouesse biologique : le végétal vert capte l'énergie lumineuse pour transformer du gaz carbonique (CO₂) et de l'eau en matière organique (glucose) et en dioxygène (O₂). Mais cette réaction ne se produit pas dans un tube à essai isolé ; elle s'inscrit dans le fonctionnement de la \textbf{biosphère}. Que deviennent le CO₂ prélevé et l'O₂ rejeté ? Comment les végétaux, par cette activité, régulent-ils la composition de l'atmosphère et assurent-ils la survie des êtres hétérotrophes (animaux, champignons, la plupart des bactéries) ? C'est l'objet de cette section : comprendre les \textbf{cycles biogéochimiques du carbone et de l'oxygène} et le rôle central des producteurs primaires.

\courssub{Le cycle du carbone : un grand cycle planétaire}

\begin{experience}[Mise en évidence des échanges de CO₂ entre un végétal et l'atmosphère]
\textbf{Problème :} Le végétal vert prélève-t-il du CO₂ en permanence ? Le rejette-t-il ?
\textbf{Matériel :} Deux bocaux hermétiques identiques, une plante verte (ex. \emph{Chlorophytum} ou jeune plant de maïs), de l'eau de chaux (solution d'hydroxyde de calcium Ca(OH)₂), une source lumineuse, un tissu opaque.
\textbf{Protocole :}
\begin{enumerate}
    \item Placer la plante dans le bocal A (lumière) et le bocal B (obscurité totale grâce au tissu).
    \item Laisser 2 à 3 heures.
    \item Prélever un échantillon d'air dans chaque bocal (seringue) et le faire bouillir dans un tube contenant de l'eau de chaux.
\end{enumerate}
\textbf{Observations :}
\begin{itemize}
    \item Bocal A (lumière) : l'eau de chaux reste \textbf{claire} (peu ou pas de trouble).
    \item Bocal B (obscurité) : l'eau de chaux \textbf{trouble} (formation de carbonate de calcium CaCO₃).
\end{itemize}
\textbf{Interprétation :} L'eau de chaux trouble en présence de CO₂ (\ce{Ca(OH)2 + CO2 -> CaCO3 v + H2O}).
\begin{itemize}
    \item Dans le noir, la plante ne photosynthétise pas ; elle \textbf{respire} : elle consomme de l'O₂ et \textbf{rejette du CO₂} (trouble).
    \item À la lumière, la plante photosynthétise \textbf{et} respire. La photosynthèse \textbf{consomme le CO₂} produit par la respiration et prélève celui de l'air du bocal. Le bilan net est une \textbf{consommation de CO₂} (eau de chaux claire).
\end{itemize}
\textbf{Conclusion :} Le végétal vert est un \textbf{puits de CO₂ net le jour} (photosynthèse > respiration) et une \textbf{source de CO₂ la nuit} (respiration seule).
\end{experience}

\begin{definition}[Cycle du carbone]
Le \cle{cycle du carbone} est l'ensemble des processus par lesquels les atomes de carbone circulent entre les différents \textbf{réservoirs} de la Terre : l'atmosphère (CO₂), l'hydrosphère (CO₂ dissous, bicarbonates), la lithosphère (roches carbonatées, combustibles fossiles) et la biosphère (êtres vivants, matière organique du sol).
\end{definition}

\begin{popfigure}
\begin{tikzpicture}[
    node distance=2.5cm and 3cm,
    reservoir/.style={draw=popdark, thick, fill=popblueL, rounded corners, minimum width=2.2cm, minimum height=1.2cm, align=center, font=\small},
    flux/.style={->, thick, popdark, font=\footnotesize, >=Stealth},
    label/.style={midway, fill=white, inner sep=1pt, font=\footnotesize, align=center}
]
% Reservoirs
\node[reservoir, align=center] (atm){Atmosphère\\ \ce{CO2} (\SI{880}{Gt C})};
\node[reservoir, right=of atm, align=center] (veg){Végétaux chlorophylliens\\ (Biomasse \SI{550}{Gt C})};
\node[reservoir, below=of veg, align=center] (sol){Sol / Matière organique\\ (\SI{1500}{Gt C})};
\node[reservoir, left=of sol] (anim) {Animaux / Hétérotrophes};
\node[reservoir, below=of atm, align=center] (ocean){Océan / Hydrosphère\\ (\SI{38000}{Gt C})};
\node[reservoir, right=of ocean, align=center] (fossil){Combustibles fossiles\\ (Charbon, pétrole, gaz)};

% Flux naturels
\draw[flux] (atm) -- node[label, above, align=center]{Photosynthèse\\ \SI{120}{Gt C/an}} (veg);
\draw[flux] (veg) -- node[label, right, align=center]{Respiration végétale\\ \SI{60}{Gt C/an}} (atm);
\draw[flux] (veg) -- node[label, left] {Mort / Chute feuilles} (sol);
\draw[flux] (anim) -- node[label, below, align=center]{Respiration animale\\ \SI{60}{Gt C/an}} (atm);
\draw[flux] (sol) -- node[label, left, align=center]{Décomposition\\ (Respiration microbienne)} (atm);
\draw[flux] (atm) -- node[label, left] {Dissolution} (ocean);
\draw[flux] (ocean) -- node[label, right] {Dégazage} (atm);
\draw[flux] (veg) -- node[label, above] {Alimentation} (anim);
\draw[flux] (anim) -- node[label, right] {Mort / Déjections} (sol);

% Flux anthropique (perturbation)
\draw[flux, poporange, dashed, very thick] (fossil) -- node[label, fill=poporangeL!50, align=center]{Combustion\\ \SI{9}{Gt C/an} (anthropique)} (atm);
\draw[flux, poporange, dashed] (sol) -- node[label, fill=poporangeL!50] {Déforestation / Labour} (atm);

% Légendes des flèches
\node[align=center, font=\tiny, text=popmuted, anchor=north] at (current bounding box.south) [yshift=-0.5cm] {Traits pleins : flux naturels équilibrés (pré-industriel).\\ Traits pointillés orange : flux anthropiques perturbateurs.};
\end{tikzpicture}
\legende{Cycle simplifié du carbone à l'échelle planétaire. Les valeurs en Gt C/an (Gigatonnes de carbone par an) sont des ordres de grandeur pour l'ère pré-industrielle (flux naturels) et actuels (flux anthropiques en orange). La photosynthèse est le flux entrant principal depuis l'atmosphère vers la biosphère vivante.}
\end{popfigure}

\begin{methode}[Lire et annoter un schéma de cycle biogéochimique]
Pour analyser un cycle (carbone, azote, eau, phosphore...) :
\begin{enumerate}
    \item \textbf{Identifiez les réservoirs (stocks) :} Ce sont les cases/ovales (ex. : Atmosphère, Végétaux, Sol, Océan). Notez leur unité (souvent Gt C ou ppm).
    \item \textbf{Identifiez les flux (transferts) :} Ce sont les flèches. Notez le \textbf{processus} (photosynthèse, respiration, combustion, dissolution...) et le \textbf{sens} (qui donne, qui reçoit).
    \item \textbf{Quantifiez si possible :} Repérez les valeurs numériques (Gt/an). Comparez l'amplitude des flèches entrantes et sortantes d'un réservoir.
    \item \textbf{Déduisez l'évolution du stock :} Si $\sum$ Entrées $>$ $\sum$ Sorties $\rightarrow$ le stock \textbf{augmente}. Si $\sum$ Entrées $<$ $\sum$ Sorties $\rightarrow$ le stock \textbf{diminue}. Si égalité $\rightarrow$ \textbf{équilibre dynamique (stationnaire)}.
    \item \textbf{Repérez les perturbations :} Cherchez les flèches d'origine humaine (souvent en pointillés ou couleur distincte) qui rompent l'équilibre.
\end{enumerate}
\end{methode}

\begin{propriete}[Équilibre dynamique pré-industriel du carbone atmosphérique]
Avant l'ère industrielle (vers 1750), le stock de carbone dans l'atmosphère était stable (\textbf{280 ppm} de CO₂). Les flux naturels d'entrée (respiration globale, décomposition, dégazage océanique, volcans) étaient \textbf{exactement compensés} par les flux de sortie (photosynthèse globale, dissolution océanique).
\[
\text{Photosynthèse globale} \approx \text{Respiration totale} + \text{Dégazage océanique} + \text{Émissions volcaniques}
\]
Cet équilibre n'est pas statique : d'immenses quantités circulent (\textbf{200 Gt C/an} échangés), mais le \textbf{stock atmosphérique ne variait que très lentement} (échelles géologiques).
\end{propriete}

\begin{exemplebox}[Le bassin du Congo : un puits de carbone vital pour le Cameroun et la planète]
La forêt dense humide du bassin du Congo (dont la partie camerounaise : Sud, Est, Littoral) est le \textbf{deuxième massif forestier tropical au monde} après l'Amazonie.
\begin{itemize}
    \item Elle stocke \textbf{30 à 40 milliards de tonnes de carbone} dans sa biomasse aérienne et ses sols (tourbières de la Cuvette Centrale).
    \item Chaque année, sa photosynthèse nette (NPP - Production Primaire Nette) prélève \textbf{1 à 1,5 Gt C} de l'atmosphère.
    \item \textbf{Rôle régulateur :} En stockant du carbone, elle freine l'augmentation du CO₂ atmosphérique et atténue le réchauffement climatique. Sa déforestation (exploitation forestière non durable, agriculture sur brûlis, agro-industrie) libère ce carbone stocké, transformant un \textbf{puits} en \textbf{source} nette de CO₂.
\end{itemize}
Protéger la forêt camerounaise (parcs nationaux comme Lobéké, Boumba Bek, Nki, Dja) est donc une action locale à impact climatique \textbf{global}.
\end{exemplebox}

\begin{definition}[Producteur primaire et Puits de carbone]
\begin{itemize}
    \item \textbf{Producteur primaire (autotrophe) :} Organisme capable de synthétiser sa matière organique à partir de matière minérale (CO₂, H₂O, sels minéraux) en utilisant une source d'énergie (lumière = \emph{photoautotrophe} ; réactions chimiques = \emph{chémoautotrophe}). Les végétaux chlorophylliens (terrestres + phytoplancton) sont les principaux producteurs primaires de la biosphère. Ils sont à la \textbf{base de toutes les chaînes alimentaires}.
    \item \textbf{Puits de carbone :} Réservoir (forêt, océan, sol) qui absorbe plus de carbone (sous forme de CO₂) qu'il n'en émet sur une période donnée, entraînant une \textbf{augmentation de son stock de carbone}. L'opposé est une \textbf{source de carbone}.
\end{itemize}
\end{definition}

\begin{attention}[Erreur fréquente : « Les plantes produisent de l'oxygène la nuit »]
\textbf{FAUX.} La nuit, \textbf{pas de lumière $\rightarrow$ pas de photosynthèse}. La plante ne fait que \textbf{respire} : elle consomme de l'O₂ et rejette du CO₂, exactement comme un animal.
\begin{itemize}
    \item \textbf{Jour :} Photosynthèse + Respiration. Bilan net souvent : \textbf{Consommation CO₂, Production O₂} (si lumière suffisante).
    \item \textbf{Nuit :} Respiration seule. Bilan : \textbf{Consommation O₂, Production CO₂}.
\end{itemize}
Au Probatoire, précisez toujours : « \emph{En présence de lumière, le bilan net de la photosynthèse est une production d'O₂} » et non « la plante produit de l'O₂ en permanence ».
\end{attention}

\courssub{Le cycle de l'oxygène : couplage étroit avec le carbone}

L'oxygène dioxygène (O₂) représente \textbf{21\%} de l'atmosphère actuelle (\textbf{1,2 $\times$ 10\textsuperscript{6} Gt O₂}). Ce taux élevé est une \textbf{anomalie géologique} : sans vie, l'O₂ réagirait vite avec les roches réduites (fer, soufre) et disparaîtrait de l'air. Sa présence massive est la \textbf{signature biologique} de la photosynthèse oxygénique.

\begin{popfigure}
\begin{tikzpicture}[
    node distance=2.5cm and 3cm,
    reservoir/.style={draw=popdark, thick, fill=popgreenL, rounded corners, minimum width=2.2cm, minimum height=1.2cm, align=center, font=\small},
    flux/.style={->, thick, popdark, font=\footnotesize, >=Stealth},
    label/.style={midway, fill=white, inner sep=1pt, font=\footnotesize, align=center}
]
% Reservoirs
\node[reservoir, align=center] (atm){Atmosphère\\ O₂ (\SI{21}{\%}, \SI{1.2e6}{Gt})};
\node[reservoir, right=of atm, align=center] (bio){Biosphère\\ (Végétaux, Animaux, Microbes)};
\node[reservoir, below=of atm, align=center] (litho){Lithosphère\\ (Oxydation roches, sédiments)};
\node[reservoir, below=of bio, align=center] (ocean){Océan\\ (O₂ dissous)};

% Flux
\draw[flux] (atm) -- node[label, above, align=center]{Respiration\\ (Végétaux, Animaux, Microbes)} (bio);
\draw[flux] (bio) -- node[label, right, align=center]{Photosynthèse\\ (Végétaux, Phytoplancton)} (atm);
\draw[flux] (atm) -- node[label, left] {Dissolution} (ocean);
\draw[flux] (ocean) -- node[label, right] {Dégazage / Photosynthèse phytoplancton} (atm);
\draw[flux] (atm) -- node[label, left, align=center]{Combustion\\ (Feux, Industrie, Respiration)} (litho);
\draw[flux, dashed, poporange] (litho) -- node[label, fill=poporangeL!50] {Altération / Sédimentation (très lent)} (atm);

\node[align=center, font=\tiny, text=popmuted, anchor=north] at (current bounding box.south) [yshift=-0.5cm] {Le cycle de l'O₂ est l'image quasi-miroir de celui du C.\\ La photosynthèse est la SEULE source majeure d'O₂ atmosphérique.};
\end{tikzpicture}
\legende{Cycle simplifié de l'oxygène. La photosynthèse (flèche verte vers le haut) est l'unique flux majeur injectant de l'O₂ dans l'atmosphère. Tous les autres flux (respiration, combustion, oxydation des roches) en sont des puits.}
\end{popfigure}

\begin{propriete}[Le phytoplancton : le « poumon bleu » de la planète (admis)]
Bien que la biomasse du phytoplancton marin soit bien plus faible que celle des forêts terrestres (\textbf{1 à 2 Gt C} contre \textbf{550 Gt C}), son \textbf{taux de renouvellement} est extrême (jours/semaines contre décennies/siècles pour les arbres).
\begin{itemize}
    \item Il assure \textbf{50 à 60\% de la production primaire mondiale} et de la production nette d'O₂ atmosphérique.
    \item Il réalise un \textbf{pompage biologique du carbone} : en fixant le CO₂ en surface et en exportant la matière organique vers les fonds (neige marine), il séquestre du carbone sur le long terme.
\end{itemize}
\emph{Protéger les océans (lutte contre l'acidification, la pollution, la surpêche qui déséquilibre les réseaux trophiques) est aussi crucial que protéger les forêts pour l'équilibre O₂/CO₂.}
\end{propriete}

\begin{savaistu}[Le lac Nyos (Cameroun) : un piège géologique au CO₂]
Le lac Nyos (région de l'Ouest, ligne volcanique du Cameroun) est un lac de cratère profond. Du CO₂ d'origine magmatique s'accumule au fond, dissous sous pression dans l'eau froide. En 1986, un renversement brutal (éruption limnique) a libéré \textbf{1,6 million de tonnes de CO₂} en quelques minutes. Ce gaz, plus lourd que l'air, a formé un nuage au ras du sol qui a asphyxié \textbf{1 746 personnes} et des milliers d'animaux dans les vallées en aval.
\begin{itemize}
    \item \textbf{Lien avec le cours :} Cet événement tragique illustre que le CO₂ est un gaz \textbf{incolore, inodore, plus dense que l'air} et \textbf{toxique/asphyxiant} à haute concentration.
    \item Des tuyaux de dégazage artificiels ont été installés pour prévenir une récidive en laissant le CO₂ s'échapper en continu et contrôlé vers l'atmosphère.
\end{itemize}
\end{savaistu}

\courssub{Les végétaux verts : base des réseaux trophiques et régulateurs climatiques}

\begin{exoresolu}[Analyse d'un bilan carbone forestier]
\textbf{Énoncé :} Une parcelle de forêt équatoriale camerounaise de 1 ha stocke 250 t de carbone dans sa biomasse aérienne (bois, feuilles). La production primaire brute (PPB) est de 30 t C/ha/an. La respiration autotrophe (Ra) de la forêt est de 15 t C/ha/an. La respiration hétérotophe (Rh, décomposeurs du sol) est de 14 t C/ha/an.
\begin{enumerate}
    \item Calculer la Production Primaire Nette (PPN = PPB - Ra) et la Production Nette d'Écosystème (PNE = PPN - Rh). Interpréter le signe de la PNE.
    \item Si cette parcelle est convertie en plantation de palmiers à huile (stock biomasse \textbf{50 t C/ha}, PPN \textbf{10 t C/ha/an}, Rh \textbf{8 t C/ha/an}), calculer la nouvelle PNE et la perte de stock de carbone initiale. Discuter l'impact climatique.
\end{enumerate}

\textbf{Solution détaillée :}
\tcblower
\begin{enumerate}
    \item \textbf{Forêt naturelle :}
    \begin{align*}
        \text{PPN} &= \text{PPB} - R_a = 30 - 15 = \mathbf{15\; t\; C/ha/an} \\
        \text{PNE} &= \text{PPN} - R_h = 15 - 14 = \mathbf{+1\; t\; C/ha/an}
    \end{align*}
    La PNE est \textbf{positive} : l'écosystème est un \textbf{puits net de carbone} (il séquestre 1 t C/ha/an dans le sol/biomasse). C'est typique d'une forêt mature en croissance lente ou en régénération.
    \item \textbf{Plantation de palmiers :}
    \begin{align*}
        \text{PNE}_{\text{palmier}} &= 10 - 8 = \mathbf{+2\; t\; C/ha/an}
    \end{align*}
    La plantation est aussi un puits net (même plus fort par an car croissance rapide), MAIS :
    \begin{itemize}
        \item \textbf{Perte de stock initiale :} $250 - 50 = \mathbf{200\; t\; C/ha}$ libérés brutalement dans l'atmosphère lors du défrichement (brûlis, décomposition des souches, oxydation du sol).
        \item \textbf{Temps de remboursement (carbon payback time) :} Pour compenser les 200 t C perdues avec un gain net de +2 t C/an (au lieu de +1 t C/an), il faut $200 / (2 - 1) = 200$ ans (si on compare au statu quo forêt) ou $200 / 2 = 100$ ans (si on compare à sol nu).
    \end{itemize}
    \textbf{Conclusion :} La conversion libère un « bombement » de carbone immédiat que la nouvelle culture met des décennies à siècles à réabsorber. C'est une \textbf{source nette massive à court/moyen terme}.
\end{enumerate}
\end{exoresolu}

\begin{aretenir}[Synthèse : Les végétaux verts, piliers des équilibres planétaires]
\begin{enumerate}
    \item \textbf{Cycle du carbone :} La photosynthèse est le \textbf{seul flux naturel majeur} transférant le carbone de l'atmosphère (CO₂ inerte) vers la biosphère (matière organique vivante). Respiration, décomposition, combustion (naturelle ou anthropique) ferment la boucle.
    \item \textbf{Cycle de l'oxygène :} La photosynthèse est la \textbf{seule source significative} d'O₂ atmosphérique. Toute respiration et combustion en sont des puits. L'équilibre O₂/CO₂ atmosphérique est le miroir de l'équilibre carbone.
    \item \textbf{Producteurs primaires :} Végétaux terrestres (forêts, savanes, cultures) + Phytoplancton marin. Base de \textbf{toutes} les chaînes alimentaires (hétérotrophes = consommateurs primaires, secondaires...).
    \item \textbf{Régulation climatique :} En stockant du C (puits de carbone), les forêts (bassin du Congo, Amazonie, Bornéo) et l'océan (pompe biologique + physique) limitent l'effet de serre. \textbf{La déforestation casse ce thermostat planétaire.}
    \item \textbf{Équilibre dynamique :} À l'échelle géologique, photosynthèse $\approx$ respiration + combustion naturelle. L'activité humaine (combustion fossile + changement d'affectation des sols) rompt cet équilibre depuis 1750 $\rightarrow$ hausse CO₂ (+50\%) $\rightarrow$ réchauffement.
\end{enumerate}
\end{aretenir}

\begin{exercice}[Application : Lire un cycle et argumenter]
\textbf{Document :} Le graphique ci-dessous (simulé) montre l'évolution saisonnière de la concentration en CO₂ à l'observatoire de Mauna Loa (Hawaï, hémisphère Nord) sur deux ans.
\begin{center}
\begin{tikzpicture}
\begin{axis}[
    width=\linewidth, height=6cm,
    xlabel={Temps (mois)}, ylabel={CO₂ (ppm)},
    xmin=0, xmax=24, ymin=410, ymax=425,
    xtick={0,3,6,9,12,15,18,21,24},
    xticklabels={Jan, Avr, Juil, Oct, Jan, Avr, Juil, Oct, Jan},
    ytick={410,415,420,425},
    grid=major, grid style={popline},
    axis line style={popdark},
    tick style={popdark},
    label style={font=\small},
    tick label style={font=\small},
]
\addplot[popcurve, domain=0:24, samples=300] {418 + 3*sin(deg(2*pi*x/12 - pi/2)) + 0.1*x};
\addplot[poporange, thick, dashed, domain=0:24] {418 + 0.1*x};
\end{axis}
\end{tikzpicture}
\end{center}
\legende{Variation saisonnière du CO₂ à Mauna Loa (courbe continue) et tendance à long terme (courbe pointillée orange).}
\begin{enumerate}
    \item Décrire l'allure de la courbe continue (oscillations + tendance).
    \item Expliquer l'oscillation annuelle par l'activité photosynthétique de la végétation de l'hémisphère Nord (forêts boréales, tempérées, cultures). Préciser les mois de minimum et maximum.
    \item Expliquer la tendance de fond (courbe pointillée) par l'activité humaine.
    \item En vous appuyant sur ce document, argumenter sur le rôle régulateur des végétaux verts sur la composition de l'atmosphère.
\end{enumerate}
\end{exercice}

\begin{corrige}[Exercice n°1]
\begin{enumerate}
    \item La courbe présente une \textbf{oscillation saisonnière régulière} (pics et creux annuels) superposée à une \textbf{tendance à la hausse continue} (environ +2 ppm/an).
    \item \textbf{Printemps/Été (Avril $\rightarrow$ Septembre) :} Végétation active (feuilles, cultures). Photosynthèse intense $>$ Respiration. \textbf{Prélèvement net de CO₂} $\rightarrow$ \textbf{Minimum en Septembre/Octobre} (\textbf{415 ppm}).
    \item \textbf{Automne/Hiver (Octobre $\rightarrow$ Mars) :} Chute des feuilles, dormance, respiration du sol > photosynthèse quasi nulle. \textbf{Relâchement net de CO₂} $\rightarrow$ \textbf{Maximum en Mai/Juin} (\textbf{421 ppm}).
    \item La courbe pointillée orange (+0,1 ppm/mois soit \textbf{1,2 ppm/an}, ici \textbf{2,4 ppm} sur 2 ans) représente l'accumulation du CO₂ d'origine \textbf{anthropique} (combustion fossile, déforestation) que les puits naturels (océan, forêts) n'absorbent qu'à \textbf{50\%}.
    \item \textbf{Argumentation :} Le document prouve que la végétation agit comme un \textbf{régulateur saisonnier puissant} (amplitude \textbf{6 ppm} = \textbf{12 Gt C} brassés en 6 mois). Sans cette « pompe » photosynthétique, le CO₂ monterait bien plus vite. Cependant, la capacité de régulation est \textbf{dépassée par l'ampleur des émissions humaines} (tendance haussière inexorable). Protéger et restaurer les forêts (bassin du Congo, reboisement au Nord) augmente l'amplitude du « souffle » planétaire et freine la hausse.
\end{enumerate}
\end{corrige}

\coursec{Déforestation et déséquilibres environnementaux}

Nous venons de voir que les végétaux verts, grâce à la photosynthèse, absorbent le \ce{CO2} atmosphérique et libèrent l'\ce{O2} indispensable à la respiration des êtres vivants : ils sont les grands régulateurs des cycles du carbone et de l'oxygène. Mais que se passe-t-il lorsque ces « poumons verts » disparaissent ? Autour de nous, les signes sont visibles : collines autrefois boisées devenues nues dans l'Ouest, savanisation croissante dans l'Adamaoua, tarissement précoce des points d'eau en saison sèche dans l'Extrême-Nord. Cette dernière partie de la leçon montre comment la \cle{déforestation} rompt les équilibres écologiques entretenus par la photosynthèse, et quelles solutions le Cameroun met en œuvre pour y remédier.

\courssub{Qu'est-ce que la déforestation ?}

\textbf{Situation concrète.} En 1990, un village de la région du Centre est entouré d'une forêt dense où l'on chasse et où l'on récolte le bois. Trente ans plus tard, il faut marcher plus de dix kilomètres pour trouver du bois de chauffe, les pluies sont devenues irrégulières et les ruisseaux tarissent plus tôt. Qu'est-il arrivé ? La forêt a reculé, parcelle après parcelle.

\begin{definition}[Déforestation]
La \cle{déforestation} est la destruction durable et à grande échelle des formations forestières, c'est-à-dire la conversion des terres forestières vers d'autres usages (agriculture, pâturage, zones urbaines, exploitation minière), sans reconstitution du couvert végétal. Elle se distingue de la simple coupe de bois, qui peut être suivie d'une régénération naturelle ou d'un reboisement.
\end{definition}

\textbf{Les causes de la déforestation au Cameroun.} Le Cameroun possède l'un des plus grands massifs forestiers d'Afrique (environ 20 millions d'hectares, soit plus de 40 \% du territoire), mais ce patrimoine recule sous l'effet de plusieurs pressions combinées :

\begin{itemize}
\item l'\cle{exploitation forestière} industrielle (grumes exportées : ayous, sapelli, iroko) et artisanale, souvent mal contrôlée ;
\item l'\cle{agriculture itinérante sur brûlis} : le paysan coupe et brûle une parcelle de forêt, la cultive deux ou trois ans, puis l'abandonne lorsque le sol s'épuise pour en défricher une nouvelle ;
\item la collecte du \cle{bois de chauffe} et la fabrication du charbon de bois, première source d'énergie domestique des ménages camerounais ;
\item l'\cle{extension urbaine} et les infrastructures (routes, lotissements) : la croissance rapide de villes comme Yaoundé et Douala grignote les forêts périurbaines ;
\item l'agriculture industrielle (palmier à huile, hévéa, cacao) et l'élevage extensif.
\end{itemize}

\begin{attention}[Déforestation et déboisement : des nuances]
Toutes les coupes de bois ne sont pas de la déforestation. Une coupe sélective suivie d'une régénération naturelle ou d'un reboisement n'appauvrit pas durablement le couvert forestier. La déforestation désigne la disparition \emph{durable} de la forêt. Au Probatoire, précise toujours ce que tu entends par le terme que tu emploies.
\end{attention}

\courssub{Évolution de la surface forestière : lecture de données}

\textbf{Démarche d'investigation.} Pour évaluer l'ampleur du phénomène, les scientifiques comparent des cartes et des images satellites prises à plusieurs décennies d'intervalle.

\begin{itemize}
\item \textbf{Observation} : les images satellites montrent un recul progressif des zones vertes continues.
\item \textbf{Hypothèse} : la surface forestière diminue au fil des décennies sous l'effet des activités humaines.
\item \textbf{Données} : on mesure la surface boisée à différentes dates (estimations basées sur les données FAO/FRA, arrondies ci-dessous).
\item \textbf{Interprétation} : l'histogramme confirme une baisse régulière, de l'ordre de 0,5 à 1 \% de la surface forestière par an dans les périodes les plus critiques.
\item \textbf{Conclusion} : la déforestation est un phénomène réel, mesurable et continu ; elle réduit d'autant la surface photosynthétisante du pays.
\end{itemize}

\begin{popfigure}
\begin{center}
\begin{tikzpicture}
\begin{axis}[
ybar, bar width=18pt,
width=0.85\linewidth, height=6.5cm,
xlabel={Année}, ylabel={Surface forestière (millions d'ha)},
symbolic x coords={1990,2000,2010,2020},
xtick=data,
ymin=0, ymax=26,
ymajorgrids=true, grid style={popline!40},
nodes near coords, every node near coord/.append style={font=\small},
]
\addplot[fill=popgreen!70, draw=popdark] coordinates {(1990,24.3) (2000,22.8) (2010,21.2) (2020,20.0)};
\end{axis}
\end{tikzpicture}
\end{center}
\legende{Évolution de la surface forestière du Cameroun entre 1990 et 2020 (estimations FAO/FRA arrondies). La perte cumulée atteint plusieurs millions d'hectares, soit autant de surface chlorophyllienne perdue pour la fixation du \ce{CO2}.}
\end{popfigure}

\courssub{Conséquence 1 : augmentation du \ce{CO2} atmosphérique et effet de serre}

\textbf{Le mécanisme.} Chaque hectare de forêt dense fixe, par photosynthèse, plusieurs tonnes de \ce{CO2} par an et stocke le carbone dans le bois. Lorsque la forêt est coupée puis brûlée ou laissée à se décomposer, deux effets se combinent : la pompe à \ce{CO2} s'arrête (moins de photosynthèse) \emph{et} le carbone stocké est relâché dans l'atmosphère (combustion, décomposition). Le \ce{CO2} atmosphérique augmente donc.

\begin{definition}[Effet de serre]
L'\cle{effet de serre} est un phénomène naturel de réchauffement de l'atmosphère : certains gaz (dits « gaz à effet de serre » : \ce{CO2}, vapeur d'eau, méthane \ce{CH4}) laissent passer le rayonnement solaire vers le sol mais absorbent une partie du rayonnement infrarouge renvoyé par la Terre, qu'ils renvoient à leur tour vers la surface. Sans effet de serre, la température moyenne de la Terre serait d'environ $-18$ °C au lieu de $+15$ °C. Le problème actuel est le \emph{renforcement} de cet effet par l'excès de \ce{CO2} d'origine humaine, qui provoque le \cle{réchauffement climatique}.
\end{definition}

\begin{popfigure}
\begin{center}
\begin{tikzpicture}[
box/.style={draw=popdark, fill=popblueL, rounded corners=2pt, align=center, font=\small, minimum height=1cm, text width=3.1cm},
warn/.style={draw=popdark, fill=poporangeL, rounded corners=2pt, align=center, font=\small, minimum height=1cm, text width=3.1cm},
fleche/.style={-{Stealth[length=2.5mm]}, thick, popdark}
]
\node[box, align=center] (def) at (0,0){\textbf{Déforestation}\\ (coupe, brûlis)};
\node[box, align=center] (photo) at (4.6,0){$\downarrow$ surface chlorophyllienne\\ $\rightarrow$ $\downarrow$ photosynthèse};
\node[warn, align=center] (co2) at (9.2,0){$\uparrow$ \ce{CO2} atmosphérique\\ (fixation réduite + carbone relâché)};
\node[warn, align=center] (ges) at (9.2,-2.6){Renforcement de\\ l'effet de serre};
\node[warn, align=center] (clim) at (4.6,-2.6){\textbf{Réchauffement climatique}\\ (sécheresses, pluies irrégulières)};
\draw[fleche] (def) -- (photo);
\draw[fleche] (photo) -- (co2);
\draw[fleche] (co2) -- (ges);
\draw[fleche] (ges) -- (clim);
\draw[fleche] (clim.west) -- ++(-2.2,0) |- node[pos=0.25, left, font=\scriptsize, align=right]{boucle\\ aggravante} (def.west);
\end{tikzpicture}
\end{center}
\legende{Chaîne de causalité reliant la déforestation au réchauffement climatique. La boucle de retour indique que le climat dégradé fragilise à son tour les forêts restantes (stress hydrique, feux).}
\end{popfigure}

\begin{exemplebox}[Ordre de grandeur : le \ce{CO2} non absorbé par hectare coupé]
Une forêt dense humide d'Afrique centrale séquestre en moyenne (bilan net d'écosystème) environ \textbf{5 à 10 tonnes de \ce{CO2} par hectare et par an}. Couper 100 ha de forêt, c'est donc renoncer à l'absorption d'environ 500 à 1000 tonnes de \ce{CO2} par an, auxquelles s'ajoutent plusieurs dizaines de milliers de tonnes de \ce{CO2} relâchées lors de la combustion ou de la décomposition du bois (une forêt dense stocke 150 à 250 tonnes de carbone par hectare dans sa biomasse aérienne). À l'échelle des centaines de milliers d'hectares perdus chaque année dans le bassin du Congo, l'impact sur le bilan carboné de la planète est considérable.
\end{exemplebox}

\courssub{Conséquence 2 : baisse de la production d'\ce{O2} et perturbation du cycle de l'eau}

\textbf{Oxygène.} En produisant moins de matière organique, une surface déforestée libère aussi moins d'\ce{O2}. Il faut cependant raisonner avec des ordres de grandeur corrects (voir l'encadré « Attention » plus bas) : le danger principal n'est pas un manque brutal d'oxygène, mais le dérèglement du \ce{CO2}.

\textbf{Cycle de l'eau.} C'est une conséquence souvent sous-estimée. Un grand arbre puise l'eau du sol par ses racines et la rejette dans l'atmosphère sous forme de vapeur par ses stomates : c'est l'\cle{évapotranspiration}. Une forêt dense comme celle du Sud-Cameroun recycle ainsi une grande partie de l'eau des pluies, humidifie l'air et favorise la formation de nouvelles précipitations en aval. Lorsque la forêt disparaît :

\begin{itemize}
\item l'évapotranspiration chute fortement, l'air devient plus sec ;
\item les pluies diminuent ou deviennent irrégulières localement (sécheresses locales, allongement de la saison sèche) ;
\item l'eau des pluies, non retenue par les racines et l'humus, ruisselle en surface au lieu d'alimenter les nappes phréatiques : les sources et les puits tarissent plus tôt.
\end{itemize}

\begin{exemplebox}[Un phénomène observable au Cameroun]
Dans les régions de l'Adamaoua et de l'Extrême-Nord, là où le couvert végétal a fortement reculé, les populations constatent un tarissement plus précoce des mares et des points d'eau en saison sèche, obligeant les femmes et les enfants à parcourir de plus longues distances pour s'approvisionner en eau. La restauration du couvert végétal autour des sources est d'ailleurs l'une des mesures recommandées par les services des eaux et forêts.
\end{exemplebox}

\courssub{Conséquence 3 : érosion des sols et perte de biodiversité}

\textbf{Érosion des sols.} Les racines des arbres maintiennent le sol ; la litière et le couvert végétal amortissent l'impact des gouttes de pluie. Sur un sol dénudé après brûlis, les pluies tropicales violentes emportent la couche fertile (lessivage, érosion en nappes puis en ravines). Le sol s'appauvrit, les rendements agricoles chutent, ce qui pousse le paysan à défricher une nouvelle parcelle : c'est un \cle{cercle vicieux}. Les sédiments arrachés colmatent en outre les rivières et perturbent la vie aquatique.

\textbf{Perte de biodiversité.} La forêt camerounaise abrite une biodiversité exceptionnelle : gorilles et chimpanzés du Sud-Est, éléphants de forêt, perroquets gris, plantes médicinales (prunier d'Afrique, \emph{Prunus africana}, de la région du Mont Cameroun), essences précieuses. La destruction de la forêt entraîne :

\begin{itemize}
\item la disparition des \cle{habitats} (niches écologiques) et la fragmentation des populations animales ;
\item l'extinction locale, voire totale, d'espèces avant même qu'elles soient connues de la science ;
\item l'appauvrissement génétique, qui réduit la capacité d'adaptation des écosystèmes.
\end{itemize}

\begin{attention}[Erreur fréquente : « la déforestation va supprimer tout l'oxygène »]
Non. Il faut nuancer avec les ordres de grandeur réels. D'une part, l'atmosphère contient un stock immense d'\ce{O2} (environ 21 \% de l'air) qui ne varie que très lentement. D'autre part, une grande partie de l'oxygène produit sur Terre provient du \emph{phytoplancton} des océans, pas seulement des forêts. Enfin, une forêt mature consomme par respiration une part importante de l'\ce{O2} qu'elle produit. Le danger majeur de la déforestation n'est donc pas une pénurie immédiate d'oxygène, mais bien : l'augmentation du \ce{CO2} (effet de serre), le dérèglement du cycle de l'eau, l'érosion et la perte de biodiversité. Au Probatoire, une argumentation qui exagère le rôle « producteur d'oxygène » sans nuance est pénalisée.
\end{attention}

\courssub{Les solutions : protéger et restaurer les forêts}

Face à ces déséquilibres, des solutions existent à toutes les échelles :

\begin{itemize}
\item le \cle{reboisement} : plantation d'arbres sur les terres dégradées (campagnes annuelles de reboisement au Cameroun, journées nationales de l'arbre) ;
\item l'\cle{agroforesterie} : associer arbres et cultures sur une même parcelle (cacaoyers plantés sous ombrage, haies vives, arbres fruitiers dans les champs), ce qui maintient un couvert végétal tout en produisant ;
\item l'\cle{exploitation forestière durable} : plans d'aménagement, coupes sélectives, quotas, certification (labels de gestion durable), lutte contre l'exploitation illégale ;
\item la création et la gestion des \cle{aires protégées} : le Cameroun compte de nombreux parcs nationaux et réserves, parmi lesquels le parc national de Waza (Extrême-Nord), le parc national de la Bénoué, le parc national de Korup (Sud-Ouest, l'une des plus anciennes forêts d'Afrique), le parc national de Lobéké et la réserve de faune du Dja (patrimoine mondial de l'UNESCO) ;
\item les énergies alternatives au bois de chauffe (gaz, foyers améliorés) pour réduire la pression sur les forêts.
\end{itemize}

\begin{savaistu}[Le Cameroun et le programme REDD+]
Le Cameroun s'est engagé dans le programme international \textbf{REDD+} (« Réduction des Émissions issues de la Déforestation et de la Dégradation des forêts »), piloté dans le cadre des Nations Unies. Le principe : les pays qui conservent et restaurent leurs forêts reçoivent des financements internationaux, car les forêts rendent un service à toute la planète en stockant le carbone. Le bassin du Congo, dont le Cameroun fait partie, est ainsi reconnu comme le « deuxième poumon vert » du monde après l'Amazonie. Le pays développe aussi des forêts communautaires, gérées par les populations locales elles-mêmes.
\end{savaistu}

\courssub{Construire une argumentation scientifique au Probatoire}

Le savoir-faire « argumenter sur l'importance des végétaux verts pour l'équilibre de l'environnement » est explicitement au programme. Une bonne argumentation scientifique n'est pas une opinion : elle s'appuie sur des faits, des mécanismes et des conséquences démontrées.

\begin{methode}[Argumenter en 3 étapes : fait, mécanisme, conséquence]
\begin{enumerate}
\item \textbf{Énoncer le fait} (donnée observable ou mesurable) : « La surface forestière du Cameroun diminue d'année en année » ou « les végétaux verts absorbent le \ce{CO2} et rejettent l'\ce{O2} ».
\item \textbf{Expliquer le mécanisme} (la relation de cause à effet, avec le vocabulaire scientifique exact) : « La photosynthèse consomme le \ce{CO2} atmosphérique ; réduire la surface chlorophyllienne réduit donc la fixation du carbone, tandis que la combustion du bois relâche le carbone stocké ».
\item \textbf{Déduire la conséquence} (et conclure) : « Le \ce{CO2} s'accumule, l'effet de serre se renforce, le climat se réchauffe ; il faut donc protéger et replanter les forêts ».
\end{enumerate}
Chaque affirmation doit être justifiée ; évite les exagérations (« sans forêt, plus d'oxygène demain ») qui affaiblissent ton propos.
\end{methode}

\begin{exoresolu}[Sujet type Probatoire : argumenter l'importance des végétaux verts]
\textbf{Énoncé.} Un chef de village déclare : « Couper la forêt n'a aucune importance, il restera toujours assez d'air pour respirer. » À l'aide de tes connaissances sur la photosynthèse et les cycles du carbone et de l'oxygène, construis en une quinzaine de lignes une argumentation scientifique montrant l'importance des végétaux verts pour l'équilibre de l'environnement.
\tcblower
\textbf{Solution détaillée (plan de réponse attendu).}

\emph{Fait 1.} Les végétaux chlorophylliens réalisent la photosynthèse : ils absorbent le \ce{CO2} de l'atmosphère et l'eau du sol pour fabriquer de la matière organique, en rejetant de l'\ce{O2} (bilan : \ce{6CO2 + 6H2O -> C6H12O6 + 6O2}).

\emph{Mécanisme 1 et conséquence.} En fixant le carbone atmosphérique dans leur biomasse, les forêts limitent la concentration en \ce{CO2}, principal gaz à effet de serre d'origine humaine. Détruire la forêt supprime cette fixation et relâche le carbone stocké : l'effet de serre se renforce, le climat se réchauffe.

\emph{Fait 2.} Par l'évapotranspiration, les arbres recyclent l'eau des pluies et humidifient l'atmosphère.

\emph{Mécanisme 2 et conséquence.} La déforestation réduit l'évapotranspiration : l'air s'assèche, les pluies deviennent irrégulières, les nappes phréatiques sont moins rechargées, les sources tarissent.

\emph{Fait 3.} Les racines fixent le sol et le couvert végétal abrite une biodiversité riche.

\emph{Conséquence.} Sans forêt : érosion des sols fertiles, disparition des habitats et des espèces.

\emph{Conclusion.} L'affirmation du chef de village est donc fausse : même si l'oxygène de l'air ne disparaît pas brutalement, la forêt est indispensable à la régulation du climat, du cycle de l'eau, à la protection des sols et au maintien de la biodiversité. Il faut la protéger (aires protégées, exploitation durable) et la restaurer (reboisement, agroforesterie).
\end{exoresolu}

\begin{exercice}[À toi de jouer]
À partir de l'histogramme de la surface forestière (figure ci-dessus), calcule la perte totale de surface entre 1990 et 2020, puis estime, en prenant une fixation moyenne de 7 tonnes de \ce{CO2} par hectare et par an, la quantité de \ce{CO2} qui n'est plus absorbée chaque année en 2020 par rapport à 1990. Conclus en deux phrases sur l'intérêt du reboisement.
\end{exercice}

\begin{corrige}[Exercice]
Perte de surface : $24{,}3 - 20{,}0 = 4{,}3$ millions d'hectares. \ce{CO2} non absorbé chaque année : $4{,}3 \times 10^{6} \times 7 \approx 30$ millions de tonnes de \ce{CO2} par an. Conclusion : la déforestation prive l'atmosphère d'une « pompe à \ce{CO2} » considérable ; reboiser, même partiellement, permet de restaurer cette fixation et de lutter contre le réchauffement climatique, tout en protégeant les sols, l'eau et la biodiversité.
\end{corrige}

\begin{aretenir}[L'essentiel de la section]
\begin{itemize}
\item La \cle{déforestation} (exploitation forestière, brûlis, bois de chauffe, urbanisation) réduit durablement la surface chlorophyllienne.
\item Conséquences : $\uparrow$ \ce{CO2} atmosphérique $\rightarrow$ renforcement de l'\cle{effet de serre} et réchauffement climatique ; $\downarrow$ évapotranspiration $\rightarrow$ perturbation du cycle de l'eau et sécheresses locales ; érosion des sols et perte de biodiversité.
\item Le danger principal n'est pas une pénurie d'\ce{O2} (nuance exigée) mais le dérèglement du \ce{CO2}, de l'eau, des sols et du vivant.
\item Solutions : reboisement, agroforesterie, exploitation durable, aires protégées (Waza, Korup, Dja...), programme REDD+.
\item Une argumentation scientifique se construit en trois étapes : \textbf{fait} $\rightarrow$ \textbf{mécanisme} $\rightarrow$ \textbf{conséquence}.
\end{itemize}
\end{aretenir}

\newpage

\coursec{Activité d'intégration}

\courssub{Objectifs de l'activité}

À l'issue de cette activité d'intégration, tu seras capable de :

\begin{itemize}
\item mobiliser le \cle{bilan global de la photosynthèse} pour effectuer des calculs simples de masses et de volumes de gaz ;
\item exploiter des données chiffrées pour quantifier le rôle des \cle{végétaux chlorophylliens} dans les cycles du \cle{carbone} et de l'\cle{oxygène} ;
\item argumenter, à l'écrit et à l'oral, sur l'importance des végétaux verts pour l'\cle{équilibre de l'environnement} et sur les conséquences de la \cle{déforestation} ;
\item produire une affiche collective comportant un schéma annoté du cycle du carbone.
\end{itemize}

\courssub{Situation}

\textbf{La forêt de ma région : un poumon à protéger.}

Le village de Nkolngock, situé dans la région du Centre, non loin du massif forestier qui s'étend vers la réserve du Dja, compte \num{1000} habitants. Depuis quelques années, les habitants observent avec inquiétude le recul de la forêt environnante : coupe de bois de chauffe, extension des champs de cacao et de manioc, exploitation clandestine du bois d'œuvre. Le chef du village, lors d'une réunion communautaire, s'interroge : \og Notre forêt nous donne du bois, du gibier et des fruits. Mais certains jeunes disent qu'elle nous donne aussi de l'air pur et qu'elle protège le climat. Est-ce vrai ? Combien vaut réellement notre forêt ? \fg

Les élèves de la classe de Première D du lycée voisin décident de répondre scientifiquement à cette question. Ils disposent des données suivantes, fournies par un document du service des eaux et forêts et par leurs recherches documentaires :

\begin{center}
\begin{tabularx}{\linewidth}{|l|X|}\hline
\rowcolor{popblueL} \textbf{Donnée} & \textbf{Valeur} \\ \hline
Absorption de \ce{CO2} par la forêt tropicale & environ 10 tonnes de \ce{CO2} par hectare et par an \\ \hline
Rejet de \ce{CO2} par habitant (activités domestiques : cuisine au bois, lampes, petits moteurs) & environ 0,3 tonne de \ce{CO2} par habitant et par an \\ \hline
Masses molaires atomiques & $M(\ce{C}) = 12$ g/mol ; $M(\ce{O}) = 16$ g/mol ; $M(\ce{H}) = 1$ g/mol \\ \hline
Volume molaire des gaz (conditions usuelles) & $V_m = 24$ L/mol \\ \hline
\end{tabularx}
\end{center}

\textbf{Problème à résoudre :} quelle masse de matière organique (glucose) et quel volume de dioxygène la forêt produit-elle chaque année par hectare ? Quelle superficie de forêt faut-il préserver ou replanter pour compenser les rejets de \ce{CO2} des \num{1000} habitants du village ? Comment convaincre le chef du village de créer un reboisement communal ?

\courssub{Consignes}

\begin{enumerate}
\item Rappelle l'\cle{équation-bilan} de la photosynthèse, en précisant les réactifs, les produits et les deux conditions nécessaires (lumière, chlorophylle). Indique l'organite cellulaire où se déroule ce phénomène.
\item Calcule les masses molaires de \ce{CO2}, du glucose \ce{C6H12O6} et de \ce{O2}.
\item En t'aidant des coefficients de l'équation-bilan, calcule la masse de glucose produite et la masse de dioxygène dégagée par un hectare de forêt en un an (pour 10 tonnes de \ce{CO2} absorbées).
\item Convertis la masse de \ce{O2} obtenue en volume de gaz dégagé (en litres puis en mètres cubes). Commente : à quoi ce dioxygène sert-il dans l'environnement ?
\item Calcule la masse totale de \ce{CO2} rejetée par les \num{1000} habitants du village en un an, puis la superficie de forêt nécessaire pour l'absorber entièrement.
\item Rédige un court texte argumentatif (environ 5 lignes) adressé au chef du village pour justifier la création d'un reboisement communal, en citant au moins deux rôles de la forêt dans les cycles du carbone et de l'oxygène et une conséquence de la déforestation.
\item En groupe, présente tes résultats sur une affiche comportant un schéma annoté du \cle{cycle du carbone} (atmosphère, végétaux, êtres vivants, déforestation/combustion).
\end{enumerate}

\courssub{Ressources à mobiliser}

\begin{aretenir}[Ce qu'il faut avoir en tête]
\begin{itemize}
\item \textbf{Équation-bilan de la photosynthèse} (en présence de lumière et de chlorophylle) :
\[ \ce{6CO2 + 6H2O -> C6H12O6 + 6O2} \]
\item La photosynthèse se déroule dans les \cle{chloroplastes}, grâce aux \cle{pigments photosynthétiques} (chlorophylles) qui captent l'énergie lumineuse.
\item Les végétaux verts sont des \cle{producteurs} : ils prélèvent le \ce{CO2} atmosphérique (cycle du carbone) et libèrent le \ce{O2} indispensable à la respiration des êtres vivants (cycle de l'oxygène).
\item Relations stœchiométriques : les quantités de matière réagissent dans les proportions des coefficients de l'équation ; $n = \dfrac{m}{M}$ et $V = n \times V_m$.
\item La déforestation réduit l'absorption de \ce{CO2} et la production de \ce{O2}, accentue l'effet de serre et appauvrit la biodiversité.
\end{itemize}
\end{aretenir}

\begin{attention}[Pièges fréquents]
N'oublie pas de tenir compte des \textbf{coefficients stœchiométriques} : 6 moles de \ce{CO2} donnent 1 mole de glucose et 6 moles de \ce{O2}. On ne peut pas écrire \og 1 tonne de \ce{CO2} donne 1 tonne de glucose \fg. Vérifie aussi l'homogénéité des unités (convertis les tonnes en grammes avant d'utiliser les masses molaires).
\end{attention}

\courssub{Démarche et corrigé commenté}

\begin{exoresolu}[La forêt de Nkolngock : résolution complète]
Énoncé : à partir des données de la situation (10 t de \ce{CO2} absorbées par hectare et par an ; 0,3 t de \ce{CO2} rejetées par habitant et par an ; \num{1000} habitants), calculer la masse de glucose et le volume de \ce{O2} produits par hectare, puis la superficie de forêt compensant les rejets du village, et rédiger l'argumentaire.
\tcblower
\textbf{Étape 1 : équation-bilan.} La photosynthèse se déroule dans les \cle{chloroplastes} des cellules des feuilles, en présence de \textbf{lumière} et de \textbf{chlorophylle} :
\[ \ce{6CO2 + 6H2O -> C6H12O6 + 6O2} \]
Les réactifs sont le dioxyde de carbone et l'eau ; les produits sont le glucose (matière organique) et le dioxygène.

\textbf{Étape 2 : masses molaires.}
\begin{align*}
M(\ce{CO2}) &= 12 + 2 \times 16 = 44 \text{ g/mol} \\
M(\ce{C6H12O6}) &= 6 \times 12 + 12 \times 1 + 6 \times 16 = 180 \text{ g/mol} \\
M(\ce{O2}) &= 2 \times 16 = 32 \text{ g/mol}
\end{align*}

\textbf{Étape 3 : masses produites par hectare.} D'après l'équation, 6 moles de \ce{CO2}, soit $6 \times 44 = 264$ g, donnent 1 mole de glucose (180 g) et 6 moles de \ce{O2} (192 g). On applique la proportionnalité à 10 tonnes de \ce{CO2} :
\begin{align*}
m(\text{glucose}) &= 10 \times \dfrac{180}{264} \approx 6{,}8 \text{ tonnes} \\
m(\ce{O2}) &= 10 \times \dfrac{192}{264} \approx 7{,}3 \text{ tonnes}
\end{align*}
Un hectare de forêt tropicale fabrique donc environ \textbf{6,8 tonnes de glucose} (matière organique : bois, feuilles, fruits) et dégage environ \textbf{7,3 tonnes de dioxygène} par an.

\textbf{Étape 4 : volume de dioxygène dégagé.} On convertit : $m(\ce{O2}) = 7{,}3 \text{ t} = 7{,}3 \times 10^{6}$ g.
\begin{align*}
n(\ce{O2}) &= \dfrac{m}{M} = \dfrac{7{,}3 \times 10^{6}}{32} \approx 2{,}28 \times 10^{5} \text{ mol} \\
V(\ce{O2}) &= n \times V_m \approx 2{,}28 \times 10^{5} \times 24 \approx 5{,}5 \times 10^{6} \text{ L} \approx \num{5500} \text{ m}^3
\end{align*}
Chaque hectare de forêt dégage donc environ \textbf{\num{5500} mètres cubes de dioxygène par an}, gaz indispensable à la respiration des hommes et des animaux.

\textbf{Étape 5 : superficie compensatrice.} Rejets annuels du village :
\[ \num{1000} \times 0{,}3 = 300 \text{ tonnes de } \ce{CO2} \]
Superficie nécessaire :
\[ S = \dfrac{300}{10} = 30 \text{ hectares} \]
Il faut donc préserver (ou replanter) au moins \textbf{30 hectares de forêt} pour absorber la totalité du \ce{CO2} rejeté par les activités domestiques des \num{1000} habitants.

\textbf{Étape 6 : texte argumentatif (modèle de réponse).} \og Grand Chef, notre forêt est un véritable poumon : chaque hectare absorbe 10 tonnes de dioxyde de carbone par an et nous rend en échange plus de \num{5000} mètres cubes d'oxygène pur, l'air que nous respirons. Elle fabrique aussi près de 7 tonnes de matière vivante qui nourrit les hommes et les animaux. Notre village rejette 300 tonnes de gaz carbonique par an : seuls 30 hectares de forêt peuvent les compenser. Si nous déboisons, ce gaz s'accumulera dans l'atmosphère, le climat se déréglera et nos sols s'épuiseront. Créons donc un reboisement communal pour nos enfants. \fg

\textbf{Étape 7 : affiche.} Le schéma du cycle du carbone doit montrer : le \ce{CO2} atmosphérique, la \textbf{photosynthèse} (flèche vers les végétaux verts), la \textbf{respiration} des êtres vivants (flèche retour vers l'atmosphère), et la \textbf{déforestation/combustion du bois} comme source supplémentaire de \ce{CO2}.
\end{exoresolu}

\courssub{Bilan de l'activité}

Cette activité t'a permis de passer du savoir abstrait (l'équation-bilan de la photosynthèse) à une \textbf{quantification concrète} du rôle des végétaux verts dans ton environnement proche. Tu as mis en évidence que :

\begin{itemize}
\item la photosynthèse est un phénomène \textbf{mesurable} : à 10 tonnes de \ce{CO2} absorbées correspondent environ 6,8 tonnes de matière organique et \num{5500} m$^3$ de \ce{O2} par hectare et par an ;
\item les végétaux chlorophylliens assurent l'\textbf{équilibre des cycles du carbone et de l'oxygène} : ils fixent le carbone atmosphérique et régénèrent le dioxygène respiré par tous les êtres vivants ;
\item la \textbf{déforestation} rompt cet équilibre : moins de \ce{CO2} absorbé, moins de \ce{O2} produit, accumulation des gaz à effet de serre, érosion des sols et perte de biodiversité ;
\item la science fournit des \textbf{arguments chiffrés} pour la prise de décision citoyenne : ici, convaincre une communauté villageoise de préserver ou de replanter 30 hectares de forêt.
\end{itemize}

Tu es désormais capable d'exploiter des résultats expérimentaux sur les échanges gazeux photosynthétiques, d'écrire et d'interpréter le bilan global de la photosynthèse, et d'argumenter sur l'importance des végétaux verts pour l'équilibre de l'environnement : trois savoir-faire exigibles au Probatoire.

\newpage

\coursec{Méthodes et exercices résolus}

\courssub{Méthodes et exercices résolus}

\begin{methode}[Mettre en évidence le dégagement de dioxygène par une plante chlorophyllienne en lumière]
    \textbf{Objectif :} Démontrer que la plante verte émet du $\ce{O2}$ uniquement en présence de lumière.
    \begin{enumerate}
        \item \textbf{Matériel :} Plante aquatique (\emph{Elodea} ou \emph{Cabomba}), eau de source, bécher, entonnoir, tube à essai gradué (ou eudiomètre), source lumineuse (lampe), écran opaque.
        \item \textbf{Préparation :} Plonger la plante dans l'eau sous l'entonnoir. Remplir le tube à essai d'eau, l'inverser sur le tuyau de l'entonnoir (pas de bulle d'air initiale).
        \item \textbf{Expérience témoin (obscurité) :} Envelopper le montage dans l'écran opaque. Laisser $\SI{30}{-}\SI{60}{min}$. Observer : \emph{aucune bulle} ne se forme dans le tube.
        \item \textbf{Expérience test (lumière) :} Retirer l'écran, éclairer à distance fixe. Laisser $\SI{30}{-}\SI{60}{min}$. Observer : \emph{des bulles} s'accumulent dans le tube (le niveau d'eau baisse).
        \item \textbf{Identification du gaz :} Prélever le gaz avec une seringue ou approcher une allumette incandescente au tube : \emph{rallumage vif} $\rightarrow$ présence de $\ce{O2}$.
        \item \textbf{Contrôle (facteurs limitants) :} Varier l'intensité lumineuse (distance lampe) ou la concentration en $\ce{CO2}$ (ajout de $\ce{NaHCO3}$) : compter les bulles/min ou mesurer le volume de gaz/temps $\rightarrow$ courbe de vitesse.
    \end{enumerate}
    \textbf{Réflexes Probatoire :} Citer le témoin (obscurité), l'identification formelle du gaz (allumette), et la quantification (volume ou nombre de bulles/unité de temps). Ne pas confondre bulles de $\ce{O2}$ (photosynthèse) et bulles d'air dissous (chauffage de l'eau).
\end{methode}

\begin{exoresolu}[Analyse de l'influence de l'intensité lumineuse sur le dégagement d'$\ce{O2}$]
    \textbf{Énoncé :} On réalise le montage classique à l'\emph{Elodea}. On place la lampe à différentes distances $d$ de la plante et on mesure le volume de gaz $V$ (en $\text{mL}$) recueilli en $t = \SI{20}{min}$. Résultats :
    \begin{center}
        \begin{tabular}{|c|c|c|c|c|c|}\hline
        Distance $d\ (\text{cm})$ & 10 & 20 & 30 & 40 & 50 \\ \hline
        Volume $V\ (\text{mL})$ & 4,8 & 3,2 & 1,8 & 0,9 & 0,3 \\ \hline
        \end{tabular}
    \end{center}
    \textbf{Questions :}
    \begin{enumerate}
        \item Tracer la courbe $V = f(d)$ et décrire l'allure.
        \item Calculer la vitesse de photosynthèse $v = V/t$ (en $\text{mL/min}$) pour chaque distance.
        \item Interpréter l'évolution de la vitesse en fonction de la distance. Quel est le facteur limitant ici ?
        \item Préciser le rôle du témoin « obscurité » dans cette démarche.
    \end{enumerate}
    \tcblower
    \textbf{Solution détaillée :}
    \begin{enumerate}
        \item \textbf{Courbe :} Dans un repère orthogonal, $d$ en abscisse (échelle $\SI{1}{cm} \rightarrow \SI{5}{cm}$), $V$ en ordonnée (échelle $\SI{1}{cm} \rightarrow \SI{0,5}{mL}$). Les points s'alignent sur une courbe \textbf{décroissante}, convexe, qui tend vers 0 quand $d$ augmente.
        \item \textbf{Vitesses :} $v = V / 20$.
        \begin{center}
            \begin{tabular}{|c|c|c|c|c|c|}\hline
            $d\ (\text{cm})$ & 10 & 20 & 30 & 40 & 50 \\ \hline
            $v\ (\text{mL/min})$ & 0,24 & 0,16 & 0,09 & 0,045 & 0,015 \\ \hline
            \end{tabular}
        \end{center}
        \item \textbf{Interprétation :} La vitesse $v$ diminue quand la distance $d$ augmente. Or, l'intensité lumineuse $I$ est proportionnelle à $1/d^2$ (loi de l'éclairement). Quand on éloigne la lampe, $I$ diminue fortement $\rightarrow$ l'énergie lumineuse fournie aux chloroplastes baisse $\rightarrow$ le taux de photolyse de l'eau (étape claire) diminue $\rightarrow$ moins de $\ce{O2}$ libéré. \textbf{La lumière est le facteur limitant} dans cette plage de distances (autres facteurs : $\ce{CO2}$, température, constants et non limitants).
        \item \textbf{Rôle du témoin :} Le montage dans l'obscurité (même plante, même eau, même durée, même température) ne montre \textbf{aucune accumulation de gaz}. Il valide que le gaz observé en lumière n'est pas de l'air dissous qui s'échaufferait, ni un dégagement respiratoire mitochondrial (qui serait faible et constant), mais bien un produit \textbf{spécifique de la réaction photochimique}. C'est la preuve du lien de causalité : \textbf{Lumière $\rightarrow$ Dégagement $\ce{O2}$}.
    \end{enumerate}
    \textbf{Conclusion :} La photosynthèse libère du $\ce{O2}$ ; sa vitesse dépend directement de l'intensité lumineuse reçue par la plante.
\end{exoresolu}

\begin{methode}[Mettre en évidence l'absorption du dioxyde de carbone par une plante en lumière (indicateur pH)]
    \textbf{Objectif :} Montrer que la plante consomme du $\ce{CO2}$ dissous lors de la photosynthèse.
    \begin{enumerate}
        \item \textbf{Principe :} Le $\ce{CO2}$ dissous acidifie le milieu ($\ce{CO2 + H2O <=> H2CO3 <=> H+ + HCO3-}$). Un indicateur pH (bromothymol bleu BBT, rouge phénol, cresol red) change de couleur selon le pH.
        \item \textbf{Matériel :} 4 tubes à essai, eau de source + indicateur (BBT : bleu pH $> 7,6$ ; vert neutre ; jaune pH $< 6,0$), $\ce{CO2}$ (expiration dans l'eau via paille ou eau gazeuse), brins d'\emph{Elodea}, papier aluminium, lampe.
        \item \textbf{Préparation de l'eau carbonatée :} Souffler dans l'eau + BBT jusqu'à couleur \textbf{jaune} (pH $\approx 5,5$--$6,0$, excès $\ce{CO2}$). Répartir dans les 4 tubes.
        \item \textbf{Montages :}
            \begin{itemize}
                \item T1 : Eau carbonatée seule (témoin abiotique) -- Lumière.
                \item T2 : Eau carbonatée + \emph{Elodea} -- \textbf{Obscurité} (papier alu).
                \item T3 : Eau carbonatée + \emph{Elodea} -- \textbf{Lumière}.
                \item T4 : Eau carbonatée + \emph{Elodea} bouillie (tuée) -- Lumière (contrôle biologique).
            \end{itemize}
        \item \textbf{Observations après $\SI{1}{-}\SI{2}{h}$ :}
            \begin{itemize}
                \item T1 (Lumière, pas plante) : \textbf{Jaune} inchangé (pas de $\ce{CO2}$ consommé abiotiquement).
                \item T2 (Plante, Obscurité) : \textbf{Jaune} inchangé ou plus jaune (respiration ajoute du $\ce{CO2}$).
                \item T3 (Plante, Lumière) : \textbf{Vert puis Bleu} (pH monte $> 7,6$) $\rightarrow$ $\ce{CO2}$ consommé $\rightarrow$ milieu basique.
                \item T4 (Plante morte, Lumière) : \textbf{Jaune} inchangé.
            \end{itemize}
        \item \textbf{Conclusion :} Seule la plante \textbf{vivante} en \textbf{lumière} fait monter le pH $\rightarrow$ elle absorbe le $\ce{CO2}$ pour la photosynthèse.
    \end{enumerate}
    \textbf{Attention :} Ne pas dire « la plante rejette de l'oxygène donc le pH monte ». C'est la \textbf{consommation de $\ce{CO2}$ (acide)} qui fait monter le pH. L'indicateur doit être choisi sensible autour de la neutralité (BBT idéal).
\end{methode}

\begin{exoresolu}[Interprétation d'une expérience à l'indicateur bromothymol bleu]
    \textbf{Énoncé :} On a réalisé l'expérience décrite ci-dessus. On note les couleurs finales :
    \begin{itemize}
        \item Tube A (Eau + BBT + $\ce{CO2}$ + Plante, Lumière) : \textbf{Bleu}
        \item Tube B (Eau + BBT + $\ce{CO2}$ + Plante, Obscurité) : \textbf{Jaune}
        \item Tube C (Eau + BBT + $\ce{CO2}$, Lumière, sans plante) : \textbf{Jaune}
        \item Tube D (Eau + BBT + $\ce{CO2}$ + Plante bouillie, Lumière) : \textbf{Jaune}
    \end{itemize}
    Rappel : BBT jaune si pH $< 6,0$ ; vert si $6,0 < \text{pH} < 7,6$ ; bleu si pH $> 7,6$.
    \textbf{Questions :}
    \begin{enumerate}
        \item Pour chaque tube, indiquer si le pH a augmenté, diminué ou est resté stable par rapport au départ (jaune). En déduire l'évolution de la concentration en $\ce{CO2}$ dissous.
        \item Expliquer au niveau cellulaire (chloroplaste / mitochondrie) ce qui se passe dans les tubes A et B.
        \item Pourquoi le tube C est-il un témoin indispensable ? Que valide-t-il ?
        \item Pourquoi le tube D reste-t-il jaune malgré la lumière ?
    \end{enumerate}
    \tcblower
    \textbf{Solution détaillée :}
    \begin{enumerate}
        \item \textbf{Évolution pH / $\ce{CO2}$ :}
            \begin{itemize}
                \item Tube A : Jaune $\rightarrow$ \textbf{Bleu} = pH \textbf{augmente} ($> 7,6$) $\rightarrow$ concentration en $\ce{CO2}$ \textbf{diminue fortement} (consommé).
                \item Tube B : Jaune $\rightarrow$ \textbf{Jaune} = pH \textbf{stable} (ou légère baisse) $\rightarrow$ concentration en $\ce{CO2}$ \textbf{stable ou augmente} (respiration only).
                \item Tube C : Jaune $\rightarrow$ \textbf{Jaune} = pH \textbf{stable} $\rightarrow$ $\ce{CO2}$ \textbf{stable} (pas de processus biologique).
                \item Tube D : Jaune $\rightarrow$ \textbf{Jaune} = pH \textbf{stable} $\rightarrow$ $\ce{CO2}$ \textbf{stable}.
            \end{itemize}
        \item \textbf{Niveau cellulaire :}
            \begin{itemize}
                \item \textbf{Tube A (Lumière) :} Dans les \textbf{chloroplastes}, la \textbf{phase claire} (photolyse $\ce{H2O}$) fournit $\ce{ATP}$ et $\ce{NADPH}$. La \textbf{phase sombre} (cycle de Calvin) utilise $\ce{ATP}$, $\ce{NADPH}$ et \textbf{fixe le $\ce{CO2}$} sur le RuBP $\rightarrow$ synthèse de glucides. \textbf{Bilan : $\ce{CO2}$ consommé $>$ $\ce{CO2}$ produit par respiration mitochondriale}. pH monte.
                \item \textbf{Tube B (Obscurité) :} Pas de phase claire $\rightarrow$ pas d'$\ce{ATP}$/$\ce{NADPH}$ photosynthétiques $\rightarrow$ \textbf{cycle de Calvin arrêté}. Seule la \textbf{respiration mitochondriale} fonctionne (glycolyse + cycle de Krebs + chaîne respiratoire) $\rightarrow$ \textbf{dégagement de $\ce{CO2}$}. pH stable ou baisse.
            \end{itemize}
        \item \textbf{Tube C (Témoin abiotique) :} Valide que la lumière \textbf{seule} ne dégrade pas l'indicateur, ne fait pas évaporer le $\ce{CO2}$ significativement, ne modifie pas le pH par effet thermique. Isoler la variable « \textbf{plante vivante} ».
        \item \textbf{Tube D (Plante bouillie) :} L'ébullition a \textbf{dénaturé les enzymes} (Rubisco, enzymes du cycle de Calvin, protéines des photosystèmes) et détruit l'intégrité membranaire (thylakoïdes). La plante est \textbf{morte} : plus d'activité métabolique (ni photosynthèse, ni respiration). La lumière n'a plus d'effet biologique.
    \end{enumerate}
    \textbf{Conclusion rigoureuse :} L'absorption nette de $\ce{CO2}$ est une signature de la \textbf{photosynthèse active} (plante vivante + lumière).
\end{exoresolu}

\begin{methode}[Mettre en évidence la synthèse de matière organique (amidon) dans les feuilles (Test à l'iode)]
    \textbf{Objectif :} Localiser la production d'amidon (réserve carbonée) dans le limbe foliaire et montrer sa dépendance à la lumière et au $\ce{CO2}$.
    \begin{enumerate}
        \item \textbf{Déstamisation préalable (CRUCIALE) :} Placer la plante au \textbf{noir complet $\SI{24}{-}\SI{48}{h}$}. La plante consomme ses réserves d'amidon par respiration $\rightarrow$ feuilles \textbf{non colorées} par l'iode (test négatif). Vérifier sur une feuille témoin.
        \item \textbf{Traitements expérimentaux (sur plante déstamisée) :}
            \begin{itemize}
                \item Feuille A : Exposée à la lumière normale (témoin positif).
                \item Feuille B : Maintenue au noir (témoin négatif).
                \item Feuille C : Partie masquée par papier aluminium opaque / cache-cartonnage (test lumière local).
                \item Feuille D : Plante sous cloche avec $\ce{KOH}$ (absorbe $\ce{CO2}$) + Lumière (test $\ce{CO2}$).
                \item Feuille E : Plante sous cloche avec eau + $\ce{NaHCO3}$ (source $\ce{CO2}$) + Lumière (contrôle $\ce{CO2}$ non limitant).
            \end{itemize}
        \item \textbf{Durée :} $\SI{4}{-}\SI{6}{h}$ de traitement.
        \item \textbf{Prélèvement et test (Lugol / Iode) :}
            \begin{enumerate}
                \item Bouillir la feuille $\SI{1}{-}\SI{2}{min}$ dans l'eau (tuer cellules, arrêter enzymes, ramollir).
                \item Bouillir dans l'\textbf{éthanol $\SI{90}{\degree}$} au bain-marie (déchlorophyllation $\rightarrow$ feuille blanche/jaune pâle). \textbf{Sécurité :} Pas de flamme directe, éthanol inflammable.
                \item Rincer à l'eau chaude (ramollir, enlever éthanol).
                \item Déposer sur lame / boîte de Pétri, ajouter quelques gouttes de \textbf{solution de Lugol} ($\ce{I2/KI}$).
            \end{enumerate}
        \item \textbf{Lecture :} \textbf{Noir/violet foncé} = Amidon présent (photosynthèse active). \textbf{Jaune/brun pâle} = Pas d'amidon (photosynthèse absente ou bloquée).
        \item \textbf{Interprétation attendue :} A (Noire), B (Jaune), C (Partie éclairée Noire / Partie masquée Jaune), D (Jaune), E (Noire).
    \end{enumerate}
    \textbf{Réflexes rédaction :} « Déstamisation » $\neq$ « étiolement ». Préciser \textbf{bain-marie pour l'éthanol}. L'iode révèle l'amidon, \textbf{pas le glucose} (soluble, diffusé, non coloré par Lugol).
\end{methode}

\begin{exoresolu}[Analyse d'une expérience de masquage partiel du limbe]
    \textbf{Énoncé :} Une plante déstamisée (noir $\SI{48}{h}$) porte une feuille dont la moitié du limbe est masquée par un papier aluminium opaque fixé des deux faces. La plante est exposée à la lumière $\SI{5}{h}$. On prélève la feuille, on la déchlorophylle à l'éthanol au bain-marie, on la teste au Lugol.
    \textbf{Questions :}
    \begin{enumerate}
        \item Décrire l'aspect de la feuille après révélation au Lugol. Faire un schéma annoté.
        \item Expliquer la différence de coloration entre les deux zones en détaillant les étapes biochimiques concernées.
        \item Pourquoi l'étape de déstamisation (noir $\SI{48}{h}$) est-elle indispensable avant l'expérience ?
        \item L'éthanol est chauffé au bain-marie. Justifier cette précaution et expliquer le rôle de l'éthanol.
    \end{enumerate}
    \tcblower
    \textbf{Solution détaillée :}
    \begin{enumerate}
        \item \textbf{Observation / Schéma :} La feuille présente \textbf{deux zones nettes} :
            \begin{itemize}
                \item Zone éclairée (non masquée) : \textbf{Coloration noir violet intense} (amidon abondant).
                \item Zone masquée (sous aluminium) : \textbf{Coloration jaune brun pâle} (couleur du Lugol seul, pas d'amidon).
            \end{itemize}
            \begin{popfigure}
                \begin{tikzpicture}[scale=1.2]
                    % Feuille
                    \draw[popgreen, thick, fill=popgreen!20] (0,0) .. controls (1,1.5) and (3,1.5) .. (4,0) .. controls (3,-1.5) and (1,-1.5) .. cycle;
                    \draw[popgreen, thick] (2,0) -- (2,0.1); % nervure centrale suggérée
                    % Masque Alu (gauche)
                    \draw[popmuted, thick, fill=popmuted!30, pattern=north east lines] (-0.2,1.6) rectangle (2.1,-1.6);
                    \node[popdark, align=center] at (1,0) {\textbf{Zone masquée}\\\textbf{Jaune/Brune}\\(Pas d'amidon)};
                    \node[popdark, align=center] at (3,0) {\textbf{Zone éclairée}\\\textbf{Noir Violet}\\(Amidon +++)};
                    % Légende aluminium
                    \draw[<-, popmuted] (0.5,1.2) -- (0.5,2) node[above, popdark] {Papier Alu opaque};
                \end{tikzpicture}
                \legende{Feuille après test au Lugol : zone masquée (jaune, pas d'amidon) vs zone éclairée (noir violet, amidon).}
            \end{popfigure}
        \item \textbf{Explication biochimique :}
            \begin{itemize}
                \item \textbf{Zone éclairée :} Lumière absorbée par chlorophylles (Photosystèmes II et I) $\rightarrow$ Photolyse $\ce{H2O}$ ($\ce{O2}$ libéré, $\ce{e-}$, $\ce{H+}$) $\rightarrow$ Synthèse \textbf{ATP} + \textbf{NADPH} (Phase claire). Dans le stroma, \textbf{Cycle de Calvin} : Fixation $\ce{CO2}$ sur RuBP (Rubisco) $\rightarrow$ 2 PGA $\rightarrow$ Réduction (ATP, NADPH) $\rightarrow$ \textbf{G3P (Triose-P)}. Une partie du G3P $\rightarrow$ Régénération RuBP. L'autre partie $\rightarrow$ Synthèse \textbf{Amidon} (polymère de glucose) dans le chloroplaste (stockage transitoire jour).
                \item \textbf{Zone masquée :} Pas de lumière $\rightarrow$ \textbf{Phase claire impossible} $\rightarrow$ Pas d'ATP/NADPH photosynthétiques $\rightarrow$ \textbf{Cycle de Calvin bloqué} (pas de réduction du PGA). Pas de synthèse de G3P $\rightarrow$ \textbf{Pas de synthèse d'amidon}. La respiration mitochondriale continue mais ne synthétise pas d'amidon.
            \end{itemize}
        \item \textbf{Rôle de la déstamisation :} Elle \textbf{vidange les réserves d'amidon préexistantes} (synthétisées les jours précédents). Sans elle, la zone masquée serait \textbf{noire} (amidon résiduel), ce qui fausserait la conclusion : on croirait à tort que la photosynthèse a eu lieu dans l'obscurité. Elle assure la \textbf{spécificité} du test : l'amidon détecté \textbf{nécessairement} synthétisé \textbf{pendant} l'expérience.
        \item \textbf{Éthanol au bain-marie :}
            \begin{itemize}
                \item \textbf{Rôle :} Solvant apolaire $\rightarrow$ extrait les \textbf{pigments chlorophylliens} (chlorophylles a, b, caroténoïdes) des thylakoïdes $\rightarrow$ feuille décolorée (blanche/jaune) $\rightarrow$ permet de \textbf{voir la couleur du complexe Iode-Amidon} (bleu noir) sans interférence du vert chlorophyllien.
                \item \textbf{Bain-marie :} L'éthanol est \textbf{très inflammable} (point d'éclair $\approx \SI{13}{\celsius}$). Le chauffer \textbf{directement à la flamme} provoquerait un \textbf{incendie/brûlure grave}. Le bain-marie limite la température à $\SI{100}{\celsius}$ (ébullition eau) $\rightarrow$ ébullition douce de l'éthanol ($\SI{78}{\celsius}$) en sécurité.
            \end{itemize}
    \end{enumerate}
\end{exoresolu}

\begin{methode}[Écrire et interpréter le bilan global de la photosynthèse]
    \textbf{Objectif :} Maîtriser l'équation bilan, sa stœchiométrie, son bilan énergétique et redox.
    \begin{enumerate}
        \item \textbf{Équation bilan (stœchiométrie standard) :}
            \[\ce{6 CO2 + 12 H2O ->[Lumiere][Chlorophylles] C6H12O6 + 6 O2 + 6 H2O}\]
            Souvent simplifiée (en annulant 6 $\ce{H2O}$ des deux côtés) :
            \[\ce{6 CO2 + 6 H2O -> C6H12O6 + 6 O2}\]
            \textbf{Attention :} La version complète (12 $\ce{H2O}$) montre l'origine de l'$\ce{O2}$ libéré (photolyse de l'eau) et le bilan électronique exact. Au Probatoire, \textbf{les deux sont acceptées} mais la complète est scientifiquement rigoureuse. Préciser : « en présence de lumière et de chlorophylles ».
        \item \textbf{Bilan redox :} $\ce{CO2}$ (C degré oxydation $+IV$) $\rightarrow$ $\ce{C6H12O6}$ (C degré oxydation 0) = \textbf{Réduction} (gain 24 e- par glucose). $\ce{H2O}$ (O $-II$) $\rightarrow$ $\ce{O2}$ (O 0) = \textbf{Oxydation} (perte 4 e- par $\ce{O2}$, 24 e- pour 6 $\ce{O2}$). Électrons transférés de l'eau au carbone via porteurs (NADP+ $\rightarrow$ NADPH).
        \item \textbf{Bilan énergétique :} Réaction \textbf{endothermique} ($\Delta H > 0$, $\Delta G^\circ \approx +\SI{2870}{kJ/mol}$ glucose). Énergie fournie par \textbf{photons} (longueur d'onde $\SI{400}{-}\SI{700}{nm}$, PAR). Rendement quantique maximal $\approx$ 8--10 photons / $\ce{CO2}$ fixé.
        \item \textbf{Interprétation des coefficients :}
            \begin{itemize}
                \item 6 $\ce{CO2}$ fixés $\rightarrow$ 1 molécule $\ce{C6}$ (glucose).
                \item 12 $\ce{H2O}$ oxydées $\rightarrow$ 6 $\ce{O2}$ libérés + 24 $\ce{H+}$/e- pour réduction.
                \item 6 $\ce{H2O}$ régénérées dans le cycle de Calvin (réactions de condensation).
            \end{itemize}
        \item \textbf{Calculs stœchiométriques types :} $n(\ce{CO2}) = 6 \times n(\ce{C6H12O6})$ ; $V(\ce{O2}) = 6 \times V(\ce{CO2})$ (mêmes T, P) ; Masse glucose $= n \times \SI{180}{g/mol}$.
    \end{enumerate}
    \textbf{Pièges :} Ne pas écrire $\ce{C6H12O6}$ comme « sucre » générique (c'est le glucose, produit direct du cycle). Ne pas oublier l'état physique : $\ce{CO2(g)}$, $\ce{H2O(l)}$, $\ce{C6H12O6(aq/s)}$, $\ce{O2(g)}$. $\Delta G > 0$ : réaction non spontanée, couplée à l'absorption lumineuse.
\end{methode}

\begin{exoresolu}[Calculs stœchiométriques et bilan énergétique à partir du bilan global]
    \textbf{Énoncé :} Une feuille de maïs (\emph{Zea mays}) de surface $S = \SI{250}{cm^2}$ fixe du $\ce{CO2}$ à un flux moyen de $J = \SI{35}{\micro mol\per\square\meter\per\second}$ pendant une journée de $\SI{12}{h}$ d'ensoleillement effectif.
    \textbf{Données :} $M(\ce{C}) = \SI{12,0}{g/mol}$ ; $M(\ce{H}) = \SI{1,0}{g/mol}$ ; $M(\ce{O}) = \SI{16,0}{g/mol}$. $M(\ce{C6H12O6}) = \SI{180,0}{g/mol}$. $1\ \text{mol}$ gaz $\approx \SI{24,0}{L}$ (CNTP). Énergie lumineuse reçue (PAR) : $E_{\text{recue}} = \SI{15}{MJ/m^2}$ sur la journée.
    \textbf{Questions :}
    \begin{enumerate}
        \item Calculer la quantité de matière de $\ce{CO2}$ fixée (en mol) par cette feuille sur la journée.
        \item En déduire la masse de glucose synthétisée (en g) et le volume de $\ce{O2}$ libéré (en L aux CNTP).
        \item Calculer l'énergie chimique stockée dans le glucose produit (enthalpie de combustion $\Delta H_c^\circ(\text{glucose}) = -\SI{2800}{kJ/mol}$).
        \item Calculer le rendement énergétique de la photosynthèse pour cette feuille (énergie stockée / énergie lumineuse PAR reçue). Commenter.
    \end{enumerate}
    \tcblower
    \textbf{Solution détaillée :}
    \begin{enumerate}
        \item \textbf{Quantité de $\ce{CO2}$ fixée :}
            Surface $S = \SI{250}{cm^2} = 250 \times 10^{-4}\ \text{m}^2 = \SI{0,0250}{m^2}$.
            Durée $t = \SI{12}{h} = 12 \times 3600\ \text{s} = \SI{43200}{s}$.
            Flux $J = \SI{35}{\micro mol\per\square\meter\per\second} = 35 \times 10^{-6}\ \text{mol}\ \text{m}^{-2}\ \text{s}^{-1}$.
            $n(\ce{CO2}) = J \times S \times t = (35 \times 10^{-6}) \times 0,0250 \times 43200$.
            $n(\ce{CO2}) = 35 \times 0,0250 \times 43,2 \times 10^{-3} = 37,8 \times 10^{-3}\ \text{mol} = \mathbf{3,78 \times 10^{-2}\ \text{mol}}$.
        \item \textbf{Masse glucose \& Volume $\ce{O2}$ :}
            Bilan : $\ce{6 CO2 -> C6H12O6 + 6 O2}$.
            $n(\text{glucose}) = n(\ce{CO2}) / 6 = 3,78 \times 10^{-2} / 6 = \mathbf{6,30 \times 10^{-3}\ \text{mol}}$.
            $m(\text{glucose}) = n \times M = 6,30 \times 10^{-3} \times 180,0 = \mathbf{1,134\ \text{g}}$.
            $n(\ce{O2}) = n(\ce{CO2}) = 3,78 \times 10^{-2}\ \text{mol}$ (coefficients 6:6).
            $V(\ce{O2}) = n \times V_m = 3,78 \times 10^{-2} \times 24,0 = \mathbf{0,907\ \text{L}}$ (soit $\approx \SI{907}{mL}$).
        \item \textbf{Énergie chimique stockée :}
            $\Delta H_c^\circ = -\SI{2800}{kJ/mol}$ (énergie libérée par combustion = énergie stockée dans les liaisons).
            $E_{\text{stock}} = n(\text{glucose}) \times |\Delta H_c^\circ| = 6,30 \times 10^{-3} \times 2800 = \mathbf{17,64\ \text{kJ}}$.
        \item \textbf{Rendement énergétique :}
            Énergie PAR reçue par la feuille : $E_{\text{recue}} = \SI{15}{MJ/m^2} \times \SI{0,0250}{m^2} = \SI{0,375}{MJ} = \mathbf{375\ \text{kJ}}$.
            Rendement $\eta = \dfrac{E_{\text{stock}}}{E_{\text{recue}}} \times 100 = \dfrac{17,64}{375} \times 100 \approx \mathbf{4,7\%}$.
            \textbf{Commentaire :} Ce rendement ($\approx 5\%$) est typique des cultures en conditions réelles (champ). Il est bien inférieur au rendement théorique maximal ($\approx 11$--$12\%$ pour la lumière PAR, $\approx 25$--$30\%$ pour l'énergie totale solaire) car : (1) Saturation lumineuse (excès photons dissipés chaleur), (2) Photorespiration (surtout plante C3, mais maïs est C4 $\rightarrow$ moins), (3) Respiration mitochondriale (perte $\approx 30$--$50\%$ du brut), (4) Facteurs limitants transitoires ($\ce{CO2}$, $\ce{H2O}$, T°), (5) Réflexion/transmission feuilles (seulement $\approx 80$--$90\%$ PAR absorbé).
    \end{enumerate}
    \textbf{Conclusion :} La photosynthèse stocke environ $\SI{1,13}{g}$ de glucose et libère $\SI{0,91}{L}$ d'$\ce{O2}$ pour cette feuille en une journée, avec un rendement énergétique d'environ $5\%$.
\end{exoresolu}

\begin{methode}[Argumenter sur le rôle des végétaux verts dans les cycles du carbone et de l'oxygène]
    \textbf{Objectif :} Construire un argumentaire structuré reliant photosynthèse, respiration, cycles biogéochimiques et enjeux environnementaux (déforestation).
    \begin{enumerate}
        \item \textbf{Cycle du Carbone (C) :}
            \begin{itemize}
                \item \textbf{Puits / Entrée atmosphère $\rightarrow$ Biosphère :} \textbf{Photosynthèse} (végétaux chlorophylliens terrestres + phytoplancton océanique). $\ce{CO2_{(atm)}} \xrightarrow{\text{photo}} \ce{C_{organique}}$ (biomasse). Flux brut $\approx \SI{120}{Gt C/an}$.
                \item \textbf{Source / Sortie Biosphère $\rightarrow$ Atmosphère :} \textbf{Respiration} (végétaux, animaux, décomposeurs) + \textbf{Décomposition} (humification/minéralisation) + \textbf{Incendies} + \textbf{Activités humaines} (combustion fossiles, déforestation). Flux brut $\approx \SI{120}{Gt C/an}$ (respiration/décomp.) + $\SI{10}{Gt C/an}$ (anthropique).
                \item \textbf{Rôle clé :} Les végétaux sont le \textbf{principal portail d'entrée} du carbone inorganique dans le vivant. Sans eux, pas de base trophique, $\ce{CO2}$ s'accumule $\rightarrow$ effet de serre.
                \item \textbf{Stockage long terme :} Bois (forêts), sols (humus), sédiments marins (carbonates, kérogène). \textbf{Forêts tropicales (bassin du Congo, Amazonie) = stocks majeurs.}
            \end{itemize}
        \item \textbf{Cycle de l'Oxygène (O) :}
            \begin{itemize}
                \item \textbf{Source $\ce{O2}$ :} \textbf{Photolyse de l'eau} lors de la photosynthèse (végétaux + cyanobactéries/algues). $\ce{2 H2O -> 4 H+ + 4 e- + O2}$. Flux $\approx \SI{300}{Gt O2/an}$.
                \item \textbf{Puits $\ce{O2}$ :} \textbf{Respiration aérobie} (quasi tout le vivant), \textbf{Combustion} (feux, industrie), \textbf{Oxydation} roches (altération), \textbf{Décomposition}.
                \item \textbf{Rôle clé :} Maintien du taux d'$\ce{O2}$ atmosphérique ($\approx 21\%$) indispensable à la respiration aérobie et à la couche d'ozone (UV).
            \end{itemize}
        \item \textbf{Structure de l'argumentation (Probatoire) :}
            \begin{enumerate}
                \item \textbf{Thèse :} « Les végétaux chlorophylliens sont les piliers des cycles du C et de l'O. »
                \item \textbf{Argument 1 (Carbone) :} Fixation $\ce{CO2}$ $\rightarrow$ Matière organique $\rightarrow$ Base chaînes alimentaires + Stockage (bois, sols). Chiffres clés (Gt C/an).
                \item \textbf{Argument 2 (Oxygène) :} Photolyse $\ce{H2O}$ $\rightarrow$ $\ce{O2}$ atmosphère $\rightarrow$ Respiration aérobie + Bouclier UV.
                \item \textbf{Argument 3 (Couplage) :} Photosynthèse/Respiration = couple redox inverse. Équilibre dynamique.
                \item \textbf{Nuance / Menace :} Déforestation $\rightarrow$ Perte puits C + Émission C (brûlis/décomp.) + Perte production $\ce{O2}$ + Érosion sols + Perte biodiversité.
                \item \textbf{Conclusion :} Préserver les forêts (reboisement, agroforesterie, aires protégées) = atténuation changement climatique + maintien air respirable.
            \end{enumerate}
    \end{enumerate}
    \textbf{Vocabulaire imposé :} \cle{Puits de carbone}, \cle{Source de carbone}, \cle{Flux brut}, \cle{Flux net}, \cle{Séquestration}, \cle{Photolyse}, \cle{Biomasse}, \cle{Déforestation}, \cle{Effet de serre}.
\end{methode}

\begin{exoresolu}[Synthèse argumentée : Déforestation au Cameroun et équilibres planétaires]
    \textbf{Énoncé :} Le Cameroun perd environ $\SI{100000}{ha}$ de forêt par an (données Global Forest Watch / MINEPDED). En vous appuyant sur le bilan de la photosynthèse et les cycles du carbone et de l'oxygène, rédigez un argumentaire structuré (20--25 lignes) expliquant les conséquences de cette déforestation sur : (1) le cycle du carbone et le climat, (2) le cycle de l'oxygène, (3) les écosystèmes locaux (sols, biodiversité, cycle de l'eau). Utilisez le vocabulaire scientifique précis.
    \tcblower
    \textbf{Solution détaillée (Modèle de rédaction Probatoire) :}
    \textbf{Introduction :} La photosynthèse ($\ce{6 CO2 + 6 H2O -> C6H12O6 + 6 O2}$) assure le couplage des cycles du carbone et de l'oxygène. Les forêts tropicales humides, comme celles du bassin du Congo au Sud-Cameroun, constituent des \textbf{puits de carbone} majeurs et des \textbf{sources d'oxygène} primaires. Leur destruction massive rompt ces équilibres.

    \textbf{1. Cycle du carbone et climat :} La déforestation transforme un \textbf{puits net de carbone} (fixation nette $\ce{CO2}$ par croissance) en \textbf{source nette}. Le brûlis des abattis (agriculture sur brûlis) et la décomposition des rémanents (bois mort, litière) libèrent brutalement le carbone stocké dans la biomasse aérienne (bois, feuilles) et les sols (humus) sous forme de $\ce{CO2}$ (et $\ce{CH4}$). Cela augmente le flux net $\ce{CO2}$ vers l'atmosphère, renforçant l'\textbf{effet de serre additionnel} et le réchauffement global. La perte de la capacité de séquestration future (arbres non remplacés) aggrave le bilan carbone national et planétaire.

    \textbf{2. Cycle de l'oxygène :} La photolyse de l'eau au niveau des photosystèmes II des chloroplastes est l'unique source significative d'$\ce{O2}$ atmosphérique. La disparition de la canopée forestière réduit d'autant la production primaire brute (PPB) et donc le flux d'$\ce{O2}$ injecté dans l'atmosphère. Bien que le stock atmosphérique d'$\ce{O2}$ soit immense ($\approx 1,2 \times 10^6\ \text{Gt}$) et que la baisse soit lente, la contribution locale à la qualité de l'air et la perte du potentiel de production sont réelles.

    \textbf{3. Écosystèmes locaux (Sols, Eau, Biodiversité) :} La mise à nu du sol expose celui-ci aux pluies tropicales intenses $\rightarrow$ \textbf{lessivage} (perte nutriments, acidification) et \textbf{ruissellement/érosion} hydrique massive (ravines, sédimentation rivières comme la Sanaga, le Wouri). La matière organique du sol (humus) s'oxyde rapidement (minéralisation accélérée par T° et humidité) $\rightarrow$ perte fertilité, libération $\ce{CO2}$ supplémentaire. La fragmentation des habitats entraîne une \textbf{érosion de la biodiversité} (espèces endémiques, grands mammifères, pollinisateurs). Le cycle de l'eau est perturbé : baisse de l'évapotranspiration $\rightarrow$ diminution précipitations locales, assèchement cours d'eau en saison sèche, perturbation microclimatique.

    \textbf{Conclusion :} La déforestation au Cameroun n'est pas qu'une perte de bois ; c'est une \textbf{rupture des services écosystémiques} (régulation climat, air, eau, sols, biodiversité). La lutte passe par la \textbf{gestion durable des forêts} (aménagement, certification FSC), l'\textbf{agroforesterie} (cacao/café sous couvert), la \textbf{restauration} (Initiative AFR100, Bonn Challenge) et la \textbf{valorisation du carbone forestier} (mécanismes REDD+).
\end{exoresolu}

\begin{methode}[Analyser un document complexe : Flux de carbone, bilan net, impact anthropique]
    \textbf{Objectif :} Savoir extraire des données d'un schéma de cycle du carbone (réservoirs en Gt C, flux en Gt C/an), calculer un bilan net, identifier les perturbations anthropiques.
    \begin{enumerate}
        \item \textbf{Lire le schéma :} Identifier les \textbf{réservoirs} (Atmosphère, Végétation, Sols, Océan, Fossiles) et leurs tailles (Gt C). Identifier les \textbf{flux} (flèches) avec leurs valeurs (Gt C/an) et sens (vers/Depuis atmosphère).
        \item \textbf{Calculer le bilan atmosphérique net :} $\Delta C_{\text{atm}} = \sum \text{Flux Entrants} - \sum \text{Flux Sortants}$.
            \begin{itemize}
                \item Entrants : Respiration/Décomp. ($\approx 120$), Océan $\rightarrow$ Atm ($\approx 90$), \textbf{Anthropique (Fossiles + Déforestation)} ($\approx 10$).
                \item Sortants : \textbf{Photosynthèse} ($\approx 120$), Atm $\rightarrow$ Océan ($\approx 92$).
            \end{itemize}
            Bilan naturel $\approx 0$ (équilibre préindustriel). Bilan actuel $= +10\ \text{Gt C/an}$ (sources anthropiques) $- \text{Puits supplémentaires}$ (fertilisation $\ce{CO2}$, réchauffement...) $\approx +5\ \text{Gt C/an}$ d'accumulation atmosphérique.
        \item \textbf{Distinction Flux Brut / Flux Net :} \textbf{Flux Brut} = échange total (ex: Photosynthèse $120$). \textbf{Flux Net} = Brut - Respiration autotrophe (RA) = \textbf{Production Primaire Nette (PPN)}. Pour la végétation : PPN $\approx 60\ \text{Gt C/an}$. C'est ce qui construit la biomasse.
        \item \textbf{Identifier l'impact « Déforestation » :} Flux « Changement d'affectation des sols » ($\approx 1,5$--$2\ \text{Gt C/an}$ source). Perte simultanée du puits « Végétation » (photosynthèse future) et émission du stock (bois/sols).
        \item \textbf{Rédiger la conclusion :} « Le système n'est plus à l'équilibre. L'apport anthropique ($+10$) dépasse la capacité d'absorption nette des puits naturels (Océan + Végétation $\approx +5$), d'où l'augmentation continue du $\ce{CO2}$ atmosphérique ($> 420\ \text{ppm}$). »
    \end{enumerate}
    \textbf{Compétence transversale :} Manipuler les puissances de 10 ($1\ \text{Gt} = 10^9\ \text{t} = 10^{15}\ \text{g}$). Convertir $\text{ppm} \leftrightarrow \text{Gt C}$ ($1\ \text{ppm} \ce{CO2} \approx 2,12\ \text{Gt C}$).
\end{methode}

\begin{exoresolu}[Exploitation d'un bilan de carbone global (Données simplifiées GIEC)]
    \textbf{Énoncé :} On donne les flux annuels moyens de carbone (en $\text{Gt C/an}$) pour la période 2010--2019 :
    \begin{itemize}
        \item Photosynthèse terrestre (PFB) : $123$ (sortie atmosphère)
        \item Respiration + Décomposition terrestres : $119$ (entrée atmosphère)
        \item Échange Océan-Atmosphère (net) : Océan absorbe $2,5$ (sortie atmosphère)
        \item Émissions combustibles fossiles + industrie : $9,4$ (entrée atmosphère)
        \item Émissions changement d'affectation des sols (déforestation nette) : $1,6$ (entrée atmosphère)
    \end{itemize}
    \textbf{Questions :}
    \begin{enumerate}
        \item Calculer le \textbf{flux net d'échange terrestre} (Puité ou Source ?) et la \textbf{Production Primaire Nette (PPN)} terrestre.
        \item Calculer le \textbf{bilan net annuel de l'atmosphère} (Accumulation en $\text{Gt C/an}$).
        \item Convertir cette accumulation en $\text{ppm/an}$ d'augmentation du $\ce{CO2}$ atmosphérique ($1\ \text{ppm} \approx 2,12\ \text{Gt C}$).
        \item Expliquer pourquoi la « Photosynthèse terrestre » ($123$) ne compense pas les émissions anthropiques ($9,4 + 1,6$), alors que $123 > 11$.
        \item Proposer deux mesures concrètes au Cameroun pour transformer le secteur « Changement d'affectation des sols » de source en puits.
    \end{enumerate}
    \tcblower
    \textbf{Solution détaillée :}
    \begin{enumerate}
        \item \textbf{Flux net terrestre / PPN :}
            Flux net = Photosynthèse (Sortie) - Respiration/Décomp. (Entrée) = $123 - 119 = \mathbf{+4\ \text{Gt C/an}}$ (Sortie nette = \textbf{Puits net} de $4\ \text{Gt C/an}$).
            \textbf{PPN} = Flux net terrestre = $\mathbf{4\ \text{Gt C/an}}$ (C'est le carbone séquestré dans la biomasse et les sols par an).
        \item \textbf{Bilan atmosphérique net :}
            Entrées totales = Respiration/Décomp ($119$) + Fossiles ($9,4$) + Déforestation ($1,6$) = $130,0\ \text{Gt C/an}$.
            Sorties totales = Photosynthèse ($123$) + Absorption Océan ($2,5$) = $125,5\ \text{Gt C/an}$.
            Accumulation nette = Entrées - Sorties = $130,0 - 125,5 = \mathbf{+4,5\ \text{Gt C/an}}$.
            (Vérification : Sources anthropiques $11,0$ - Puits naturels terrestres $4,0$ - Puits océan $2,5$ = $+4,5$).
        \item \textbf{Augmentation en ppm :}
            $\Delta \text{ppm} = \dfrac{4,5\ \text{Gt C}}{2,12\ \text{Gt C/ppm}} \approx \mathbf{2,1\ \text{ppm/an}}$.
            (Cohérent avec l'observé $\approx 2,4\ \text{ppm/an}$ récemment).
        \item \textbf{Explication :} La photosynthèse brute ($123$) est quasi entièrement compensée par la respiration autotrophe (plantes) et hétérotrophe (décomposeurs) + feux ($119$). Le \textbf{puits net} (PPN = $4$) est bien plus petit que le flux brut. Les émissions anthropiques ($11$) s'ajoutent \textbf{par-dessus} ce cycle naturel quasi équilibré. La nature n'absorbe que $\approx 40\%$ de nos émissions ($4,5$ accumulés sur $11$ émis). Le reste ($6,5$) est absorbé par les puits (terre $4$ + océan $2,5$) grâce à l'effet de fertilisation par le $\ce{CO2}$ et le réchauffement (allongement saison de croissance).
        \item \textbf{Mesures Cameroun :}
            \begin{enumerate}
                \item \textbf{Zéro déforestation nette / Restauration (AFR100) :} Arrêter la conversion forêt $\rightarrow$ agriculture extensive (cacao, huile de palme, vivriers sur brûlis). Promouvoir l'\textbf{agroforesterie} (arbres d'ombrage + cultures) et le \textbf{reboisement} d'espèces locales à croissance rapide sur terres dégradées (Nord, Adamaoua). Cela stoppe l'émission ($1,6$ source) et crée un puits (bois + sols).
                \item \textbf{Gestion durable des forêts / Certification :} Appliquer strictement les plans d'aménagement (UFA) avec rotation longue, exploitation à impact réduit (EIR), protection des zones à haute valeur de conservation (ZHVC). Valoriser le \textbf{carbone forestier} via marchés volontaires ou REDD+ pour financer la conservation. Lutter contre l'exploitation illégale (bois d'œuvre, bois de chauffe urbain).
            \end{enumerate}
    \end{enumerate}
\end{exoresolu}

\begin{attention}[Les 5 erreurs qui coûtent des points au Probatoire]
    \begin{enumerate}
        \item \textbf{Confondre « Déstamisation » et « Étiolement ».} \textit{Erreur :} « On met la plante au noir pour l'étioler. » \textbf{Correction :} La déstamisation ($\SI{24}{-}\SI{48}{h}$ noir) vide les réserves d'amidon \textbf{sans détruire les chloroplastes}. L'étiolement (semis au noir prolongé) détruit les chloroplastes (pas de protochlorophyllide $\rightarrow$ pas de chlorophylle) $\rightarrow$ plante incapable de photosynthèse ensuite. \textbf{Sanction :} 0 pt à la question « rôle de l'étape préliminaire ».
        \item \textbf{Omettre le témoin « Obscurité » ou « Plante bouillie » dans l'interprétation.} \textit{Erreur :} « La feuille devient noire, donc la lumière fait de l'amidon. » \textbf{Correction :} Il faut \textbf{comparer} : Feuille éclairée (noire) \textbf{vs} Feuille masquée/bouillie/obscurité (jaune). Seule la comparaison valide la causalité « Lumière $\rightarrow$ Amidon ». \textbf{Sanction :} Perte de 2--3 pts sur 5 (démarche scientifique incomplète).
        \item \textbf{Écrire le bilan global sans états physiques ni conditions.} \textit{Erreur :} $\ce{CO2 + H2O -> C6H12O6 + O2}$. \textbf{Correction :} $\ce{6 CO2(g) + 6 H2O(l) ->[Lumiere][Chlorophylles] C6H12O6(aq) + 6 O2(g)}$. Préciser « en présence de lumière et de chlorophylles » (ou « au niveau des chloroplastes »). \textbf{Sanction :} -1 à -2 pts (rigueur chimique).
        \item \textbf{Dire que l'$\ce{O2}$ libéré vient du $\ce{CO2}$.} \textit{Erreur :} « Le $\ce{CO2}$ perd son oxygène qui devient $\ce{O2}$. » \textbf{Correction :} L'$\ce{O2}$ libéré provient \textbf{exclusivement de la photolyse de l'eau} ($\ce{2 H2O -> 4 H+ + 4 e- + O2}$), prouvé par l'expérience au $\ce{H2^{18}O}$ (Ruben \& Kamen, 1941). Le $\ce{CO2}$ fournit le carbone et l'oxygène des glucides. \textbf{Sanction :} Erreur conceptuelle majeure, 0 pt à la question « origine de l'oxygène ».
        \item \textbf{Confondre Flux Brut et Flux Net (PPN) dans le cycle du carbone.} \textit{Erreur :} « La photosynthèse absorbe 120 Gt C, donc elle compense nos 10 Gt d'émissions. » \textbf{Correction :} La photosynthèse \textbf{brute} (PFB) est de $123\ \text{Gt C/an}$, mais la respiration/décomposition renvoie $119\ \text{Gt C/an}$. Le \textbf{puits net} (PPN) n'est que de $\approx 4\ \text{Gt C/an}$. C'est ce flux net qui peut absorber une partie des émissions anthropiques. \textbf{Sanction :} Incompréhension du cycle, perte de points sur le calcul de bilan et l'argumentation déforestation.
    \end{enumerate}
\end{attention}

\newpage

\coursec{Exercices}

\courssub{Je vérifie mes connaissances}

\begin{exercice}[QCM : les bases de la photosynthèse]
Pour chaque question, une seule réponse est correcte. Recopie la lettre correspondante.
\begin{enumerate}
\item Le siège de la photosynthèse dans la cellule végétale est :
\begin{enumerate}
\item[A.] la mitochondrie \quad B. le chloroplaste \quad C. le noyau \quad D. le ribosome
\end{enumerate}
\item Le dioxygène dégagé lors de la photosynthèse provient :
\begin{enumerate}
\item[A.] du dioxyde de carbone \quad B. du glucose \quad C. de l'eau \quad D. de la chlorophylle
\end{enumerate}
\item La photosynthèse se déroule :
\begin{enumerate}
\item[A.] jour et nuit \quad B. uniquement la nuit \quad C. uniquement à la lumière \quad D. uniquement en l'absence de \ce{CO2}
\end{enumerate}
\item Le pigment majoritaire des feuilles vertes est :
\begin{enumerate}
\item[A.] le carotène \quad B. la chlorophylle $a$ \quad C. la xanthophylle \quad D. l'hémoglobine
\end{enumerate}
\end{enumerate}
\end{exercice}

\begin{exercice}[Vrai ou faux, à justifier]
Réponds par vrai ou faux, puis justifie chaque réponse en une ou deux phrases.
\begin{enumerate}
\item Les végétaux verts ne respirent pas : ils ne font que de la photosynthèse.
\item La photosynthèse consomme du \ce{CO2} et rejette du \ce{O2}.
\item Sans lumière, une plante verte placée dans une enceinte fermée finit par enrichir l'air en \ce{CO2}.
\item La déforestation n'a aucune influence sur la teneur en \ce{CO2} de l'atmosphère.
\end{enumerate}
\end{exercice}

\begin{exercice}[Définitions essentielles]
Définis brièvement et précisément chacun des termes suivants :
\begin{enumerate}
\item Photosynthèse.
\item Chloroplaste.
\item Pigment photosynthétique.
\item Puits de carbone.
\end{enumerate}
\end{exercice}

\begin{exercice}[Calcul direct : équation-bilan]
On donne l'équation-bilan non équilibrée de la photosynthèse :
\[ \ce{CO2 + H2O ->[$\text{lumière}$][\text{chlorophylle}] C6H12O6 + O2} \]
\begin{enumerate}
\item Équilibre cette équation.
\item Calcule la masse molaire du glucose \ce{C6H12O6}.
\item Combien de moles de \ce{CO2} sont nécessaires pour produire 1 mole de glucose ?
\end{enumerate}
Données : $M(\ce{C}) = 12$ g/mol ; $M(\ce{H}) = 1$ g/mol ; $M(\ce{O}) = 16$ g/mol.
\end{exercice}

\courssub{Je m'entraîne}

\begin{exercice}[Expérience de l'élodée]
On place un rameau d'élodée (plante aquatique) dans un tube rempli d'eau additionnée d'hydrogénocarbonate de potassium (source de \ce{CO2}), éclairé par une lampe. On observe un dégagement de bulles de gaz qui se rassemblent au sommet du tube.
\begin{enumerate}
\item Quel est le gaz recueilli ? Propose un test permettant de l'identifier.
\item Que se passe-t-il si l'on place le dispositif dans l'obscurité ? Interprète.
\item Que se passe-t-il si l'on supprime la source de \ce{CO2} ? Interprète.
\item Conclus sur les conditions nécessaires à la photosynthèse.
\end{enumerate}
\end{exercice}

\begin{exercice}[Chromatographie des pigments foliaires]
On réalise la chromatographie sur papier d'un extrait de feuilles vertes d'épinard. Après migration, on observe quatre taches, du bas vers le haut : une tache vert-jaunâtre (chlorophylle $b$), une tache vert-bleu (chlorophylle $a$), une tache jaune (xanthophylle) et une tache orange (carotène).
\begin{enumerate}
\item Quel est le principe de la chromatographie sur papier ?
\item Pourquoi les pigments se séparent-ils au cours de la migration ?
\item Les chlorophylles absorbent surtout la lumière rouge et la lumière bleue-violette. Explique alors pourquoi les feuilles nous apparaissent vertes.
\item En automne, les chlorophylles sont détruites avant les autres pigments. Quelle couleur prennent alors les feuilles ? Justifie.
\end{enumerate}
\end{exercice}

\begin{exercice}[Masse de \ce{CO2} fixée]
Un jeune plant de maïs a synthétisé 90 g de glucose au cours d'une période de croissance.
\begin{enumerate}
\item Rappelle l'équation-bilan équilibrée de la photosynthèse.
\item Calcule le nombre de moles de glucose produites.
\item Déduis-en le nombre de moles, puis la masse de \ce{CO2} fixée par le plant.
\item Calcule le volume de \ce{CO2} correspondant dans les conditions où le volume molaire vaut $V_m = 24$ L/mol.
\end{enumerate}
Données : $M(\ce{C}) = 12$ g/mol ; $M(\ce{H}) = 1$ g/mol ; $M(\ce{O}) = 16$ g/mol.
\end{exercice}

\begin{exercice}[Photosynthèse et respiration : ne pas confondre]
Recopie et complète le tableau comparatif suivant :
\begin{center}
\begin{tabularx}{\linewidth}{|l|X|X|}
\hline
\rowcolor{popblueL} \textbf{Critère} & \textbf{Photosynthèse} & \textbf{Respiration cellulaire} \\
\hline
Gaz absorbé & & \\
\hline
Gaz dégagé & & \\
\hline
Condition de lumière & & \\
\hline
Organite concerné & & \\
\hline
Bilan énergétique (stockage ou libération) & & \\
\hline
\end{tabularx}
\end{center}
\end{exercice}

\courssub{Je me prépare au Probatoire}

\begin{exercice}[Type Probatoire : étude expérimentale sur l'élodée]
Au laboratoire du lycée de Bafoussam, un groupe d'élèves de Première D étudie l'influence de l'éclairement sur la photosynthèse d'une élodée. Le dispositif comprend un rameau d'élodée dans de l'eau enrichie en \ce{CO2}, placé sous une lampe dont on fait varier la distance $d$ au rameau. On mesure le volume de gaz dégagé pendant 10 minutes pour chaque distance. Résultats :

\begin{center}
\begin{tabularx}{\linewidth}{|l|X|X|X|X|X|}
\hline
\rowcolor{popblueL} Distance $d$ (cm) & 10 & 20 & 40 & 60 & 80 \\
\hline
Volume de gaz (mL) & 12,0 & 9,5 & 4,0 & 1,5 & 0,5 \\
\hline
\end{tabularx}
\end{center}

\begin{enumerate}
\item Nomme le gaz recueilli et donne le test qui permet de l'identifier formellement.
\item Trace sur papier millimétré la courbe du volume de gaz dégagé en fonction de la distance $d$. Échelle : 1 cm pour 10 cm de distance ; 1 cm pour 2 mL de gaz.
\item Décris l'allure de la courbe obtenue.
\item Sachant que l'éclairement diminue lorsque la distance à la lampe augmente, que peux-tu conclure sur le rôle de la lumière dans la photosynthèse ?
\item Le volume de gaz dégagé n'est jamais rigoureusement nul, même à 80 cm. En tenant compte de la respiration de la plante, explique pourquoi le volume mesuré à 80 cm est faible.
\item Rédige une conclusion générale sur les conditions de la photosynthèse mises en évidence par cette expérience.
\end{enumerate}
\end{exercice}

\begin{exercice}[Type Probatoire : \ce{CO2} atmosphérique et déforestation]
Le tableau ci-dessous donne la concentration moyenne en \ce{CO2} de l'atmosphère et la superficie forestière du Cameroun à différentes dates (données simplifiées).

\begin{center}
\begin{tabularx}{\linewidth}{|l|X|X|X|X|}
\hline
\rowcolor{popblueL} Année & 1970 & 1990 & 2010 & 2020 \\
\hline
Concentration en \ce{CO2} (ppm) & 325 & 354 & 390 & 415 \\
\hline
Superficie forestière (millions d'ha) & 24,0 & 22,5 & 20,0 & 19,0 \\
\hline
\end{tabularx}
\end{center}

\begin{enumerate}
\item Compare l'évolution des deux grandeurs entre 1970 et 2020. Quelle corrélation peux-tu dégager ?
\item En t'appuyant sur le rôle des végétaux chlorophylliens dans le cycle du carbone, explique le lien entre la diminution de la superficie forestière et l'augmentation de la concentration en \ce{CO2}.
\item Dégage deux conséquences de l'augmentation du \ce{CO2} atmosphérique sur l'environnement et la santé humaine.
\item Propose deux solutions concrètes, applicables au Cameroun, pour limiter ce déséquilibre, en précisant le mécanisme biologique sur lequel chacune s'appuie.
\end{enumerate}
\end{exercice}

\begin{exercice}[Type Probatoire : synthèse et argumentation]
\textbf{Partie A.} On réalise trois expériences avec des plantes vertes :
\begin{itemize}
\item Expérience 1 : une feuille de géranium partiellement couverte d'une feuille d'aluminium est exposée à la lumière, puis testée à l'eau iodée après décoloration. Seule la partie éclairée se colore en bleu-noir.
\item Expérience 2 : une plante placée sous cloche avec de la chaux sodée (qui absorbe le \ce{CO2}) ne produit pas d'amidon.
\item Expérience 3 : une plante aquatique éclairée dégage un gaz qui ravive une braise.
\end{itemize}
\begin{enumerate}
\item Interprète chacune des trois expériences et dégage les facteurs et les produits de la photosynthèse qu'elles mettent en évidence.
\item Écris l'équation-bilan équilibrée de la photosynthèse et indique l'origine de chaque réactif et le devenir de chaque produit.
\end{enumerate}
\textbf{Partie B.} À l'occasion d'une campagne de reboisement dans la région de l'Adamaoua, un élève déclare : « Les végétaux verts sont indispensables à la vie sur Terre ».
\begin{enumerate}
\item[3.] Rédige un paragraphe argumenté d'une quinzaine de lignes défendant cette affirmation, en mobilisant au moins trois savoirs de la leçon (production de matière organique, dégagement de dioxygène, régulation du \ce{CO2} atmosphérique, base des chaînes alimentaires).
\end{enumerate}
\end{exercice}

\newpage

\coursec{Corrigés des exercices}

\begin{corrige}[Exercice 1 — QCM : les bases de la photosynthèse]
\begin{enumerate}
\item \textbf{Réponse B : le chloroplaste.} La photosynthèse se déroule dans les chloroplastes, organites des cellules végétales contenant les pigments photosynthétiques. La mitochondrie est le siège de la respiration cellulaire ; le noyau contient l'ADN ; le ribosome assure la synthèse des protéines.
\item \textbf{Réponse C : de l'eau.} Le dioxygène dégagé provient de la photolyse de l'eau (clivage des molécules d'\ce{H2O} sous l'effet de la lumière), et non du \ce{CO2} : le carbone du \ce{CO2} se retrouve dans le glucose.
\item \textbf{Réponse C : uniquement à la lumière.} La photosynthèse exige l'énergie lumineuse ; elle s'arrête dans l'obscurité (seule la respiration se poursuit alors, jour et nuit).
\item \textbf{Réponse B : la chlorophylle $a$.} C'est le pigment majoritaire et le seul indispensable à la conversion de l'énergie lumineuse ; les autres pigments (chlorophylle $b$, carotène, xanthophylle) sont des pigments accessoires.
\end{enumerate}
\end{corrige}

\begin{corrige}[Exercice 2 — Vrai ou faux, à justifier]
\begin{enumerate}
\item \textbf{Faux.} Toutes les cellules vivantes de la plante respirent en permanence, jour et nuit, pour libérer l'énergie du glucose : $\ce{C6H12O6 + 6O2 -> 6CO2 + 6H2O}$. La photosynthèse ne s'ajoute à la respiration que dans les cellules chlorophylliennes éclairées.
\item \textbf{Vrai.} Le bilan global de la photosynthèse est $\ce{6CO2 + 6H2O -> C6H12O6 + 6O2}$ : le \ce{CO2} est le réactif carboné consommé et le \ce{O2} est le produit gazeux dégagé.
\item \textbf{Vrai.} Dans l'obscurité, la photosynthèse s'arrête mais la respiration se poursuit : la plante consomme du \ce{O2} et rejette du \ce{CO2}, ce qui enrichit progressivement l'air de l'enceinte en \ce{CO2}.
\item \textbf{Faux.} Les forêts sont des puits de carbone : en les détruisant, on supprime des fixateurs de \ce{CO2} et, si le bois est brûlé ou se décompose, on libère le carbone stocké. La déforestation contribue donc à l'augmentation du \ce{CO2} atmosphérique.
\end{enumerate}
\end{corrige}

\begin{corrige}[Exercice 3 — Définitions essentielles]
\begin{enumerate}
\item \textbf{Photosynthèse} : ensemble des réactions par lesquelles les végétaux chlorophylliens, à la lumière, utilisent le \ce{CO2} et l'eau pour synthétiser de la matière organique (glucose) en dégageant du dioxygène.
\item \textbf{Chloroplaste} : organite des cellules végétales vertes, délimité par une double membrane, contenant les pigments photosynthétiques logés dans les thylakoïdes ; c'est le siège de la photosynthèse.
\item \textbf{Pigment photosynthétique} : molécule colorée (chlorophylles $a$ et $b$, carotène, xanthophylle) capable d'absorber l'énergie lumineuse et de la convertir en énergie chimique.
\item \textbf{Puits de carbone} : réservoir naturel (forêt, océan, sol) qui absorbe et stocke durablement du carbone atmosphérique, retirant ainsi du \ce{CO2} de l'atmosphère.
\end{enumerate}
\end{corrige}

\begin{corrige}[Exercice 4 — Calcul direct : équation-bilan]
\begin{enumerate}
\item On équilibre d'abord le carbone : 6 \ce{CO2} pour 1 \ce{C6H12O6}. Puis l'hydrogène : il faut 6 \ce{H2O} (12 H). Enfin l'oxygène : à gauche $6\times 2 + 6 = 18$ O ; le glucose en contient 6, il reste 12 O soit 6 \ce{O2}. D'où :
\[ \ce{6CO2 + 6H2O ->[$\text{lumière}$][\text{chlorophylle}] C6H12O6 + 6O2} \]
Vérification : C : 6 = 6 ; H : 12 = 12 ; O : 18 = 6 + 12 = 18. L'équation est équilibrée.
\item Masse molaire du glucose :
\[ M(\ce{C6H12O6}) = 6\times 12 + 12\times 1 + 6\times 16 = 72 + 12 + 96 = 180 \text{ g/mol}. \]
\item D'après les coefficients stœchiométriques, il faut \textbf{6 moles de \ce{CO2}} pour produire 1 mole de glucose.
\end{enumerate}
\end{corrige}

\begin{corrige}[Exercice 5 — Expérience de l'élodée]
\begin{enumerate}
\item Le gaz recueilli est du \textbf{dioxygène \ce{O2}}. Test d'identification : on approche une braise (une allumette presque éteinte) de l'orifice du tube ; si le gaz est du \ce{O2}, la braise \textbf{se ravive} (se rallume).
\item Dans l'obscurité, le dégagement de bulles \textbf{s'arrête}. Interprétation : la lumière est indispensable à la photosynthèse ; sans lumière, l'élodée ne produit plus de \ce{O2} (elle ne fait plus que respirer).
\item Si l'on supprime la source de \ce{CO2}, le dégagement de bulles \textbf{diminue puis cesse}, même à la lumière. Interprétation : le \ce{CO2} est un réactif indispensable ; sans lui, la photosynthèse ne peut avoir lieu.
\item \textbf{Conclusion} : la photosynthèse nécessite simultanément la \textbf{lumière} et le \textbf{dioxyde de carbone} (en présence d'eau et de chlorophylle) ; elle produit du \textbf{dioxygène}.
\end{enumerate}
\end{corrige}

\begin{corrige}[Exercice 6 — Chromatographie des pigments foliaires]
\begin{enumerate}
\item \textbf{Principe} : la chromatographie sur papier est une technique de séparation des constituants d'un mélange fondée sur leur migration différentielle : un solvant (éluant) monte par capillarité le long du papier et entraîne les pigments à des vitesses différentes.
\item Les pigments se séparent parce qu'ils n'ont pas la même \textbf{solubilité} dans l'éluant ni la même \textbf{affinité} pour le papier : les plus solubles et les moins retenus (carotène) migrent le plus haut ; les moins solubles (chlorophylle $b$) restent près de la ligne de dépôt.
\item Les chlorophylles absorbent le rouge et le bleu-violet mais \textbf{n'absorbent pas le vert}, qu'elles réfléchissent (et transmettent) : c'est cette lumière verte réfléchie qui atteint nos yeux, d'où la couleur verte des feuilles.
\item En automne, les chlorophylles étant détruites, les pigments restants (carotène orange, xanthophylle jaune) deviennent visibles : les feuilles prennent des teintes \textbf{jaunes à orangées}.
\end{enumerate}
\end{corrige}

\begin{corrige}[Exercice 7 — Masse de \ce{CO2} fixée]
\begin{enumerate}
\item Équation-bilan équilibrée :
\[ \ce{6CO2 + 6H2O ->[$\text{lumière}$][\text{chlorophylle}] C6H12O6 + 6O2} \]
\item Nombre de moles de glucose produites :
\[ n(\ce{C6H12O6}) = \frac{m}{M} = \frac{90}{180} = 0{,}5 \text{ mol}. \]
\item D'après l'équation, $n(\ce{CO2}) = 6 \times n(\ce{C6H12O6}) = 6 \times 0{,}5 = 3$ mol.
Masse molaire du \ce{CO2} : $M(\ce{CO2}) = 12 + 2\times 16 = 44$ g/mol.
Masse de \ce{CO2} fixée :
\[ m(\ce{CO2}) = n \times M = 3 \times 44 = 132 \text{ g}. \]
\item Volume de \ce{CO2} (aux conditions usuelles de température et de pression, $V_m = 24$ L/mol) :
\[ V = n \times V_m = 3 \times 24 = 72 \text{ L}. \]
\end{enumerate}
\textbf{Vérification} : 0,5 mol de glucose (90 g) correspondent bien à 3 mol de \ce{CO2}, soit $3 \times 44 = 132$ g et $3 \times 24 = 72$ L. Le plant de maïs a donc fixé \textbf{132 g de \ce{CO2}}, soit \textbf{72 L}.
\end{corrige}

\begin{corrige}[Exercice 8 — Photosynthèse et respiration : ne pas confondre]
\begin{center}
\begin{tabularx}{\linewidth}{|l|X|X|}
\hline
\rowcolor{popblueL} \textbf{Critère} & \textbf{Photosynthèse} & \textbf{Respiration cellulaire} \\
\hline
Gaz absorbé & \ce{CO2} & \ce{O2} \\
\hline
Gaz dégagé & \ce{O2} & \ce{CO2} \\
\hline
Condition de lumière & Uniquement à la lumière & Jour et nuit (en permanence) \\
\hline
Organite concerné & Chloroplaste & Mitochondrie \\
\hline
Bilan énergétique & Stockage d'énergie (synthèse de glucose) & Libération d'énergie (dégradation du glucose) \\
\hline
\end{tabularx}
\end{center}
\end{corrige}
\begin{attention}[Point de vigilance]
Chez une plante éclairée, les deux processus ont lieu \emph{simultanément} dans les cellules chlorophylliennes ; le dégagement net de \ce{O2} observé signifie seulement que la photosynthèse est plus intense que la respiration.
\end{attention}

\begin{corrige}[Exercice 9 — Type Probatoire : étude expérimentale sur l'élodée]
\begin{enumerate}
\item Le gaz recueilli est le \textbf{dioxygène \ce{O2}}. Test formel : il \textbf{ravive une braise} (une allumette présentant un point incandescent se rallume au contact du gaz).
\item Courbe (à reproduire sur papier millimétré avec les échelles imposées) :
\begin{popfigure}
\begin{center}
\begin{tikzpicture}
\begin{axis}[
width=0.85\linewidth, height=6.5cm,
xlabel={Distance $d$ (cm)}, ylabel={Volume de gaz (mL)},
xmin=0, xmax=90, ymin=0, ymax=13,
xtick={0,10,20,40,60,80}, ytick={0,2,4,6,8,10,12},
grid=major, grid style={popline!50},
axis lines=left]
\addplot[popcurve, mark=*, mark options={fill=poporange}] coordinates {(10,12) (20,9.5) (40,4) (60,1.5) (80,0.5)};
\end{axis}
\end{tikzpicture}
\end{center}
\legende{Volume de gaz dégagé par l'élodée en 10 min en fonction de la distance à la lampe.}
\end{popfigure}
\item La courbe est \textbf{décroissante} : le volume de gaz diminue rapidement quand la distance augmente, puis tend vers une valeur très faible sans s'annuler.
\item L'éclairement diminuant avec la distance, on conclut que \textbf{l'intensité de la photosynthèse augmente avec l'éclairement} : la lumière est un facteur indispensable et limitant de la photosynthèse.
\item À 80 cm, la photosynthèse est très faible mais la plante \textbf{respire} : elle consomme une partie du \ce{O2} qu'elle produit. Le volume mesuré est le dégagement \emph{net} (photosynthèse $-$ respiration), d'où une valeur faible (0,5 mL) mais non nulle.
\item \textbf{Conclusion générale} : en présence de \ce{CO2} et d'eau, une plante verte éclairée réalise la photosynthèse et dégage du \ce{O2} ; l'intensité de ce dégagement croît avec l'éclairement. La lumière est donc une condition indispensable de la photosynthèse.
\end{enumerate}
\textbf{Barème indicatif (sur 10)} : identification du gaz et test (2 pts) ; courbe soignée, axes légendés, échelles respectées (3 pts) ; description (1 pt) ; rôle de la lumière (1,5 pt) ; prise en compte de la respiration (1,5 pt) ; conclusion générale (1 pt).

\textbf{Points de vigilance} : citer le test de la braise et non « le gaz qui s'enflamme » (confusion fréquente avec le dihydrogène) ; préciser que le volume mesuré est un dégagement \emph{net} ; ne pas écrire que la photosynthèse est nulle à 80 cm.
\end{corrige}

\begin{corrige}[Exercice 10 — Type Probatoire : \ce{CO2} atmosphérique et déforestation]
\begin{enumerate}
\item Entre 1970 et 2020, la concentration en \ce{CO2} passe de 325 à 415 ppm (augmentation de 90 ppm, soit environ $+28\,\%$) tandis que la superficie forestière passe de 24,0 à 19,0 millions d'ha (perte de 5 millions d'ha). \textbf{Corrélation} : quand la superficie forestière diminue, la concentration en \ce{CO2} augmente : les deux grandeurs évoluent en sens inverse.
\item Les végétaux chlorophylliens fixent le \ce{CO2} atmosphérique par la photosynthèse ($\ce{6CO2 + 6H2O -> C6H12O6 + 6O2}$) et stockent le carbone dans leur biomasse (troncs, racines) : les forêts constituent des \textbf{puits de carbone}. La déforestation supprime ces fixateurs de \ce{CO2} ; de plus, la combustion ou la décomposition du bois relâche le carbone stocké sous forme de \ce{CO2}. Double effet : moins de fixation, plus de libération, d'où l'accumulation du \ce{CO2} dans l'atmosphère.
\item Deux conséquences :
\begin{itemize}
\item \textbf{Environnement} : le \ce{CO2} est un gaz à effet de serre ; son accumulation provoque le réchauffement climatique (modification des saisons, sécheresses, déséquilibres des écosystèmes).
\item \textbf{Santé humaine} : dégradation des rendements agricoles et insécurité alimentaire, extension des zones de maladies (paludisme), déplacements de populations liés aux catastrophes climatiques.
\end{itemize}
\item Deux solutions applicables au Cameroun :
\begin{itemize}
\item \textbf{Reboisement et agroforesterie} (par exemple dans l'Adamaoua et la zone soudanienne) : les jeunes arbres photosynthétisent activement et fixent le \ce{CO2} dans leur biomasse — mécanisme : photosynthèse, création de puits de carbone.
\item \textbf{Lutte contre les feux de brousse et exploitation forestière durable} (coupe sélective, permis contrôlés dans le bassin du Congo) : on évite la libération massive de \ce{CO2} par combustion et on préserve les puits de carbone existants.
\end{itemize}
\end{enumerate}
\textbf{Barème indicatif (sur 10)} : comparaison chiffrée et corrélation (2,5 pts) ; explication par le cycle du carbone (3 pts) ; deux conséquences (2 pts) ; deux solutions avec mécanisme (2,5 pts).

\textbf{Points de vigilance} : appuyer la comparaison sur des valeurs chiffrées du tableau ; expliciter le double mécanisme (moins de fixation \emph{et} libération du carbone stocké) ; chaque solution doit être reliée à un mécanisme biologique, pas seulement énoncée.
\end{corrige}

\begin{corrige}[Exercice 11 — Type Probatoire : synthèse et argumentation]
\textbf{Partie A.}
\begin{enumerate}
\item \textbf{Interprétations} :
\begin{itemize}
\item \emph{Expérience 1} : l'amidon (mis en évidence par la coloration bleu-noir à l'eau iodée \emph{après décoloration de la feuille à l'alcool chaud}) ne se forme que dans la partie éclairée. La \textbf{lumière} est donc un facteur indispensable de la photosynthèse, dont l'\textbf{amidon} est un produit de stockage.
\item \emph{Expérience 2} : en l'absence de \ce{CO2} (absorbé par la chaux sodée), la plante ne produit pas d'amidon malgré la lumière. Le \textbf{\ce{CO2}} est donc un réactif indispensable.
\item \emph{Expérience 3} : le gaz dégagé ravive une braise : c'est du \textbf{dioxygène \ce{O2}}, produit de la photosynthèse.
\end{itemize}
Bilan : facteurs mis en évidence = lumière et \ce{CO2} ; produits = matière organique (amidon) et \ce{O2}.
\item Équation-bilan équilibrée :
\[ \ce{6CO2 + 6H2O ->[$\text{lumière}$][\text{chlorophylle}] C6H12O6 + 6O2} \]
Origine des réactifs : le \ce{CO2} provient de l'air (absorbé par les stomates des feuilles) ; l'eau est absorbée par les racines et transportée par la sève brute. Devenir des produits : le glucose sert à la respiration, à la croissance ou est stocké sous forme d'amidon ; le \ce{O2} est rejeté dans l'atmosphère.
\end{enumerate}
\textbf{Partie B.}
\begin{enumerate}
\item[3.] \textbf{Paragraphe argumenté (exemple de réponse attendue)} : Les végétaux verts sont effectivement indispensables à la vie sur Terre, et ce pour plusieurs raisons convergentes. D'abord, ils sont les seuls êtres capables, grâce à la photosynthèse, de transformer la matière minérale (\ce{CO2} et eau) en matière organique (glucose) en utilisant l'énergie lumineuse : ils constituent ainsi la \textbf{base de toutes les chaînes alimentaires}, car herbivores, carnivores et décomposeurs dépendent directement ou indirectement de cette production primaire. Ensuite, la photosynthèse dégage le \textbf{dioxygène} indispensable à la respiration de la quasi-totalité des êtres vivants, y compris l'Homme : l'air que nous respirons est un héritage de milliards d'années d'activité photosynthétique. Enfin, en fixant le \ce{CO2} atmosphérique, les végétaux chlorophylliens — et en particulier les grandes forêts comme celle du bassin du Congo — jouent le rôle de \textbf{puits de carbone} qui régulent la composition de l'atmosphère et limitent l'effet de serre. Leur destruction massive, par la déforestation, rompt ces équilibres : moins de \ce{O2} produit, moins de \ce{CO2} fixé, davantage de carbone libéré. Protéger et replanter les végétaux verts, comme le fait la campagne de reboisement de l'Adamaoua, c'est donc préserver les conditions mêmes de la vie sur notre planète.
\end{enumerate}
\textbf{Barème indicatif (sur 10)} : Partie A — interprétation des trois expériences (3 pts), équation-bilan et origine/devenir (2 pts) ; Partie B — paragraphe mobilisant au moins trois savoirs, structuré et argumenté (5 pts).

\textbf{Points de vigilance} : pour l'expérience 1, préciser que la feuille doit être décolorée (ébullition dans l'eau puis alcool) avant le test à l'eau iodée ; ne pas confondre glucose (produit immédiat) et amidon (forme de stockage) ; dans le paragraphe, chaque savoir mobilisé doit être relié à une conséquence concrète pour la vie.
\end{corrige}

\newpage

\coursec{Fiche bilan}

\begin{popfigure}
\begin{center}
\begin{tikzpicture}[
  mindmap,
  grow cyclic,
  every node/.style={concept, font=\small},
  root concept/.append style={concept color=popgreen, font=\small\bfseries, minimum size=2.6cm},
  level 1/.append style={level distance=3.4cm, sibling angle=60},
  level 1 concept/.append style={font=\scriptsize\bfseries, minimum size=1.9cm},
]
\node[root concept]{La photosynthèse}
  child[concept color=popblue]{ node {Conditions : lumière, \ce{CO2}, \ce{H2O}, chlorophylle} }
  child[concept color=poporange]{ node {Siège : le chloroplaste (grana + stroma)} }
  child[concept color=poppurple]{ node {Pigments : chlorophylles a, b, caroténoïdes} }
  child[concept color=popgold]{ node {Bilan : \ce{6CO2 + 6H2O -> C6H12O6 + 6O2}} }
  child[concept color=poppink]{ node {Cycles du carbone et de l'oxygène} }
  child[concept color=popmuted]{ node {Déforestation : déséquilibres} };
\end{tikzpicture}
\end{center}
\legende{Carte mentale de la leçon : la photosynthèse, ses conditions, son siège, son bilan et son rôle dans les équilibres de l'environnement.}
\end{popfigure}

\begin{bilan}
\textbf{Définitions clés}
\begin{itemize}
  \item \cle{Photosynthèse} : ensemble des réactions par lesquelles les végétaux chlorophylliens, en présence de lumière, utilisent le \ce{CO2} de l'air et l'eau du sol pour fabriquer de la matière organique (glucose) en rejetant du dioxygène \ce{O2}.
  \item \cle{Chloroplaste} : organite cytoplasmique des cellules végétales vertes, siège de la photosynthèse ; il contient les pigments dans les \cle{grana} (empilements de thylakoïdes) et un milieu fondamental, le \cle{stroma}.
  \item \cle{Pigments photosynthétiques} : molécules qui absorbent l'énergie lumineuse — chlorophylle $a$ (vert-bleu), chlorophylle $b$ (vert-jaune), caroténoïdes (jaunes-orangés).
  \item \cle{Végétaux chlorophylliens} : végétaux possédant de la chlorophylle (plantes vertes, algues), dits \emph{autotrophes} car ils fabriquent leur propre matière organique.
\end{itemize}

\textbf{Conditions de la photosynthèse (à connaître par cœur)}
\begin{center}
\begin{tabularx}{\linewidth}{|l|X|}\hline
\rowcolor{popgreenL} \textbf{Facteur} & \textbf{Rôle} \\ \hline
Lumière & Source d'énergie indispensable (absorbée par les pigments) \\ \hline
Dioxyde de carbone \ce{CO2} & Source de carbone, prélevé dans l'air par les stomates \\ \hline
Eau \ce{H2O} & Matière première absorbée par les racines \\ \hline
Chlorophylle & Pigment captant l'énergie lumineuse (présent dans les chloroplastes) \\ \hline
\end{tabularx}
\end{center}

\textbf{Bilan global de la photosynthèse (formule à savoir écrire et équilibrer)}
\[ \ce{6CO2 + 6H2O ->[$\text{lumière}$][\text{chlorophylle}] C6H12O6 + 6O2} \]
Interprétation : absorption de \ce{CO2} et d'eau $\rightarrow$ production de glucose (matière organique) et rejet de \ce{O2}. C'est l'inverse du bilan de la respiration cellulaire.

\textbf{Méthodes express}
\begin{itemize}
  \item \emph{Mettre en évidence le rejet de \ce{O2}} : expérience de l'Élodée sous entonnoir éclairé $\rightarrow$ dégagement gazeux qui ravive une braise (c'est du dioxygène).
  \item \emph{Mettre en évidence la fixation du \ce{CO2}} : deux plantes à la lumière, l'une avec potasse (absorbe le \ce{CO2}) $\rightarrow$ seule la plante avec \ce{CO2} produit de l'amidon (test à l'eau iodée : bleu-noir).
  \item \emph{Rôle de la lumière et de la chlorophylle} : feuille partiellement masquée ou panachée $\rightarrow$ l'amidon n'apparaît que dans les parties vertes éclairées.
  \item \emph{Argumenter sur les cycles} : photosynthèse = pompage du \ce{CO2} et source de \ce{O2} ; respiration et combustion = rejets de \ce{CO2}. Les végétaux verts régulent la composition de l'atmosphère.
  \item \emph{Déforestation} : moins de chloroplastes $\Rightarrow$ moins de \ce{CO2} fixé et moins de \ce{O2} produit $\Rightarrow$ accumulation de \ce{CO2} atmosphérique (effet de serre, réchauffement climatique), érosion des sols, perte de biodiversité.
\end{itemize}
\end{bilan}

\begin{aretenir}[Checklist avant le Probatoire]
\begin{itemize}
  \item[$\square$] Je sais définir la photosynthèse en une phrase complète.
  \item[$\square$] Je cite les quatre conditions nécessaires : lumière, \ce{CO2}, eau, chlorophylle.
  \item[$\square$] Je sais écrire et équilibrer le bilan : \ce{6CO2 + 6H2O -> C6H12O6 + 6O2}.
  \item[$\square$] Je sais décrire et légender un chloroplaste (grana, thylakoïdes, stroma, double membrane).
  \item[$\square$] Je cite les pigments photosynthétiques et leur rôle dans l'absorption de la lumière.
  \item[$\square$] Je sais décrire l'expérience de l'Élodée et interpréter le test à la braise.
  \item[$\square$] Je sais interpréter le test de l'amidon à l'eau iodée (coloration bleu-noir).
  \item[$\square$] Je sais exploiter un dispositif expérimental avec témoin (avec/sans \ce{CO2}, lumière/obscurité).
  \item[$\square$] J'explique le rôle des végétaux verts dans les cycles du carbone et de l'oxygène.
  \item[$\square$] Je distingue photosynthèse (jour, rejet de \ce{O2}) et respiration (jour et nuit, rejet de \ce{CO2}).
  \item[$\square$] Je cite au moins trois conséquences de la déforestation sur l'équilibre de l'environnement.
  \item[$\square$] Je sais argumenter avec des exemples camerounais (forêt du bassin du Congo, reboisement).
\end{itemize}
\end{aretenir}

\findecours

\end{document}
